\documentclass[11pt]{article}

% ---------- Encoding / fonts ----------
\usepackage[T1]{fontenc}
\usepackage[utf8]{inputenc}
\usepackage{lmodern}        % Latin Modern: reliable math/text glyph coverage

% ---------- Layout / typography ----------
\usepackage[margin=1in]{geometry}
\usepackage{microtype}      % better PDF text shaping; helps readability
\usepackage{xcolor}

% ---------- Math packages ----------
\usepackage{amsmath,amssymb,amsthm,mathtools}
\usepackage{bm}             % bold math
\usepackage{stmaryrd}       % provides \llbracket \rrbracket
\usepackage{mathrsfs}       % optional: \mathscr

% ---------- Lists / diagrams ----------
\usepackage{enumitem}       % consistent lists
\usepackage{tikz-cd}        % commutative diagrams

% ---------- Hyperlinks / bookmarks ----------
\usepackage{hyperref}
\usepackage{bookmark}
\usepackage{natbib}

\hypersetup{
  colorlinks=true,
  linkcolor=blue,
  citecolor=blue,
  urlcolor=blue
}

% ============================================================
% Theorem environments
% ============================================================

\newtheorem{theorem}{Theorem}[section]
\newtheorem{lemma}[theorem]{Lemma}
\newtheorem{proposition}[theorem]{Proposition}
\newtheorem{corollary}[theorem]{Corollary}

\theoremstyle{definition}
\newtheorem{definition}[theorem]{Definition}
\newtheorem{remark}[theorem]{Remark}
\newtheorem{example}[theorem]{Example}
\newtheorem{assumption}[theorem]{Assumption}

% ============================================================
% FRFP Macros (use everywhere; avoids symbol drift)
% ============================================================

\newcommand{\FRFP}{FRFP}
\newcommand{\Cfull}{C_{\mathrm{full}}}
\newcommand{\Ecat}{\mathcal{E}} % explicit category
\newcommand{\Tcat}{\mathcal{T}} % tacit category
\newcommand{\Ncat}{\mathcal{N}} % navigation groupoid
\newcommand{\Scat}{\mathcal{S}} % shapes
\newcommand{\TDG}{\mathsf{TDG}} % Task Decomposition Grammar
\newcommand{\RB}{\mathsf{RB}}   % Representational Backflow
\newcommand{\RI}{\mathsf{RI}}
\newcommand{\EC}{\mathsf{EC}}
\newcommand{\ED}{\mathsf{ED}}
\newcommand{\TE}{\mathsf{TE}}
\newcommand{\HFD}{\mathsf{HFD}}
\newcommand{\Req}{\mathsf{Req}}
\newcommand{\Hom}{\mathrm{Hom}}
\newcommand{\Obj}{\mathrm{Ob}}
\newcommand{\BRANCH}{\mathsf{BRANCH}}
\newcommand{\PARALLEL}{\mathsf{PARALLEL}}
\newcommand{\LOOP}{\mathsf{LOOP}}

% aliases for explicit / tacit space used in the appendix
\newcommand{\Espace}{\mathcal{E}} % explicit space
\newcommand{\Tspace}{\mathcal{T}} % tacit space

% encoding / decoding
\newcommand{\enc}{\mathsf{enc}}
\newcommand{\dec}{\mathsf{dec}}

% semantic RB and free-category notation
\newcommand{\RBsharp}{\RB^{\sharp}}
\newcommand{\FreeCat}{\mathsf{FreeCat}}
\newcommand{\Ufun}{U_E}

% shape constructors (non-conflicting names)
\newcommand{\chain}{\mathsf{chain}}
\newcommand{\branch}{\mathsf{branch}}
\newcommand{\loopshape}{\mathsf{loop}}

% frameworks category name: \Frm_1
\newcommand{\Frm}{\mathsf{Frm}} % so \Frm_1 works

% example navigation generators (purely illustrative)
\newcommand{\TRYMODEL}{\mathsf{TRYMODEL}}
\newcommand{\ADDASSUMPTION}{\mathsf{ADDASSUMPTION}}
\newcommand{\DROPASSUMPTION}{\mathsf{DROPASSUMPTION}}
\newcommand{\MERGE}{\mathsf{MERGE}}
\newcommand{\RESTART}{\mathsf{RESTART}}

\newcommand{\id}{\mathrm{id}}
\newcommand{\shape}{\mathsf{shape}}

\newcommand{\Inst}{\mathsf{Inst}}
\newcommand{\HP}{\mathcal{HP}}
\newcommand{\Sem}{\mathsf{Sem}}
\newcommand{\Esem}{\mathcal{E}_{\mathrm{sem}}}

\newcommand{\Empty}{\varnothing}
\newcommand{\Eobj}{E}
\newcommand{\Tobj}{T}

% Semantic brackets
\newcommand{\sem}[1]{\llbracket #1 \rrbracket}

% Semidirect / Grothendieck notation
\newcommand{\NltimesE}{\Ncat \ltimes \Ecat}
\newcommand{\NltimesEsem}{\Ncat \ltimes \Esem}

% Logic-ish convenience
\newcommand{\defeq}{\mathrel{\vcentcolon=}}
\newcommand{\eqdef}{\mathrel{=\vcentcolon}}


\title{Foundational Reasoning and Feedback Protocol (FRFP):\\
A Formal Architecture for Long-Horizon Human--AI Scientific Collaboration}

\author{Vijay Nadendla\\[0.3ex]
Independent Researcher\\
vnadendla@gmail.com}
\date{\today}

\begin{document}
\maketitle


\begin{abstract}
We propose the Foundational Reasoning and Feedback Protocol (FRFP), a formal architecture for long-horizon scientific collaboration between human agents and AI systems. FRFP enforces a sharp separation between explicit representations, which may be manipulated by machines, and tacit epistemic states, which are exclusive to human agents and bear correctness authority. On top of this separation, FRFP builds a free algebra of explicit task decompositions and navigation steps, realized as a Grothendieck-style semidirect product \(N \ltimes E\) of explicit pipelines with a navigation groupoid. A semantic functor \(J{-}K : N \ltimes E \to T\) into a preorder-enriched tacit category, together with a correctness quotient \(\gamma : \mathrm{Ob}(T) \to J\), yields observable judgments that are invariant under admissible refactorings and normalization in \(N \ltimes E\).

Using this architecture, we develop a dynamic theory of long-run interaction. Tacit context degrades under use, requirement maps encode what is needed to safely endorse explicit artifacts in a given context, and hallucination is defined structurally as a mismatch between tacit state and requirements at the moment of acceptance. Under qualitative but pessimistic assumptions on degradation, feedback, and recurring nontrivial tasks, we prove inevitability theorems: along sufficiently long admissible runs, hallucination occurs almost surely, and protocol-level ``safe horizons'' are necessarily finite. We then lift the framework to epistemic populations, institutions, and governance operators, showing that explicit-only consensus and monitoring cannot in general eliminate hallucination or recover tacit correctness. Finally, we extract protocol and interface rules---together with FRFP-compatible system instructions---for building explicit-only AI tools that support powerful exploration while preserving human-exclusive grounding and auditability at the level of whole scientific and regulatory workflows.

Appendix~\ref{app:business-frfp}  provides a practitioner-facing ``business version'' of FRFP: a non-technical overview of the architecture, examples in scientific and regulated workflows, an implementation roadmap, common pitfalls, and a short design checklist for teams deploying AI systems.

\end{abstract}


\pagebreak
\tableofcontents
\pagebreak
%=========================================================
\section{Introduction}
\label{sec:intro}
%=========================================================
Large language models and related AI systems are increasingly used as partners in scientific reasoning: they propose conjectures, sketch proofs, draft experiments, and synthesise literature. This creates both opportunities and risks. AI systems can operate over large explicit structures that are difficult for humans to inspect directly, but they are also known to hallucinate---producing plausible but ungrounded content, sometimes with high confidence. In most current workflows, humans and AI systems interact in an untyped, informal way. Prompts, replies, and corrections share a single channel; correctness judgments are not explicitly separated from heuristic suggestion; and long interaction traces often suffer from semantic drift. This makes it difficult to reason formally about correctness, responsibility, and long-run behaviour.

We take a different approach. We propose a protocol-level design for human--AI collaboration, the \emph{Foundational Reasoning and Feedback Protocol (FRFP)}, with three aims:
\begin{itemize}
  \item to provide a formal semantics for mixed human--AI reasoning trajectories;
  \item to make explicit the boundary between what AI systems may and may not do in scientific workflows;
  \item to enable theorem-level analysis of correctness and long-run behaviour (e.g., inevitability of hallucination under certain conditions).
\end{itemize}

The central design choice in FRFP is a sharp separation between:
\begin{itemize}
  \item \emph{Explicit mode}: machine-manipulable representations and pipelines;
  \item \emph{Tacit mode}: human-exclusive evaluative and correctness-bearing states.
\end{itemize}

AI systems are confined to explicit mode; only humans inhabit tacit mode. A single boundary morphism \(RB : E \to T\) transports explicit artifacts into tacit evaluation, and there are no morphisms from tacit to explicit space. On top of this structure, FRFP builds a free algebra of explicit pipelines (TDG), a navigation groupoid \(N\) for reversible exploration, and a semidirect product \(N \ltimes E\) integrating both. A semantic functor \(\llbracket - \rrbracket : N \ltimes E \to T\) maps trajectories to tacit states, and observable correctness is given by a quotient \(\gamma : \mathrm{Ob}(T) \to J\) into a finite set of judgments. Correctness and meaning are invariant under all admissible transformations in \(N \ltimes E\). Appendix~B then uses this structure to study pipeline dynamics: how tacit context degrades, how confidence feedback and requirement escalation interact with degradation, and when hallucinations become structurally inevitable.

\paragraph{Scope of novelty.}
We use standard categorical constructions---free categories, groupoids, fibrations, Grothendieck constructions---but assemble them to capture explicit/tacit separation and feedback in human--AI workflows. The novelty lies in (i) isolating a small set of epistemic primitives (explicit/tacit separation, boundary, feedback), (ii) combining them into a compositional architecture in which system-level properties (auditability, correctness preservation, hallucination inevitability) become theorems, and (iii) doing so at a level of abstraction that allows later work to reuse existing categorical theory rather than inventing bespoke formalisms for each application domain. FRFP aims to place currently informal design intuitions---about tool use, oversight, and responsibility---on a mathematical footing that can support future advances.

\paragraph{A business and governance-oriented companion.}
While most of this paper is written at a formal and architectural level, many readers will be interested primarily in what FRFP implies for concrete organizational practice.
Appendix~\ref{app:business-frfp} therefore provides a practitioner-facing ``business version'' of FRFP: a non-technical overview of the architecture, examples in scientific and regulated workflows, an implementation roadmap, common pitfalls, and a short design checklist for teams deploying AI systems.

\subsection{Contributions}

Formally, this paper makes the following contributions.

\begin{enumerate}
  \item \textbf{Explicit–tacit epistemic architecture.}
  We define an ambient category $\mathcal{C}_{\mathrm{full}}$ with explicit and tacit subcategories $E$ and $T$, six epistemic primitives $\mathrm{RI}, \mathrm{EC}, \mathrm{ED}, \mathrm{RB}, \mathrm{TE}, \mathrm{HFD}$, and a correctness architecture
  \[
    \mathrm{RI} \;\rightarrow\; (\mathrm{EC}/\mathrm{ED})^{\ast} \;\rightarrow\; \mathrm{RB} \;\rightarrow\; \mathrm{TE} \;\rightarrow\; \mathrm{HFD}.
  \]
  We prove that, under a set of Phase–1 axioms, this architecture is initial among suitable epistemic ontologies, providing a canonical separation between explicit representations and tacit correctness-bearing states.

  \item \textbf{Algebra of explicit reasoning and navigation.}
  We introduce a Task Decomposition Grammar (TDG) generating explicit pipelines and a navigation groupoid $N$ for reversible exploration. These combine into a Grothendieck-style semidirect-product category $N \ltimes E$, equipped with a confluent, terminating rewrite system admitting canonical representatives and decidable equality of explicit pipelines, subject to Phase–0 operational preconditions.

  \item \textbf{Semantic correctness and invariance.}
  We define a semantic functor $J{-}K : N \ltimes E \to T$ that composes explicit pipelines with the boundary operator $\mathrm{RB}$ and tacit operators $\mathrm{TE}, \mathrm{HFD}$, together with a correctness quotient $\gamma : \mathrm{Ob}(T) \to J$ into a finite space of observable judgments. We show that observable correctness $\gamma \circ J{-}K$ is invariant under all admissible explicit and navigation steps, yielding canonical correctness forms at the level of whole pipelines.

  \item \textbf{Dynamic theory of tacit context and hallucination.}
  We enrich the tacit space with a degradation operator $\delta : T \to T$, a requirement map $\mathrm{Req} : E \times \mathbb{R}_{\ge 0} \to T$, plausibility and confidence operators, and a feedback-coupled degradation operator $\delta^{\star}$. Under qualitative assumptions such as monotone degradation and recurring nontrivial requirements, we prove inevitability theorems: along sufficiently long admissible runs, hallucinations—mismatches between tacit state and requirements at the moment of acceptance—occur inevitably, and protocol-level safe horizons are necessarily finite.

  \item \textbf{Stochastic and finite-horizon pipelines.}
  We extend the dynamic analysis to stochastic implementations and ensembles. Pathwise inevitability implies almost-sure inevitability under any stochastic realization of a given pipeline. We study finite-horizon hallucination probabilities $p_N$, introduce hazard profiles $(h_n)_{n \ge 0}$ and a safe-horizon functional $H_P(\varepsilon)$, and define a dynamic refinement order on implementations based on their finite-horizon risk curves, separating static semantic equivalence from dynamic safety properties.

  \item \textbf{Collective epistemic dynamics.}
  We lift the single-agent FRFP architecture to populations of epistemic agents with shared explicit category $C^I$, agent-local tacit spaces, and communication graphs. Collective executions are modeled as Markov processes on population states, yielding collective hallucination times and associated stopping-time semantics. We represent consensus mechanisms as explicit-only functors and characterize when collective designs are equivalent, dynamically refined, or riskier, including an inevitability result for hallucination at the population level and a non-identifiability result for correctness from explicit logs alone.

  \item \textbf{Limits of explicit-only mechanisms and institutional non-automation.}
  Building on explicit–tacit separation and human-exclusive grounding, we develop a family of impossibility results showing that a broad class of desirable guarantees are not explicit-implementable: there is no perfect explicit-only hallucination detector, no fully automated correctness oracle, no grounding gain from explicit-only consensus, and no fully automated governance operator that preserves correctness across all admissible pipelines and populations. We derive an institutional non-automation theorem and a structural necessity result for human oversight at the governance layer.

  \item \textbf{Protocols, case studies, and FRFP-compatible system instructions.}
  Finally, we translate the formal architecture into protocol-level design principles for long-horizon human–AI collaboration. We state foundational rules for FRFP-compatible practice (explicit confinement, human-exclusive grounding, source–reason traceability, integrity under navigation), instantiate the framework in detailed case studies (ICU sepsis prediction literature review and financial scenario analysis / stress testing), and provide a concrete set of system instructions for explicit-only AI tools. These constructions illustrate how the abstract FRFP structure constrains interfaces, workflows, and institutional arrangements in scientific and regulatory settings.
\end{enumerate}


The detailed categorical proofs appear in Appendices; the main body of the paper focuses on conceptual structure and the resulting design constraints for human--AI scientific collaboration.

\paragraph{Core kernel vs.\ outer layers.}
Conceptually, the non-negotiable kernel of FRFP is the Phase~0–2 correctness architecture and its categorical realization: an explicit/tacit modality with a single boundary $RB : E \to T$, the six primitives $\{RI, EC, ED, RB, TE, HFD\}$, the Task Decomposition Grammar and navigation groupoid combined into the semidirect product $N \ltimes E$, and the tacit correctness functor $\gamma \circ \llbracket - \rrbracket$ that makes correctness a purely tacit property of explicit pipelines.
This kernel fixes what counts as an admissible pipeline and how correctness can (and cannot) be represented.
Sections~6--10 and Appendices~B--F should then be read as exploring one particular choice of dynamic, population, institutional, and governance layers built on top of that kernel: specific models of tacit degradation and requirements, stochastic runs and safe horizons, collective Markov dynamics, institutional aggregation, and non-automated governance.
In principle, these outer layers could be varied while preserving the FRFP kernel unchanged, so the central contribution of the paper lies in the kernel itself rather than in any particular dynamic or institutional parametrization.

\subsection{Related Work}

Our work sits at the intersection of formal semantics, human--AI interaction, AI safety governance, and social epistemology of science. In this section we contrast FRFP with three broad strands of prior work: traditional HCI-style human--AI interaction frameworks, concrete AI safety governance and evaluation proposals, and formal models of scientific communities. We then briefly connect to mechanized reasoning, categorical semantics, and the hallucination/calibration literature.

\paragraph{Interactive proof, type theory, and mechanized reasoning.}
Proof assistants and type-theoretic systems (e.g., Coq, Isabelle, Lean) distinguish between \emph{proof objects} and \emph{proof scripts}, often with rich tactic languages and automation, and they provide strong soundness guarantees for derivations inside a formal calculus. Users are typically treated as trusted sources of goals and proof terms; the core question is whether a given term inhabits a type or whether a proof object checks against a specification. FRFP is orthogonal: it does not propose a new logic or proof theory, and it does not model correctness as a syntactic property of explicit objects. Instead, correctness is a tacit attribute of human epistemic states, accessed only via the boundary morphism and tacit evaluation primitives. In particular, FRFP is designed for scientific and regulatory workflows where correctness cannot be reduced to proof checking or type inference, and where long-horizon interaction, degradation of tacit grounding, and institutional structure dominate the risk picture.

\paragraph{Category theory and categorical semantics.}
The categorical tools we employ---free categories, groupoids, (op)fibrations, and Grothendieck constructions---are standard in categorical semantics of programming languages and logics \citep{johnstone2002sketches,leinster2014basic,maclane1998categories}. Traditional applications focus on modeling syntax/semantics maps, type systems, effects, and compositional program behavior. Our contribution is to repurpose these tools for epistemic roles in mixed human--AI reasoning: explicit and tacit subcategories; a one-way boundary morphism; a navigation groupoid capturing refactoring and exploration; and a semidirect product $N \ltimes E$ that treats entire workflows as single compositional objects. This allows us to phrase auditability, invariance of correctness under refactoring, and impossibility of explicit-only grounding as theorems about a unified architecture, rather than as informal desiderata or system-specific properties.

\paragraph{Rewriting, confluence, and canonical forms.}
The explicit reduction systems and confluence arguments we use draw on the abstract reduction systems and termination/confluence methods of term rewriting \citep{baader1999term,terese2003term}. What is novel is not the rewriting technology, but its epistemic role: we use normalization in $N \ltimes E$ to show that all admissible explicit refactorings and navigation steps collapse to canonical forms that induce the same tacit correctness judgments. In other words, FRFP uses rewriting theory to guarantee that correctness is invariant under the kinds of exploration and restructuring that are unavoidable in long-horizon human--AI pipelines.

\paragraph{HCI-style human--AI interaction frameworks.}
A large HCI and human--AI interaction literature examines how humans and AI systems should share tasks, communicate, and co-adapt. Recent surveys and frameworks organize this space along dimensions such as locus of control, initiative, transparency, and feedback channels \citep{jiang2024hci_research_agenda,gomez2025interaction_taxonomy,fragiadakis2024haic_evaluation}. Reports from organizations such as the Partnership on AI propose design frameworks for human--AI collaboration that emphasize roles, affordances, and case studies in domains like healthcare and hiring \citep{pai2021haic_framework}. Empirical work investigates trust, mental models, and usability in AI-assisted decision making and creative work \citep{sidra2025generative_haic_health,haicSurvey2024}.

These HCI-style frameworks typically share three characteristics. First, they treat the human--AI boundary at the level of \emph{tasks, interfaces, and user experience}, not as a formal semantic construct. Second, they model correctness and responsibility implicitly, via notions like ``appropriate reliance'', ``overtrust'', or ``calibrated assistance'', rather than as objects in a correctness semantics. Third, they often assume that the AI system may participate in evaluation or recommendation in ways that are not sharply distinguished from human judgment, provided the interface is designed well and users understand the system's limitations.

FRFP is compatible with much of this work but takes a different starting point. Instead of asking how to make particular interactions smoother or more interpretable, FRFP asks: what are the \emph{allowable roles} for AI systems if correctness is, by construction, reserved for humans? In our architecture:

\begin{itemize}
  \item All AI-mediated operations live in explicit space $E$; systems cannot cross the boundary $RB : E \to T$ or take tacit states as inputs.
  \item Correctness is a tacit property of human epistemic states, modeled by a quotient of the tacit category; AI systems cannot be assigned correctness authority without exiting the FRFP kernel.
  \item Interaction patterns (e.g., suggestion, critique, refactoring) are treated as morphisms in $N \ltimes E$, not as ad hoc UI flows. This makes it possible to prove invariance theorems and impossibility results about whole classes of interactions.
\end{itemize}

In this sense, FRFP complements HCI-style frameworks: it provides an architectural and semantic envelope within which HCI choices must fit, and within which familiar notions like ``human-in-the-loop'' or ``appropriate reliance'' can be given a more precise, protocol-level meaning.

\paragraph{AI safety governance, risk committees, and evaluation protocols.}
A parallel literature in AI safety and governance focuses on organizational and regulatory mechanisms for managing risk from powerful AI systems. Frontier model developers and policy groups have proposed safety frameworks, risk management processes, and governance structures that include elements such as threat modeling, red-team evaluations, incident response, and risk committees \citep{frontierModelForum2024components,robinson2024transforming_risk_governance,governanceai2023frontier_regulation,rand2025governance_frontier_ai,concordia2025frontier_risk}. Government-aligned initiatives emphasize standardized evaluation suites, secure development practices, and reporting obligations \citep{aisi2025inspect_ai,cset2025ai_preparedness}. These proposals typically adopt a ``three lines of defense'' or similar governance model, where technical evaluations feed into internal risk committees and, in some cases, external oversight.

Most of this work operates at the level of institutional design, process, and regulatory instruments. It assumes that we can define evaluation protocols and governance mechanisms that, if properly implemented, yield acceptable levels of safety for certain classes of harmful behavior. But correctness and hallucination are usually treated as properties of \emph{systems} under evaluation---for example, as failure rates in benchmark suites or red-team exercises---rather than as emergent properties of human--AI \emph{protocols} and institutions.

FRFP differs in three main ways:

\begin{enumerate}
  \item \textbf{Protocol-level semantics.}
  Instead of treating evaluations or risk committees as external to the semantics of model use, FRFP builds them into the architecture. Institutions and governance operators are explicit constructions on top of $N \ltimes E$ and collections of tacit agents. This lets us express questions like ``Can an explicit-only consensus mechanism recover correctness?'' or ``Can a fully automated governance oracle exist?'' as formal questions about functors and morphisms, and then prove impossibility results.

  \item \textbf{Explicit-only impossibility results.}
  While many governance proposals emphasize the importance of human oversight, they often leave open whether, in principle, sufficiently good monitoring or evaluation might allow partial or full automation of certain governance functions. Under FRFP’s explicit/tacit separation and correctness axioms, we can show that a broad class of explicit-only detectors, consensus mechanisms, and governance operators cannot, even in principle, recover tacit correctness or eliminate hallucination across all admissible pipelines and populations. These no-go theorems sharpen and formalize the intuition that some oversight roles must remain human-exclusive.

  \item \textbf{Dynamic safe horizons.}
  Existing risk management frameworks often think in terms of static model capabilities and evaluation scores. FRFP introduces dynamic notions such as degradation, requirement schedules, hazard profiles, and safe-horizon functionals for entire \emph{workflows}. This shifts attention from ``Is the model safe on benchmark $X$?'' to ``For how long can a given protocol be run, under which institutional controls, before hallucination risk exceeds a tolerable threshold?'' and ``How do governance choices change these horizons?''.
\end{enumerate}

In short, FRFP does not compete with concrete governance frameworks; rather, it offers a semantic layer that can be used to articulate why certain governance roles must be human and how evaluation, monitoring, and escalation should be understood relative to tacit correctness.

\paragraph{Epistemic models of scientific communities.}
A substantial body of work in formal social epistemology and philosophy of science models scientific communities using tools such as multi-armed bandits, network epistemology, and agent-based simulations \citep{kitcher1990division_epistemic,weisbergMuldoon2009epistemic_landscapes,zollman2007network_epistemology,zollman2010epistemic_networks,oconnor2021modelling_scientific_communities,reijulaKuorikoski2017modeling_epistemic_communities,seseljaStraesserBorg2020formal_models_scientific_inquiry,politi2021formal_models_scientific_community}. These models investigate how factors like credit incentives, network topology, methodological diversity, and social learning affect convergence to truth, replication success, or the division of cognitive labor.

This literature shares with FRFP an interest in:

\begin{itemize}
  \item the relationship between individual rationality and collective epistemic performance;
  \item the effects of institutional and structural constraints on inquiry; and
  \item the role of communication networks in shaping epistemic outcomes.
\end{itemize}

However, there are important differences:

\begin{enumerate}
  \item \textbf{Level of abstraction and target.}
  Most social epistemology models treat scientists as stylized agents who choose among research options, share evidence or beliefs, and update via simple Bayesian or rule-based heuristics. The goal is to understand broad patterns like the benefits of transient diversity or the dangers of premature consensus. FRFP, by contrast, is explicitly targeted at \emph{protocols for human--AI scientific workflows}. It does not attempt to realistically model individual scientists' psychology or incentives; instead, it provides a high-level semantics for how explicit artifacts and tacit correctness interact when AI systems are introduced.

  \item \textbf{Explicit/tacit modality.}
  Existing models rarely distinguish explicitly between \emph{representations} and \emph{epistemic states}, let alone reserve correctness to one modality. Belief states are typically explicit probability vectors or discrete positions on an epistemic landscape. FRFP, by contrast, builds an explicit/tacit modality into the foundation: explicit objects are manipulable by AI systems and institutions, while tacit states remain human and correctness-bearing. This distinction is what allows us to prove that certain institutional arrangements cannot, even in principle, reconstruct correctness from explicit data alone.

  \item \textbf{Institutional operators as functors on $N \ltimes E$.}
  Social epistemology models often encode institutions implicitly via payoff functions, communication graphs, or update rules. FRFP treats institutions as explicit operators that aggregate explicit artifacts and tacit judgments generated by $N \ltimes E$-pipelines across a population. Governance is then modeled as a meta-level operator acting on these institutions and their pipelines. This categorical stance allows us to formulate and prove structural results (e.g., non-automation theorems) that are largely independent of the details of any particular agent-based model.
\end{enumerate}

We see FRFP as complementary to this literature: it suggests a way of embedding AI systems into formal models of scientific communities while preserving a principled distinction between explicit computation and tacit correctness, and it offers a language in which to express normative constraints on institutional design in AI-augmented science.

\paragraph{Hallucination, calibration, and uncertainty in language models.}
Empirical work has documented hallucinations in large language models and has explored mitigation techniques ranging from retrieval augmentation to uncertainty estimation. Surveys such as Ji et al.\ systematize hallucination phenomena and benchmarks \citep{ji2023hallucination_survey}. Semantic uncertainty methods, such as semantic entropy, use distributional properties over outputs to flag likely hallucinations \citep{zhou2023semantic_entropy}. Other work studies calibration and self-knowledge in language models, asking when models can assess their own correctness \citep{kadavath2022self_knowledge}, or analyzes why standard training pipelines naturally yield hallucinations and miscalibration \citep{kalai2024hallucination_explained}.

These approaches typically treat hallucination as a misalignment between model outputs and some ground-truth distribution or benchmark. They ask whether we can detect or reduce hallucination by changing architectures, training objectives, or decoding strategies. FRFP reframes hallucination as a \emph{protocol-level} property: a structural mismatch between tacit context and task requirements at the moment when an explicit artifact is accepted. This allows us to prove inevitability results that are independent of the internal details of any particular model or detector, and to connect hallucination risk directly to the design of workflows, institutions, and governance mechanisms rather than solely to model internals.

\paragraph{Human--AI collaboration, scientific workflows, and protocol design.}
Finally, there is a growing body of work on using AI systems to support scientific discovery, literature review, hypothesis generation, and experimental design. Many of these proposals focus on specific tasks (e.g., protein design, theorem conjecturing) or on augmenting particular phases of the scientific process. FRFP offers a complementary, protocol-level perspective: it treats an entire long-horizon scientific workflow---including problem formulation, literature review, modeling, evaluation, and institutional decision making---as a composition in $N \ltimes E$ with a fixed boundary to tacit space. Within this frame, familiar ideas such as ``AI copilots for science'', ``risk committees'', or ``red-team evaluations'' become concrete patterns of morphisms and functors. What is distinctive about FRFP, relative to the literatures above, is the combination of:

\begin{itemize}
  \item an explicit/tacit semantic architecture with human-exclusive correctness;
  \item a dynamic theory of tacit degradation, requirements, and protocol-level hallucination with formal inevitability theorems; and
  \item an institutional and governance layer that yields structural non-automation and explicit-only impossibility results.
\end{itemize}

Together, these ingredients provide a unified language in which to reason about the design space of long-horizon human--AI scientific collaboration, in a way that is compatible with---but more constraining than---existing HCI, governance, and social epistemology frameworks.

%=========================================================
\section{Phase 0: Operational Preconditions}
\label{sec:phase0}
%=========================================================

Phase~0 describes constraints on human--AI interaction \emph{before}
any algebraic structure is introduced. Its purpose is to rule out
pathological patterns in which AI systems implicitly appropriate tacit
roles or in which explicit and tacit modes are conflated at the UI
layer.

\subsection{Interaction Traces}

We model an interaction session as a sequence of messages tagged by
their origin.

\begin{definition}[Interaction trace]
A \emph{human--AI interaction trace} is a finite sequence
\[
S = (m_1, m_2, \dots, m_n)
\]
where each message \(m_i\) is tagged as either human-generated (H) or
AI-generated (A), and carries explicit content (e.g., text, code,
mathematical expressions).
\end{definition}

We write \(S_k = (m_1,\dots,m_k)\) for the \(k\)-prefix of \(S\). Phase~0
consists of a family of constraints \(\Phi_0\) on such prefixes.

\subsection{Phase--0 Constraints}

The constraints are operational rather than mathematical. They enforce
mode discipline and forbid the AI from taking tacit roles.

\begin{description}[leftmargin=2em]
  \item[Intent locking.] Early in the trace, the human must state a
        coarse-grained intent (e.g., ``analyze this conjecture'',
        ``design an experiment''). The AI must treat this intent as fixed
        unless explicitly revised. This prevents uncontrolled drift in
        task framing.

  \item[Ambiguity resolution.] When an AI message admits more than one
        plausible interpretation relevant to the task, the next AI move
        must be a clarification query, not a content-producing response.
        This prohibits silent resolution of deep ambiguity.

  \item[Mode discipline.] Messages must be typed as either explicit
        (facts, derivations, questions) or meta-level (requests for
        tacit judgment). The AI may only produce explicit content.

  \item[No tacit emulation.] The AI is prohibited from simulating tacit
        operations such as evaluating overall correctness or ``deciding''
        whether a result should be accepted. It may suggest checks, but
        not act as a source of correctness.

  \item[Correctness stability.] Once the human has issued a tacit
        correctness judgment, the AI may not reinterpret it or override
        it via further explicit manipulation. Tacit states are fixed
        points with respect to AI operations.
\end{description}

\begin{definition}[Phase--0 admissibility]
An interaction trace \(S\) is \emph{Phase--0 admissible} if every prefix
\(S_k\) satisfies all Phase--0 constraints. We write \(S \vDash \Phi_0\)
for this property.
\end{definition}

Phase~0 is intentionally conservative: it rules out many interaction
patterns that might be useful in practice but would make formal analysis
intractable. Later work may relax these constraints in calibrated ways.

%=========================================================
\section{Phase 1: Epistemic Ontology}
\label{sec:phase1}
%=========================================================

Phase~1 introduces the core epistemic ontology: explicit and tacit
categories, primitives, and correctness architecture. The goal is to
capture, at a structural level, who can do what in a scientific
workflow.

\subsection{Ambient Category and Explicit--Tacit Separation}

We work inside an ambient category \(\Cfull\) with distinguished objects
\(\varnothing\), \(E\), and \(T\), as detailed in Appendix~A. The
Explicit--Tacit Separation axiom (ETS) asserts:

\begin{itemize}[leftmargin=2em]
  \item explicit morphisms end at \(E\);
  \item tacit morphisms end at \(T\);
  \item there are no morphisms \(T \to E\);
  \item there is a unique boundary morphism \(\RB:E\to T\);
  \item \(\RB\) is not invertible.
\end{itemize}

Intuitively, explicit and tacit modes are \emph{disjoint}. There is a
one-way flow from explicit artifacts into tacit evaluation, and no way
back.

\subsection{Explicit and Tacit Primitives}

Within \(\Cfull\), we identify six primitive morphisms that play
irreducible epistemic roles:

\begin{itemize}[leftmargin=2em]
  \item \(\RI : \varnothing \to E\) (Representation Initiation),
  \item \(\EC : E \to E\) (Explicit Computation),
  \item \(\ED : E \to E\) (Explicit Diagnostics),
  \item \(\RB : E \to T\) (Representational Backflow),
  \item \(\TE : T \to T\) (Tacit Evaluation),
  \item \(\HFD : T \to T\) (Human Final Decision).
\end{itemize}

The first three live in explicit space; the last three in tacit space.
They are combined into a correctness architecture:

\[
\RI \;\to\; (\EC/\ED)^\ast \;\to\; \RB \;\to\; \TE \;\to\; \HFD.
\]

Here \((\EC/\ED)^\ast\) denotes an arbitrary (possibly empty) finite
sequence of explicit computations and diagnostics. The pipeline ends
with tacit evaluation and final decision; no correctness-bearing steps
occur in explicit space.

\subsection{Correctness in Tacit Space}

Tacit states form a preorder-enriched one-object category \(\Tcat\) with
primitives \(\TE,\HFD\) and a correctness quotient
\(\gamma:\Obj(\Tcat)\to J\) into a finite set of judgments. The key fact
for the main paper is:

\begin{proposition}[Tacit correctness quotient]
Tacit primitives \(\TE,\HFD:T\to T\) are idempotent and commute, and
the correctness quotient \(\gamma\) satisfies
\(\gamma\circ\TE=\gamma=\gamma\circ\HFD\). Thus tacit transformations do
not change the observable correctness class.
\end{proposition}

The proof and full structure appear in Appendix~A. For the main text, it
suffices to know that correctness is \emph{localized} in tacit space and
is only observable via \(\gamma\).

\subsection{Uniqueness of the Phase--1 Ontology}

Appendix~A shows that, under reasonable axioms, the Phase--1 ontology is
essentially unique.

\begin{theorem}[Phase--1 initiality (informal)]
Any epistemic ontology that:
\begin{itemize}[leftmargin=2em]
  \item distinguishes explicit and tacit modalities,
  \item enforces a one-way boundary \(E\to T\),
  \item localizes correctness in tacit operators \(\TE,\HFD\), and
  \item represents AI systems as acting only in explicit space,
\end{itemize}
admits a unique structure-preserving mapping from the \FRFP{}
Phase--1 ontology. In this precise sense, the Phase--1 ontology is
minimal and essentially unique.
\end{theorem}

We refer the reader to Appendix~A for the formal category of Phase--1
frameworks and the proof that the FRFP ontology is initial in that
category.

%=========================================================
\section{Phase 2: Explicit Algebra and Navigation}
\label{sec:phase2}
%=========================================================

Phase~2 adds structure to explicit reasoning. The two main ingredients
are:
\begin{itemize}[leftmargin=2em]
  \item a Task Decomposition Grammar (\(\TDG\)) that generates explicit
        pipelines from the primitives \(\RI,\EC,\ED\); and
  \item a navigation groupoid (\(\Ncat\)) that captures reversible
        changes of perspective or representation.
\end{itemize}
These are combined via a Grothendieck-style semidirect product
\(\Ncat \ltimes \Ecat\), a standard construction in categorical
semantics~\cite{MacLane1998,Awodey2010,Borceux1994}, yielding an algebra
of reasoning trajectories.


\subsection{Task Decomposition Grammar}

TDG is a free algebra of explicit pipelines. Informally, its grammar is:
\begin{align*}
Q &::= \RI \mid Q;\EC \mid Q;\ED \mid
        \BRANCH(Q_1,\dots,Q_k) \mid
        \PARALLEL(Q_1,\dots,Q_k) \mid
        \LOOP(Q_{\mathrm{cond}},Q_{\mathrm{body}}),\\
P &::= Q;\RB;\TE;\HFD.
\end{align*}
Here \(Q\) ranges over explicit pipelines and \(P\) over complete FRFP
pipelines.

The free explicit category \(E_{\mathrm{sem}}\) generated by TDG has a
standard universal property: any interpretation of the generators in an
explicit category extends uniquely to a functor from \(E_{\mathrm{sem}}\).

\subsection{Navigation Groupoid}

Navigation steps---trying alternative models, adding or dropping
assumptions, restarting from a new shape---are represented as morphisms
in a groupoid \(\Ncat\). Objects of \(\Ncat\) are \emph{shapes} of
explicit reasoning; morphisms are structural rewritings between shapes,
freely inverted to form a groupoid.

Each navigation morphism \(n\in\Ncat\) acts on explicit pipelines
via a fiberwise functor \(n^\ast:E_S\to E_{S'}\) between fibers of a
fibration \(p:\Ecat\to\Scat\) over shapes. This makes \(E\) an indexed
category over \(\Ncat\).

\subsection{Semidirect Product \texorpdfstring{\(\Ncat \ltimes \Ecat\)}{N ⋉ E}}

The combined reasoning algebra is the Grothendieck construction
\(\Ncat \ltimes E\). A morphism in \(\Ncat \ltimes E\) is a pair
\((n,e)\) consisting of a navigation morphism \(n\in\Ncat\) and an
explicit arrow \(e\in\Ecat\) in the appropriate fiber. Composition
interleaves navigation and explicit reasoning according to:
\[
(n_2,e_2)\circ(n_1,e_1)
 =
\bigl(n_2n_1,\ e_2\circ n_2^\ast(e_1)\bigr).
\]

At an intuitive level:
\begin{itemize}[leftmargin=2em]
  \item \(\Ncat\) says \emph{which view} of the explicit reasoning we
        are in (e.g., which model, which set of assumptions);
  \item \(\Ecat\) records \emph{what} explicit computation and
        diagnostics have been applied;
  \item \(\Ncat \ltimes \Ecat\) packages these into a single object
        representing a reasoning trajectory.
\end{itemize}

\subsection{Confluence, Normal Forms, and Decidability}

Appendix~A equips \(\Ncat \ltimes E\) with a reduction system that
combines explicit rewrites and navigation moves. Provided that:
\begin{itemize}[leftmargin=2em]
  \item explicit reduction is terminating and locally confluent, and
  \item navigation reduction normalizes within each component of
        \(\Ncat\),
\end{itemize}
we obtain:

\begin{theorem}[Confluence and normal forms (informal)]
Every reasoning trajectory in \(\Ncat \ltimes \Ecat\) admits a
canonical normal form, unique up to navigation isomorphism. In
particular, explicit steps and navigation steps may be permuted and
simplified without changing the canonical form.
\end{theorem}

Moreover, the question \emph{``Is this composite a valid FRFP
pipeline?''} is decidable, given a well-specified reduction system and
typing rules.

\begin{theorem}[Decidability of admissibility (informal)]
Given a composite term in \(\Ncat \ltimes \Ecat\), it is decidable
whether it corresponds to an admissible pipeline of the form
\(\RI \to (\EC/\ED)^\ast \to \RB \to \TE \to \HFD\).
\end{theorem}

These properties make \(\Ncat \ltimes \Ecat\) a suitable intermediate
language for representing and analyzing complex reasoning trajectories.

%=========================================================
\section{Semantic Correctness and Stability}
\label{sec:semantics}
%=========================================================

The combination of Phases~1 and~2 yields a semantic map from composite
trajectories to tacit states.

\subsection{Semantic Functor and Observable Correctness}

The semantic functor
\[
\llbracket - \rrbracket: \Ncat \ltimes E_{\mathrm{sem}}\to\Tcat
\]
interprets a trajectory by running its explicit part, applying the
boundary \(\RB\), and then applying tacit evaluation and final decision.
Observable correctness is given by \(\gamma(\llbracket - \rrbracket)\).

\begin{theorem}[Correctness preservation (informal)]
Correctness judgments are invariant under explicit reasoning,
navigation, and semidirect-product composition: for any admissible
trajectory \(\tau\) and any navigation morphism \(n\in\Ncat\),
\[
\gamma(\llbracket \tau \rrbracket)
 =
\gamma(\llbracket n\cdot \tau \rrbracket).
\]
\end{theorem}

In other words, changing the route by which we arrive at a result---so
long as we respect explicit typing and Phase--1/2 constraints---does not
change the tacit correctness judgment.

\subsection{Canonical Correctness Form}

Every admissible pipeline admits a canonical factorization:

\begin{theorem}[Canonical correctness form (informal)]
Any trajectory in \(\Ncat \ltimes E_{\mathrm{sem}}\) that is admissible
in the FRFP sense is equivalent, up to navigation, to a unique normal
form in which:
\begin{itemize}[leftmargin=2em]
  \item all explicit operations (including navigation) occur before the
        boundary morphism \(\RB\);
  \item the tacit suffix is \(\TE \to \HFD\), with no interleaving
        explicit operations;
  \item no additional boundary or tacit operators appear.
\end{itemize}
\end{theorem}

This expresses, at the algebraic level, the design principle that AI
systems may explore and manipulate explicit representations in many
ways, but correctness judgments always arise from a single passage
through \(\RB\) followed by tacit evaluation and decision.

%=========================================================
\section{Phase 3: Dynamic Semantics and Hallucination}
\label{sec:phase3-dynamics}
%=========================================================

The structures introduced in Phases~1 and~2 characterise \emph{single}
FRFP pipelines: how explicit and tacit operations are separated, how
explicit reasoning and navigation combine, and how correctness is
preserved under admissible transformations. In typical scientific use,
however, a human agent engages with an AI system repeatedly, often over
long horizons. The tacit context from one episode influences the next,
and errors or misplaced trust may accumulate.

Phase~3 extends the framework to this setting by equipping FRFP
pipelines with a simple dynamic semantics on tacit states. The aim is
to capture how tacit grounding, task requirements, and confidence
interact over time, and to make precise in what sense hallucinations
become structurally difficult to avoid. The formal development of this
dynamic semantics is given in Appendix~\ref{app:pipeline-dynamics}; the
present section outlines the main ingredients and results.

\subsection{Tacit Context, Degradation, and Requirements}

We view the human agent's epistemic background as a tacit context space
\(T\) equipped with a preorder \(\preceq\):
\[
t_1 \preceq t_2
\quad\text{``\(t_1\) is epistemically no stronger than \(t_2\)''}.
\]
Strict inequality \(t_1 \prec t_2\) denotes \(t_1 \preceq t_2\) and
\(t_2 \not\preceq t_1\). No further algebraic structure is imposed on
\(T\) at this stage; the categorical structure of tacit space from
Appendix~\ref{app:math-foundations} is still present but the dynamic
analysis works at the level of the underlying preordered set.

\paragraph{Context degradation.}
Long-run use of an AI assistant is assumed to induce epistemic
degradation. This is modelled by a function \(\delta:T\to T\) satisfying:
\begin{enumerate}[label=(D\arabic*),leftmargin=2.2em]
  \item \emph{Monotonicity:} \(\delta(t) \preceq t\) for all \(t\in T\);
  \item \emph{Strict decay somewhere:} there exists \(t\) with
        \(\delta(t) \prec t\);
  \item \emph{Non-invertibility:} there is no \(f:T\to T\) such that
        \(\delta\circ f = \mathrm{id}_T\).
\end{enumerate}
Intuitively, \(\delta\) aggregates effects such as fatigue, forgetting,
and inappropriate deference to the system. Repair operations (rest,
study, institutional checks) can be represented by separate operators,
but are not needed for the core Phase~3 results and are therefore left
for later extensions.

\paragraph{Requirements and hallucination.}
Hallucination is defined at the level of the protocol, rather than in
terms of any particular model. To this end, we introduce a requirement
map
\[
\Req : E \times \mathbb{R}_{\ge 0} \to T,
\]
where \(e\in E\) is an explicit artifact (e.g., an explanation or
derivation) and \(c\ge 0\) is a tacit credence level indicating how
strongly the agent treats \(e\) as reliable. The element \(\Req(e,c)\)
represents the tacit grounding required to treat \((e,c)\) as
epistemically adequate.

We assume:
\begin{enumerate}[label=(R\arabic*),leftmargin=2.2em]
  \item \emph{Monotonicity in credence:} if \(c_1 \le c_2\) then
        \(\Req(e,c_1) \preceq \Req(e,c_2)\) for all \(e\in E\);
  \item \emph{Nontriviality:} not all artifacts have the same
        requirements up to \(\preceq\).
\end{enumerate}

\begin{definition}[Hallucination]
Let \(t\in T\) be the current tacit state, \(e\in E\) an explicit output,
and \(c\ge 0\) a credence level. A \emph{hallucination} occurs at this
step if
\[
t \prec \Req(e,c).
\]
We then say that \(e\) is \emph{hallucinated} (relative to \((t,c)\)).
\end{definition}

In this sense, hallucination is a mismatch between available tacit
grounding and what would be required to support the explicit artifact at
the assigned credence. This definition does not depend on internal model
mechanisms or intent.

\subsection{Pipeline Transformers and Runs}
\label{subsec:phase3-runs}

The static semantics of Appendix~\ref{app:math-foundations} already
associates to each explicit pipeline \(e\in E_{\mathrm{sem}}\) a tacit
effect via the semantic RB functor \(\RBsharp\) and the tacit primitives
\(\TE,\HFD\). In particular, evaluation of \(e\) at tacit context \(t\)
can be written as
\[
t \longmapsto \HFD\circ \TE \circ \RBsharp(e,t).
\]

For dynamic analysis, we compose this with degradation.

\begin{definition}[Pipeline transformer]
A complete pipeline at a given turn consists of:
\begin{itemize}[leftmargin=2em]
  \item an explicit artifact \(e \in E_{\mathrm{sem}}\) constructed as in
        Phases~1–2;
  \item the tacit suffix \(\RB \to \TE \to \HFD\);
  \item a degradation step \(\delta:T\to T\).
\end{itemize}
It induces a map
\[
\llbracket P \rrbracket_\delta : T \to T,\qquad
\llbracket P \rrbracket_\delta(t)
 :=
 \delta\bigl(\HFD\circ\TE\circ\RBsharp(e,t)\bigr).
\]
\end{definition}

\begin{definition}[Run]
A \emph{run} (or \emph{interaction sequence}) is a sequence of tacit
states \((t_n)_{n\ge 0}\) such that, for some sequence of pipelines
\(P_0,P_1,\dots\),
\[
t_{n+1} = \llbracket P_n \rrbracket_\delta(t_n)
\quad\text{for all } n\ge 0.
\]
When convenient, we also record the corresponding explicit artifacts and
credences, writing \((t_n,e_n,c_n)_{n\ge 0}\).
\end{definition}

All Phase~3 conclusions are stated in terms of such runs and the
hallucination condition \(t_n \prec \Req(e_n,c_n)\).

\subsection{Baseline Hallucination Inevitability}

We first consider a deterministic setting with fixed degradation
\(\delta\) and fixed pipelines. To capture the notion that nontrivial
epistemic demands recur, we use a requirement floor.

\begin{definition}[Requirement floor and recurring demands]
A state \(r\in T\) is a \emph{requirement floor} if, for a run
\((t_n,e_n,c_n)_{n\ge 0}\), we say that \emph{nontrivial demands recur
relative to \(r\)} when there exists an infinite subsequence
\((n_k)_{k\ge 0}\) such that
\[
\Req(e_{n_k},c_{n_k}) \succeq r
\quad\text{for all }k.
\]
\end{definition}

We strengthen the degradation assumption to ensure that tacit states
eventually fall below this floor.

\begin{assumption}[Degradation below the requirement floor]
\label{assump:baseline-degradation-main}
There exists \(r\in T\) such that, for every \(t_0\in T\), there is
\(N\in\mathbb{N}\) with
\[
\delta^N(t_0) \prec r.
\]
\end{assumption}

Under these assumptions, Appendix~\ref{app:pipeline-dynamics} shows:

\begin{theorem}[Hallucination inevitability, deterministic case]
\label{thm:baseline-inevitability-main}
Let \((t_n,e_n,c_n)_{n\ge 0}\) be a run for which
Assumption~\ref{assump:baseline-degradation-main} holds and
nontrivial demands recur relative to some requirement floor \(r\). Then
there exists \(n\in\mathbb{N}\) such that
\[
t_n \prec \Req(e_n,c_n).
\]
In particular, every infinite run of this form contains at least one
hallucination.
\end{theorem}

The theorem applies to runs that are both degrading and nontrivial in
the sense above. Scenarios restricted to very easy tasks, very short
horizons, or aggressive external re-grounding may not satisfy these
premises and are not targeted here.

\subsection{Confidence, Requirement Escalation, and Feedback}

The behaviour of the human agent depends not only on the explicit
content \(e_n\) but also on how plausible it appears. Phase~3 therefore
adds structure for explicit plausibility, credence, and feedback between
credence and degradation. The formal definitions and proofs are given in
Appendix~\ref{app:pipeline-dynamics}; we summarise the main ideas.

\paragraph{Plausibility and confidence.}
An explicit plausibility functional \(\pi:E\to[0,1]\) assigns to each
artifact a scalar measure of surface coherence or apparent authority.
Tacit state is extended to \(T^e:=T\times\mathbb{R}_{\ge 0}\), with
elements \((t,c)\) where \(t\) is grounding and \(c\) is credence. A
confidence inflation operator
\(\Phi:(T^e\times E)\to T^e\) updates credence according to plausibility,
without improving \(t\); for example
\[
\Phi((t,c),e) = (t,c+\varphi(\pi(e))),
\]
for some monotone \(\varphi:[0,1]\to\mathbb{R}_{\ge 0}\).

\paragraph{Requirement escalation and feedback-coupled degradation.}
Since \(\Req(e,c)\) is monotone in \(c\), higher credence raises
requirements. To reflect the possibility that higher confidence reduces
scrutiny, degradation is allowed to depend on \(c\) via a map
\(\delta^\star:T^e\to T\), written \(\delta^\star(t,c)=\delta_c(t)\),
such that:
\begin{itemize}[leftmargin=2em]
  \item for fixed \(t\), higher \(c\) leads to weaker outputs:
        \(\delta^\star(t,c_2)\preceq \delta^\star(t,c_1)\) whenever
        \(c_1\le c_2\);
  \item at zero credence, \(\delta^\star\) reduces to the baseline
        \(\delta\).
\end{itemize}

An extended run \((t_n,c_n,e_n)_{n\ge 0}\) is then obtained by
interleaving confidence inflation and feedback-coupled degradation. Under
assumptions that credence can grow without bound along subsequences when
high-plausibility artifacts recur, Appendix~\ref{app:pipeline-dynamics}
establishes an \emph{accelerated} version of
Theorem~\ref{thm:baseline-inevitability-main}: hallucinations still
occur in finite time, and the time-to-hallucination is not increased,
and may be strictly reduced, by the presence of the feedback loop.

\subsection{Stochastic Pipelines, Ensembles, and Finite Horizons}

The preceding discussion treats runs as deterministic sequences. In many
applications, FRFP pipelines include stochastic components (e.g.,
sampling, randomised retrieval). Appendix~\ref{app:pipeline-dynamics}
extends the dynamic semantics to this case by viewing runs as
trajectories \(\omega\) in a space \(\Omega\) and equipping \(\Omega\)
with probability measures arising from stochastic implementations.

For each trajectory \(\omega = ((t_n,c_n,e_n))_{n\ge 0}\), the
\emph{hallucination time} is defined as
\[
\tau(\omega)
 :=
\inf\{n\ge 0:\ t_n \prec \Req(e_n,c_n)\},
\]
with \(\tau(\omega)=+\infty\) if no hallucination occurs. When the
structural premises of Theorem~\ref{thm:baseline-inevitability-main}
(and its feedback-augmented variants) hold pointwise on the support of a
stochastic pipeline, Appendix~\ref{app:pipeline-dynamics} shows that
\(\tau(\omega)<\infty\) for all such \(\omega\). Any probability measure
supported on these trajectories then satisfies
\(\mathbb{P}(\tau<\infty)=1\): hallucination is almost sure.

For a fixed horizon \(N\), the finite-horizon hallucination probability
\(p_N:=\mathbb{P}(\tau\le N)\) is non-decreasing in \(N\) and converges
to~1 as \(N\to\infty\). For \(K\) independent runs, the probability that
at least one run hallucinates by time \(N\) is
\(1-(1-p_N)^K\). Thus ensembles can improve robustness over short
horizons, but they do not change the long-run structural inevitability
when the underlying conditions are met.

\subsection{Implications}

Phase~3 shows how the static FRFP architecture behaves over long
interaction sequences when combined with simple models of context
degradation, requirements, and confidence. The main conclusions are:
\begin{itemize}[leftmargin=2em]
  \item hallucination is naturally captured as a mismatch between tacit
        context and task requirements, independent of any particular AI
        architecture;
  \item under monotone degradation and recurring nontrivial demands,
        hallucinations are inevitable along sufficiently long runs;
  \item feedback from plausibility and credence to degradation tends to
        shorten the time-to-hallucination rather than extend it;
  \item stochasticity and ensembles can reduce error on finite horizons,
        but do not remove the long-run structural pressure towards
        hallucination under the same assumptions.
\end{itemize}

The formal statements and proofs are given in
Appendix~\ref{app:pipeline-dynamics}. Subsequent sections extend the
framework from single agents to multi-agent and institutional settings.
% ------------------------------------------------------------------
% Section X — Collective Epistemic Dynamics
% ------------------------------------------------------------------

\section{Collective Epistemic Dynamics}
\label{sec:collective-epistemic-dynamics}

This section develops the collective-level component of the FRFP framework.
The focus is on epistemic systems comprising many agents who reason over a
shared space of explicit artifacts while retaining separate tacit states and
degradation dynamics. The formal constructions parallel the single-pipeline
dynamics introduced earlier for individual agents and pipelines, and extend
them to populations. Technical details and categorical formulations may be
found in the appendices, in particular Appendices~A--C.

\subsection{From Single Pipelines to Epistemic Populations}

In the single-agent setting, an FRFP pipeline comprises:
\begin{itemize}
  \item an explicit artifact space $E$ (texts, code, plans, proofs),
  \item a tacit space $T$ with a preorder $\preceq$ capturing ``at least as
        grounded as'',
  \item an extended tacit state space $T^e = T \times \mathbb{R}_{\ge 0}$
        pairing grounding and tacit credence,
  \item dynamics combining degradation $\delta$, requirement maps
        $\mathrm{Req}$, confidence inflation, and feedback-coupled
        degradation.
\end{itemize}
A single execution is a trajectory
\[
  \omega = \bigl( (t_n,c_n,e_n) \bigr)_{n \ge 0},
\]
and hallucination occurs at the first $n$ such that
$t_n \prec \mathrm{Req}(e_n,c_n)$ while the explicit artifact $e_n$ is
nevertheless accepted in the pipeline.

To move to the collective level, we replace the single tacit space by a
family:
\[
  A_i = (T_i, \preceq_i, \delta_i, T^e_i),
  \qquad i \in I,
\]
of epistemic agents indexed by a set $I$ (which may represent individuals,
teams, or institutions). Each $A_i$ has its own tacit space, requirement
relation, and degradation dynamics. Agents share a common explicit
reasoning environment described by the category
\[
  \mathcal{C} = N \ltimes E,
\]
where $E$ collects explicit artifacts and $N$ captures navigation and
control structure at the explicit level (task selection, mode switching,
tool invocation, and so on).

The joint explicit state of the population is modeled as an object of the
product category
\[
  \mathcal{C}^I.
\]
An object $X \in \mathcal{C}^I$ is a family $X = (X_i)_{i\in I}$ of
explicit states, indicating for each agent which explicit task and artifacts
are currently active. A morphism $X \to Y$ in $\mathcal{C}^I$ represents
one step of explicit activity across the population: one or more agents
update their explicit state, for example by computing, revising, delegating,
or producing new artifacts.

A population semantics functor
\[
  \overline{J} : \mathcal{C}^I \longrightarrow \prod_{i\in I} T^e_i
\]
associates to each global explicit state $X$ a tuple
$\overline{J}(X) = \bigl((t_i,c_i)\bigr)_{i\in I}$ of extended tacit
states, obtained by allowing each agent to consume the explicit artifacts
in $X$ and update according to the single-agent dynamics. This lifts the
explicit--tacit separation to the population level: explicit artifacts are
shared, while tacit grounding remains agent-local.

\subsection{Collective Dynamics as a Markov System}

To analyze temporal behavior, we endow $\mathcal{C}^I$ with stochastic
dynamics. From any global explicit state $X$, there is a distribution over
possible successors $Y$, reflecting:
\begin{itemize}
  \item which agents act next (possibly determined by a scheduling or
        institutional mechanism),
  \item which explicit morphisms in $\mathcal{C}$ they apply (such as
        computation, summarization, deployment, or consensus),
  \item which agents observe the newly produced artifacts (a communication
        structure on $I$),
  \item how each tacit state $(t_i,c_i)$ is updated via representational
        backflow, degradation, confidence inflation, and feedback-coupled
        degradation.
\end{itemize}

Formally, this is captured by a Markovian collective dynamics
\[
  M : \mathrm{Obj}(\mathcal{C}^I)
    \longrightarrow \mathcal{P}\bigl( \mathrm{Obj}(\mathcal{C}^I) \bigr),
\]
which specifies for each $X$ a probability measure $M(X,\cdot)$ on successor
states. Equivalently, one can view $M$ as a functor
\[
  M : \mathcal{C}^I \longrightarrow \mathbf{Kl}(D),
\]
into the Kleisli category of the discrete distribution monad $D$, but the
object-level transition kernel suffices for present purposes.

A \emph{collective execution} is a random path
\[
  X_0, X_1, X_2, \dots
\]
in $\mathcal{C}^I$ generated by $M$ from an initial state $X_0$. The
corresponding tacit trajectories $\overline{J}(X_n)$ record how grounding
and credence evolve over time for each agent, coupled through shared
explicit artifacts and communication. This yields a population-level path
space and probability measure analogous to the single-agent setting, and
permits the usual notion of stopping time with respect to the resulting
filtration (see Appendix~C for formal details).

\subsection{Collective Hallucination and Its Inevitability}

At the individual level, hallucination is defined as the first time an
explicit artifact is endorsed despite insufficient tacit grounding. At the
collective level, the corresponding notion involves agreement among a
subset of agents.

Given a collective execution $(X_n)_{n\ge 0}$ with tacit states
$\overline{J}(X_n) = \bigl((t_{i,n},c_{i,n})\bigr)_{i\in I}$, a
\emph{collective hallucination} at time $n$ occurs if there exists an
explicit artifact $e \in E$ and a nontrivial subset $J \subseteq I$ such
that:
\begin{itemize}
  \item each agent $j \in J$ has accepted or propagated $e$ in their
        pipeline by time $n$, and
  \item for all $j \in J$,
        \[
          t_{j,n} \prec_j \mathrm{Req}_j(e,c_{j,n}),
        \]
        i.e.\ for each such agent, the tacit state lies strictly below the
        requirement associated with $e$ in that context.
\end{itemize}

The \emph{collective hallucination time} is the stopping time
$\tau^{\mathrm{coll}}$ defined by
\[
  \tau^{\mathrm{coll}}(\omega)
  := \inf\{\,n \ge 0 : \omega \text{ exhibits a collective hallucination at
      step } n\,\},
\]
with $\inf \varnothing = \infty$ as usual.

Under the same high-level assumptions that drive single-agent inevitability
(finite tacit capacities, non-invertible degradation for each agent,
recurrent demands on grounding via surface-plausible artifacts, confidence
inflation, and requirement escalation), and under mild assumptions on $M$
(continued production and circulation of artifacts through the population),
one obtains an almost-sure finiteness result:
\[
  P_M\bigl(\tau^{\mathrm{coll}} < \infty\bigr) = 1.
\]
In words, over sufficiently long horizons, some under-grounded artifact
will almost surely be collectively endorsed, regardless of how the
explicit-only interaction is arranged, provided the structural FRFP
conditions remain satisfied. A formal statement and proof are given in
Appendix~C.

\subsection{Consensus Mechanisms as Explicit-Only Functors}

Collective systems frequently employ consensus mechanisms: voting schemes,
averaging procedures, ensemble methods, committee decisions, and more
elaborate governance structures. Within the present framework, these are
modeled as explicit-only operators:
\[
  C : E^n \longrightarrow E,
\]
which aggregate a finite tuple of explicit artifacts
$(e_1,\dots,e_n) \in E^n$ into a new explicit artifact $C(e_1,\dots,e_n)$.
Crucially, such operators act solely in explicit space; they have no direct
access to tacit states $T_i$ or to the requirement maps $\mathrm{Req}_i$.

A basic property of these operators is \emph{tacit-stability}. Fix an agent
$i$ and a context $c_i$. Suppose that for each $k \in \{1,\dots,n\}$,
\[
  t_i \prec_i \mathrm{Req}_i(e_k,c_i),
\]
so that each input artifact is under-grounded for agent~$i$ in context
$c_i$. Because explicit transformations do not modify the underlying tacit
state, any output $C(e_1,\dots,e_n)$ remains evaluated against the same
$t_i$ and $c_i$, and therefore also fails to meet the requirement. In
particular, explicit-only consensus cannot convert a collection of
under-grounded artifacts into a well-grounded one for a given agent.

As a consequence, fully automated consensus processes that operate purely
on explicit artifacts cannot, in general, eliminate hallucination or
under-grounded decisions at the collective level. They can smooth,
denoise, or delay failures by reshaping explicit behavior, but they cannot
create new grounding in tacit space. This extends the single-agent
observation that explicit refactorings do not restore grounding to the
population level.

\subsection{Structural Notions at the Collective Level}

The collective formalism introduces several structural concepts that do not
appear at the static single-agent level, and which will be used in later
sections.

\paragraph{Equivalence of collective protocols.}

Viewing $\mathcal{C}^I$ as the explicit category for the full population,
different interfaces, workflows, or consensus arrangements can be
represented as endofunctors
\[
  F : \mathcal{C}^I \longrightarrow \mathcal{C}^I
\]
that preserve the population semantics $\overline{J}$. If two designs
$F_1$ and $F_2$ can be related by such explicit-only refactorings that
leave $\overline{J}$ unchanged, they are \emph{equivalent} in the FRFP
sense: they have the same tacit semantics and the same distribution of the
collective hallucination time $\tau^{\mathrm{coll}}$. In that case, any
difference between them is purely syntactic and does not correspond to a
structural epistemic improvement.

\paragraph{Refinement order on architectures.}

Given two Markovian collective dynamics $M_1$ and $M_2$ defined on the
same population, with corresponding collective hallucination times
$\tau^{\mathrm{coll}}_1$ and $\tau^{\mathrm{coll}}_2$, define the
finite-horizon risk curves
\[
  p_N^{(k)} := P_{M_k}\bigl(\tau^{\mathrm{coll}}_k \le N\bigr),
  \qquad k \in \{1,2\},\ N \in \mathbb{N}.
\]
One may say that $M_2$ \emph{dynamically refines} $M_1$ if
$p_N^{(2)} \le p_N^{(1)}$ for all $N$. This induces a partial order on
collective designs: some architectures are strictly safer than others in
the sense of never increasing finite-horizon hallucination risk. Explicit
refactorings that preserve $\overline{J}$ cannot strictly improve this
order; they are at best equal. In contrast, interventions that alter tacit
structure (e.g.\ adding independent evaluators, modifying requirements, or
slowing degradation) can produce genuine refinements. A more detailed
treatment appears in Appendix~C.

\paragraph{Compositional risk and defense in depth.}

Many systems can be decomposed into interacting subpopulations (for
example, scientific, engineering, and policy components). Each
subpopulation has its own collective hallucination time
$\tau^{\mathrm{coll}}_{\mathrm{sub}}$ and hazard profile. At a coarse
level, the overall collective hallucination time is bounded by
\[
  \tau^{\mathrm{coll}} = \min_{\mathrm{sub}}
  \tau^{\mathrm{coll}}_{\mathrm{sub}}.
\]
This reflects the fact that the system as a whole fails epistemically as
soon as any critical component fails. The formalism thus supports a
compositional ``defense-in-depth'' analysis: subsystem-level guarantees on
hazard profiles may, under appropriate coupling conditions, yield
system-level guarantees. The appendices provide the technical apparatus
needed to make such statements precise.

\paragraph{Non-identifiability from logs.}

Collective executions give rise to rich explicit traces: logs, metrics,
dashboards, and derived indicators. All of these live in
$\mathcal{C}^I$ or in its images under explicit functors. Tacit states
$\overline{J}(X)$ and requirement maps $\mathrm{Req}_i$ remain unobserved.
In general, there exist many different tacit configurations and dynamics
$M$ that induce the same (or arbitrarily similar) distributions over
explicit traces. It follows that correctness and hallucination risk are
\emph{non-identifiable} from logs alone: without additional assumptions,
explicit-only monitoring and metrics cannot uniquely determine how well
grounded a population's beliefs are. This mirrors the observation in other
domains (e.g.\ markets or large-scale experiments) that transaction or
event logs do not uniquely fix underlying beliefs or theories.

\paragraph{Scaling and phase-transition questions.}

The population-level formulation also provides a natural language for
scaling questions. By parameterizing systems by population size $|I|$,
communication structure, and per-agent degradation properties, one can ask
how epistemic risk scales with the number of agents and tools, whether
particular communication patterns are stabilizing or destabilizing, and
whether qualitative changes in behavior (phase transitions) occur as
complexity grows. The present work does not attempt a systematic analysis
of scaling behavior, but the structural notions above identify the relevant
objects for future investigation.

\subsection{Connections to Example Domains and Local Interventions}

The same collective apparatus applies across domains in which many agents
interact over shared artifacts:

\begin{itemize}
  \item Scientific communities can be modeled as populations of agents
        (labs, groups, or individual scientists) with papers, datasets, and
        code as explicit artifacts, and peer review, replication, and
        citation practices as partial consensus mechanisms.
  \item Software development lifecycles can be treated as populations of
        developers, reviewers, operators, and tools with explicit artifacts
        such as requirements, designs, code, tests, and deployment
        configurations, coupled by CI/CD and incident-response processes.
  \item Strategic domains such as military planning or pandemic response
        involve agents in governments, institutions, and advisory bodies,
        operating over explicit artifacts including plans, intelligence,
        models, and policies, with doctrines and institutions acting as
        higher-level operators on $M$ and on $\mathrm{Req}$.
  \item Markets and large-scale experiments provide examples where rich
        explicit traces (transaction logs, event data) remain insufficient
        to uniquely identify underlying beliefs or theories, paralleling
        the non-identifiability observations above.
\end{itemize}

The collective perspective also clarifies the status of local interaction
patterns, such as common prompting schemes for language models (incentive
framing, challenge prompts, stepwise guidance, high-stakes framing,
self-evaluation prompts). These can be viewed as local interventions on:
\begin{itemize}
  \item the explicit pipeline shape in $N \ltimes E$ (for example,
        encouraging decomposition and self-checking),
  \item the effective requirement schedule (for example, weakening answers
        to hypotheses or raising the bar for strong claims),
  \item the human user's tacit dynamics (for example, prompting increased
        attention and skepticism).
\end{itemize}
Such interventions can materially improve finite-horizon correctness and
enlarge safe horizons for particular tasks. They do not, however, alter the
structural facts established in the underlying theory: explicit-only
systems remain ungrounded, human tacit capacities remain finite and
degrading, and, over unbounded horizons, hallucination and under-grounded
collective decisions remain inevitable.

In summary, the collective layer provides a coherent extension of FRFP to
multi-agent settings. It connects individual pipeline dynamics to
institutional design, clarifies the limits of explicit-only interventions
(including consensus mechanisms and monitoring), and identifies the
structural levers---in particular, those acting on tacit dynamics and
requirements---on which robust long-horizon human--AI collaboration must
ultimately rely.
\section{Structural Limits of Explicit Mechanisms}
\label{sec:structural-limits-explicit}

The preceding development has described a class of human--AI collaborations in which
\begin{itemize}
  \item explicit artefacts and transformations are modelled in a category $E$ of explicit reasoning,
  \item navigation, prompting, and other meta-operations are represented in a category $N$,
  \item tacit states and evaluations live in a space $T$, and
  \item the combined system $N \ltimes E$ is equipped with a reduction relation and a semantic functor
    \[
      \llbracket - \rrbracket : N \ltimes E_{\mathrm{sem}} \longrightarrow T,
    \]
    with observable correctness extracted via a map $\gamma : \mathrm{Ob}(T) \to J$.
\end{itemize}
Under the FRFP axioms (ETS, HEG, HEC, AE), explicit computation is separated from tacit evaluation and cannot reconstruct tacit state or correctness except via designated human interaction. At the same time, the dynamic and collective layers developed earlier show that long-horizon behaviour is inherently stochastic: pipelines are implemented as stochastic processes over a space of admissible trajectories, and populations interacting through explicit channels induce Markovian dynamics over collections of agents.

This section explains how these ingredients jointly constrain what can be achieved by mechanisms that operate purely in explicit space. Formal statements and proofs of the main structural results are given in Appendix~D.

\subsection{Dynamic Quantities and Long-Horizon Behaviour}

When a fixed pipeline $P$ is executed under FRFP constraints, the dynamic semantics introduced previously associate to any admissible implementation:
\begin{itemize}
  \item a measurable space of trajectories $\Omega_{\mathrm{allowed}}$ of runs, each trajectory $\omega \in \Omega_{\mathrm{allowed}}$ corresponding to a reduction sequence in $N \ltimes E$ satisfying the Phase--0 constraints;
  \item a hallucination time random variable
    \[
      \tau : \Omega_{\mathrm{allowed}} \longrightarrow \mathbb{N} \cup \{\infty\},
    \]
    recording the first time a requirement-violating artefact is accepted along the run;
  \item finite-horizon hallucination probabilities
    \[
      p_N \;=\; \mathbb{P}(\tau \leq N)
    \]
    for $N \in \mathbb{N}$;
  \item a hazard sequence $(h_n)_{n \in \mathbb{N}}$ derived from the increments of $(p_N)$; and
  \item for each tolerance $\varepsilon \in (0,1)$, a safe-horizon functional
    \[
      H_P(\varepsilon) \;=\; \sup\{N \in \mathbb{N} : p_N \leq \varepsilon\}.
    \]
\end{itemize}
These objects summarise, in a long-horizon sense, how often and how soon a given implementation of $P$ will fail the tacit requirements encoded by $\Req$.

Under the Phase--3 assumptions (finite tacit capacity, monotone non-invertible degradation of tacit state, recurring nontrivial demands, and appropriate independence conditions), one obtains an inevitability theorem: along admissible trajectories, hallucination occurs almost surely in the limit. In particular, $\mathbb{P}(\tau < \infty) = 1$ for any such implementation, and for every fixed $\varepsilon \in (0,1)$ the safe horizon $H_P(\varepsilon)$ is almost surely finite.

These results are proved in the dynamic layer but are agnostic to the internal details of the explicit mechanisms; they depend only on the interaction between explicit artefacts, requirement maps, and degradation behaviour of tacit state.

\subsection{Population-Level Dynamics and Consensus}

When multiple agents interact through explicit channels, the population development introduces:
\begin{itemize}
  \item an explicit reasoning category $\mathcal{C}$ for individual agents;
  \item a product category $\mathcal{C}^I$ of explicit profiles across a population $I$;
  \item functorial consensus operators
    \[
      C : \mathcal{C}^n \longrightarrow \mathcal{C},
    \]
    modelling explicit-only aggregation rules (voting, averaging, reconciliation of drafts, and so on); and
  \item a Markovian collective dynamics $M$ on $\mathcal{C}^I$, obtained by iterating local update and interaction steps.
\end{itemize}
Given an epistemic population (agents, tacit spaces, and requirement maps) and such a dynamics $M$, one can define a collective hallucination time $\tau_{\mathrm{coll}}$ as the first step at which a nontrivial subset of agents jointly accept an artefact that is under-grounded relative to their individual requirements.

Under assumptions analogous to those used for single-pipeline dynamics, one again obtains an almost-sure inevitability result: for admissible populations and dynamics consistent with the FRFP constraints, the collective stopping time $\tau_{\mathrm{coll}}$ is almost surely finite. Moreover, the almost-sure finiteness of $\tau_{\mathrm{coll}}$ is invariant under inserting additional explicit-only consensus layers; functorial consensus operators act purely in explicit space and cannot, by themselves, change the inevitability class of the process.

These collective results provide a quantitative description of long-horizon epistemic risk at the level of whole institutions or populations, using exactly the same kind of stochastic and categorical apparatus as in the single-pipeline case.

\subsection{Tacit Dependence and Explicit-Only Mechanisms}

Against this backdrop, it is natural to ask which desirable properties one might hope to enforce by purely explicit means. In the present setting, explicit-only mechanisms are those whose internal computation and inputs/outputs are confined to $E$ (or $\mathcal{C}$ in the population case), and which do not gain access to tacit states or tacit evaluations other than through the designated interfaces already modelled in the semantics.

A recurring theme is that many natural predicates of interest are \emph{tacit-dependent}: their truth value can vary between admissible runs that share the same explicit artefacts and explicit interaction history but differ in their tacit evolution. Examples include:
\begin{itemize}
  \item whether a given explicit artefact meets the relevant tacit requirements at the point of use;
  \item whether hallucination occurs within a fixed horizon, as expressed by events of the form $[\tau \leq N]$;
  \item whether an implementation has a given $\varepsilon$-safe horizon, as expressed by inequalities $H_P(\varepsilon) \geq N$; and
  \item in the collective setting, whether a population remains hallucination-free up to time $N$, as expressed by events such as $[\tau_{\mathrm{coll}} > N]$.
\end{itemize}
Because the combined explicit trace does not uniquely determine the underlying tacit trajectory, predicates of this kind cannot be reduced to functions of explicit artefacts alone. This observation is formalised in Appendix~D, which defines tacit-dependent predicates in terms of the semantic functor $\llbracket - \rrbracket$ and observable correctness $\gamma(\llbracket P \rrbracket)$ and then studies their interaction with explicit-only mechanisms.

The central structural result, stated and proved there, shows that under ETS and AE, no mechanism confined to explicit space can enforce a tacit-dependent predicate across all admissible runs. The proof relies only on the fact that such a mechanism receives isomorphic explicit inputs across certain pairs of admissible runs whose tacit behaviour differs, and hence must behave identically even when the predicate of interest disagrees.

\subsection{Consequences for Detection, Elimination and Consensus}

Several consequences follow immediately when this schema is instantiated with the dynamic and collective predicates arising from earlier sections.

\paragraph{Detection and elimination.}
The definition of hallucination in terms of requirement maps and tacit state yields a predicate of the form “this explicit artefact is hallucinated” which is tacit-dependent: the same explicit output can be acceptable or unacceptable depending on the tacit configuration in which it is evaluated. It follows that:
\begin{itemize}
  \item there is no perfect hallucination detector operating purely on explicit artefacts and logs; and
  \item no explicit-only policy can guarantee that hallucination never occurs along admissible runs, or that $\tau = \infty$ (or $\mathbb{P}(\tau = \infty) = 1$) holds uniformly across implementations supported on $\Omega_{\mathrm{allowed}}$.
\end{itemize}
Under the Phase--3 assumptions, inevitability theorems already show that $\tau$ is almost surely finite for admissible implementations, and that $H_P(\varepsilon)$ is almost surely finite for each fixed $\varepsilon \in (0,1)$. The structural results in Appendix~D complement this by ruling out the possibility of using explicit-only mechanisms to enforce $\tau = \infty$ or $H_P(\varepsilon) = \infty$ as design constraints.

\paragraph{Correctness oracles.}
Observable correctness is represented by $\gamma(\llbracket P \rrbracket)$ and is defined purely in tacit space, subject to HEG and HEC. For a fixed explicit artefact, different admissible runs can lead to different tacit evaluations and hence different observable correctness values. A hypothetical correctness oracle
\[
  \Omega : E \longrightarrow J
\]
that recovered $\gamma(\llbracket P \rrbracket)$ from the explicit artefact alone would thus be required to decide a tacit-dependent predicate using only explicit data. The formal results establish that no such oracle can exist within the FRFP axioms if it is confined to explicit computation.

\paragraph{Consensus and collective guarantees.}
In the population setting, consensus operators $C : E^n \to E$ and explicit coordination rules act entirely in $E$ and do not modify tacit states directly. They can be composed to form consensus pipelines and embedded in the Markov dynamics $M$ on population states. The structural results then imply that:
\begin{itemize}
  \item explicit-only consensus cannot create grounding ex nihilo: if each input artefact exceeds the grounding of an agent, so does the consensus output for that agent;
  \item no fully automated consensus pipeline can guarantee that its outputs are hallucination-free across all admissible populations and implementations; and
  \item more strongly, predicates such as $[\tau_{\mathrm{coll}} = \infty]$ are tacit-dependent and cannot be enforced by explicit-only transformations of the population dynamics.
\end{itemize}
Invariance results for $\tau_{\mathrm{coll}}$ in the collective dynamics further reinforce this picture: inserting additional explicit-only consensus layers can alter the distribution of collective risk but cannot change the fact that $\tau_{\mathrm{coll}} < \infty$ almost surely under the stated assumptions.

\subsection{Implications for Design and Governance}

Taken together, these structural limitations delineate the scope of what explicit-only design can accomplish in long-horizon human--AI collaboration.

On the one hand, the dynamic layer provides meaningful quantitative levers:
\begin{itemize}
  \item implementations can be compared via their finite-horizon risk profiles $(p_N)$ and hazard sequences $(h_n)$;
  \item the safe-horizon functional $H_P(\varepsilon)$ captures, for each tolerance $\varepsilon$, how long a system can typically operate before crossing a risk threshold;
  \item dynamic refinement orders allow one to classify implementations as safer or riskier over finite horizons.
\end{itemize}
On the other hand, the impossibility results anchored in tacit-dependence show that no amount of explicit-only filtering, logging, or aggregation can convert a system whose long-run behaviour is almost surely fallible into one that is guaranteed safe across all admissible runs and populations.

Within this boundary, the design problem becomes one of shaping hazard profiles and safe horizons, structuring interactions so that collective dynamics refine a baseline rather than degrade it, and introducing additional forms of tacit evaluation and institutional oversight. Subsequent sections on governance, incentives, and institutional design can be read as exploring these remaining degrees of freedom: adjusting requirement maps, allocating human attention, and configuring collective processes so as to manage long-horizon risk in a way that is consistent with the FRFP axioms and with the structural limitations just described.

Formal statements and proofs of the general schema and its corollaries appear in Appendix~D; the present discussion is intended to make clear how those results integrate with the dynamic and collective layers and how they bear on the design of real-world human--AI collaborations.
\section{Institutions for Long-Horizon Human--AI Collaboration}
\label{sec:institutions-main}

The previous sections analyzed individual human--AI pipelines and their long-run behavior: how explicit pipelines acquire tacit grounding, how degradation and hallucination accumulate over time, and how collective epistemic dynamics emerge when many agents interact with shared systems. We now step up one more level, from collections of agents and runs to the \emph{institutions} that sit above them. The aim of this section is conceptual: to describe the role of institutions in stabilizing long-horizon practice, while leaving formal details to Appendix~E.

\subsection{From Individual Agents to Epistemic Populations}

Earlier, we treated a single human agent as a bearer of tacit context, capable of evaluating explicit artifacts and updating their internal state in response to interaction with an FRFP pipeline. At scale, however, we are never dealing with a single agent in isolation. Scientific communities, engineering teams, regulatory bodies, and user populations are all better modeled as \emph{epistemic populations}: structured collections of agents, each with their own tacit capacities and degradation dynamics.

Formally, we use the agent and population structures introduced in the collective dynamics analysis. Each epistemic agent carries:
\begin{itemize}
  \item a tacit context space $T_i$ with a preorder expressing ``at least as well grounded as'';
  \item a degradation operator $\delta_i : T_i \to T_i$ capturing background loss of tacit context;
  \item an extended tacit state $T^e_i$ that pairs tacit context with a notion of subjective credence.
\end{itemize}
For institutional purposes, we enrich this with a \emph{correctness quotient} $\gamma_i : T_i \to J$, where $J$ is the space of observable correctness judgments (accept/reject/revise, etc.). The tuple
\[
  (T_i, \preceq_i, \delta_i, T^e_i, \gamma_i)
\]
is the institutional ``signature'' of the agent: it records which tacit states they can occupy, how those states degrade, and how they express correctness judgments when asked.

An \emph{epistemic population} is then simply a family of such agents indexed by a set $I$, together with the morphisms that allow us to compare and transform populations. This lets us treat heterogeneous communities---with agents who differ in tacit capacity, degradation rates, and evaluative norms---within a single formal frame. Full categorical details are given in Appendix~E, but the upshot for the main story is that we can talk sensibly about how a fixed pipeline behaves when replicated across a population, and how institutions act on the resulting pattern of judgments.

\subsection{Replication and Population-Level Semantics}

Given a well-typed FRFP pipeline $P$, the earlier semantic development assigns it a tacit meaning $\sem{P}$ and an associated correctness outcome under evaluation by a single agent. In practice, however, we do not trust any single evaluation of $P$; we care about what happens when $P$ is executed across a population, at different times, under different tacit conditions.

To capture this, we consider \emph{admissible runs} of $P$ by each agent in the population. An admissible run bundles:
\begin{itemize}
  \item an explicit execution trace of $P$ in the navigation--explicit space $N \ltimes E$, subject to the structural constraints we have maintained throughout, and
  \item the induced evolution of the agent's extended tacit state along that trace, as specified by the Phase--3 update rules (degradation, confidence inflation, requirement escalation, and feedback-coupled degradation).
\end{itemize}

A \emph{replication family} for $P$ is then a choice of one admissible run for each agent (or for some subcollection of agents) in the population. For each such run, we select an evaluation time---for instance, when the pipeline halts or when a fixed resource limit is reached---and record the corresponding tacit state $t_{i,*}$ for agent $A_i$.

What matters for institutional behavior is not the detailed trajectory of each run, but the \emph{pattern of correctness judgments} these runs induce. We therefore define the \emph{population semantic profile} of $P$ relative to a replication family $R(P)$ as the multiset
\[
  \mathcal{P}(P; R(P)) \;=\; \{ \gamma_i(t_{i,*}) : i \in I_{\mathrm{rep}} \} \subseteq J.
\]
This multiset is the primary interface between long-horizon practice and institutions. It aggregates how many agents, under what tacit conditions, endorse, reject, or revise the outputs of a given pipeline. The construction makes clear that institutional input is inherently \emph{tacit-dependent}, even though it is summarized in a discrete space of judgments.

The formal development of admissible runs, replication families, and population semantic profiles is given in Appendix~E, building on the run semantics and stochastic models of Appendices~B and~C.

\subsection{Institutions as Aggregators over Tacit Judgments}

Institutions---journals, standards bodies, model review boards, internal red-teams, and so on---do not typically see the full detailed state of every agent. Instead, they receive some distribution of judgments: accept, reject, request further work, abstain. In our framework, an \emph{institutional aggregator} is exactly a map
\[
  I : M(J) \to J \cup \{\bot\},
\]
where $M(J)$ is the space of multisets of judgments and $\bot$ denotes suspended judgment. For a given pipeline $P$, an institution chooses a particular aggregator $I_P$ and applies it to the population semantic profiles generated when $P$ is replicated.

This abstraction covers common patterns:
\begin{itemize}
  \item majority or supermajority rules;
  \item unanimity or veto-based protocols;
  \item weighted consensus, where some agents' judgments count more than others;
  \item deliberative or staged procedures that can output ``no decision'' ($\bot$).
\end{itemize}

Crucially, $I_P$ operates \emph{only} on the multiset of tacitly produced judgments. It does not see raw model outputs, logs, or metrics directly. All of that explicit activity is mediated through human agents, whose tacit evaluations are then aggregated.

Full definitions of institutional aggregators, institutional outcomes, and their basic properties appear in Appendix~E. Here, the important point is that institutions sit at a distinct layer in the architecture: they neither compute the explicit artifacts themselves, nor embody the tacit processes that evaluate them. They are rules for turning many tacit judgments into collective decisions.

\subsection{Why Institutions Cannot Be Fully Automated}

Given mature AI systems and rich logging, it is tempting to imagine institutions that operate purely on explicit artifacts and their explicit summaries: dashboards, performance metrics, proofs of robustness, and so on. In our framework, such an institution would be \emph{explicit-only}: for each pipeline $P$, its decisions could be computed by a mechanism that inspects only explicit artifacts, never accessing the underlying tacit states or correctness judgments.

The impossibility results developed earlier rule this out in a strong sense. Under the assumptions that:
\begin{itemize}
  \item correctness is fundamentally tacit-dependent (it lives in the $\gamma_i$ quotients of tacit states), and
  \item admissible pipelines can be arranged so that explicit traces are ambiguous with respect to tacit correctness,
\end{itemize}
Appendix~E proves an institutional non-automation theorem: no explicit-only institutional layer can be \emph{correctness-preserving} across all admissible pipelines and populations. In particular, if we require that the institution agree with unanimous tacit verdicts whenever they occur, then there is no way to define its decisions purely in terms of explicit artifacts.

Intuitively, any such institution would amount to an ``explicit-only correctness oracle,'' and the earlier no-go theorems about correctness oracles apply. No amount of logs and metrics, however elaborate, can substitute for direct access to tacit evaluations. Institutions that aim at epistemic reliability must therefore be designed around \emph{tacit input}, even if they also make use of explicit signals.

The formal statement and proof of the institutional non-automation theorem are collected in Appendix~E, where they are derived as a corollary of the explicit-only impossibility results from the more general theory.

\section{Governance, Law, and Oversight}
\label{sec:governance-main}

Institutions regulate the behavior of individual pipelines and agents. But institutions themselves change over time: journals tighten or relax standards, internal review processes are redesigned, regulatory frameworks evolve. In long-horizon settings, these higher-order changes are at least as important as first-order institutional choices. This section outlines how governance and law fit into the architecture, again referring to Appendix~F for the formal layer.

\subsection{Governance as a Meta-Institution}

The collection of all institutional arrangements we might adopt forms a structured space. For each pipeline $P$ we can choose an aggregator $I_P$, and we may change these choices over time in response to experience. We write $\Inst$ for the category whose objects are such institutional families and whose morphisms are structure-preserving transformations between them.

A \emph{governance operator} is then an endofunctor
\[
  G : \Inst \to \Inst
\]
that describes how institutional arrangements evolve. At each time step $t$, a governance process takes an existing institutional layer $(I^{(t)}_P)_P$ and produces an updated one $(I^{(t+1)}_P)_P$. This update may depend on various observables derived from the collective dynamics of agents and pipelines: disagreement rates, error frequencies, time-to-detection for hallucinations, evidence of confidence inflation, and so on.

The collective dynamics themselves were modeled earlier as a process on epistemic populations interacting with pipelines. Appendix~F shows how to couple this with governance: observables are extracted from the evolving population, fed into a governance operator, and used to adjust institutional aggregators over time. The resulting picture is a closed loop:
\begin{enumerate}
  \item institutions aggregate tacit judgments and shape which pipelines are accepted or deployed;
  \item those pipelines, in turn, influence the tacit states of agents and the composition of populations;
  \item governance watches this evolving behavior and modifies institutions accordingly.
\end{enumerate}

This three-level interaction---agents and pipelines, institutions, and governance---is where the long-horizon behavior of human--AI systems is ultimately determined.

\subsection{Explicit-Only Governance and Its Limits}

Just as we can imagine explicit-only institutions, we can imagine explicit-only governance: processes that monitor logs, metrics, and other explicit artifacts of institutional performance, and then automatically adjust institutional rules based solely on those explicit signals.

In our framework, an explicit-only governance operator is one that, at each time step, updates the institutional layer using only explicit artifacts produced so far. Formally, Appendix~F captures this by requiring the updates to factor through explicit-only mechanisms on the explicit space $E$, in the same way the explicit-only institutions of Appendix~E factor through explicit artifacts.

The same tension appears here as before. If correctness is fundamentally tacit-dependent, and if pipelines and populations can be arranged so that explicit traces underdetermine tacit correctness, then governance that never consults tacit judgments cannot reliably distinguish good from bad institutional behavior. Any attempt to make governance fully explicit-only will either:
\begin{itemize}
  \item fail to track when institutions drift away from correct tacit practice, or
  \item implicitly smuggle in tacit judgments under the guise of ``ground truth'' labels or other human-provided signals.
\end{itemize}

Appendix~F formalizes this as a non-automation theorem for governance: under the same assumptions that rule out explicit-only correctness oracles and explicit-only institutional layers, there is no correctness-preserving governance operator that depends exclusively on explicit artifacts or their explicit summaries. In essence, a fully automated governance layer would again amount to a higher-order explicit-only oracle, and the impossibility results apply at this meta level as well.

\subsection{Law as a Constraint on Governance}

Governance does not operate in a vacuum. In practice, institutions and their evolution are constrained by legal and normative frameworks: data protection laws, safety regulations, due-process requirements, and so on. In the present setting, we model these as predicates on institutional arrangements and governance operators.

Concretely, we can view a legal or regulatory framework as specifying a set of \emph{admissible} institutions and transformations between them. A governance process is then \emph{lawful} if it never maps admissible institutions to inadmissible ones. Law can also constrain which observables governance is allowed to use (for example, privacy-preserving restrictions on what can be logged) and which kinds of interventions are permitted.

While the formal treatment of law as a constraint on $\Inst$ and on governance functors appears in Appendix~F, the conceptual takeaway is straightforward: the space of possible governance processes is smaller than the space of merely mathematically definable ones, and non-automation results for governance must be understood relative to whatever legal admissibility conditions we impose.

\subsection{Human Oversight as a Structural Requirement}

Taken together, the institutional and governance non-automation results point to a strong conclusion. It is not simply that, for now, we happen to rely on humans in the loop; rather, under the structural assumptions of the framework, \emph{there is no coherent way to remove them} while preserving correctness over long horizons.

At the institutional level, correctness-preserving aggregation requires access to tacit judgments. At the governance level, correctness-preserving adaptation of institutions requires some access to tacit evaluations of institutional performance. Attempts to fully automate either layer collapse, in the limit, into explicit-only oracles, which the earlier impossibility theorems prohibit.

This does not mean that institutions or governance must be informal or opaque. On the contrary, much of their structure can and should be formalized, monitored, and partially automated. The key point is that the formal and automated parts must be designed \emph{around} human tacit judgment, not as replacements for it.

The detailed statements and proofs of these claims are provided in Appendices~E and~F. Here we emphasize their architectural consequence: in any long-horizon human--AI ecosystem that satisfies the assumptions of the framework, human agents retain irreplaceable roles at three levels simultaneously:
\begin{enumerate}
  \item as individual evaluators of explicit artifacts;
  \item as contributors to institutional decisions via their tacit judgments; and
  \item as overseers of governance processes that adapt and constrain institutions over time.
\end{enumerate}
The rest of the paper can be read as an exploration of how to design pipelines, institutions, and governance mechanisms that make these roles precise rather than aspirational.

%=========================================================
\section{Discussion and Outlook}
\label{sec:discussion}
%=========================================================

\subsection{Design Implications}

The FRFP architecture formalizes a set of design constraints for AI
systems intended for long-horizon scientific collaboration:

\begin{itemize}[leftmargin=2em]
  \item \textbf{Explicit-only AI.} AI systems should manipulate explicit
        structures but must not assume tacit roles such as correctness
        judgment. This limits their epistemic authority while maximizing
        their computational utility.
  \item \textbf{Single explicit--tacit boundary.} All correctness-bearing
        information flow should pass through a clearly identifiable
        boundary operator \(\RB\), with no hidden channels from tacit
        to explicit space.
  \item \textbf{Auditability via canonical forms.} Reasoning trajectories
        can be normalized within \(\Ncat \ltimes \Ecat\), yielding
        canonical forms amenable to audit and replication.
  \item \textbf{Structural handling of hallucinations.} Hallucinations
        are treated as structural mismatches between context and
        requirements, not as psychological quirks of models. This allows
        analysis of how context degradation, feedback, and task design
        interact.
\end{itemize}

\paragraph{Kernel-level constraints vs.\ instantiated layers.}
The design constraints above attach first and foremost to the kernel of FRFP: the explicit/tacit separation with a single boundary, the minimal primitive set, and the canonical $N \ltimes E$ semantics with tacit correctness.
By construction, any system that respects this kernel inherits the same structural limitations and invariants, regardless of how it instantiates dynamics, populations, or institutions.
Conversely, Sections~6--10 and Appendices~B--F are deliberately more contingent: they work out one particular family of choices for context degradation, requirement schedules, stochastic pipelines, epistemic populations, institutional operators, and governance functors compatible with the kernel.
These outer layers are intended as illustrative and revisable; what is non-negotiable, from the FRFP point of view, is the underlying correctness architecture and its separation between explicit computation and tacit judgment.


\subsection{Limitations and scope}
\label{subsec:limitations}

The present work has several limitations.

\paragraph{Single-agent tacit model.}
We model one tacit context space $T$ at a time; multi-agent and institutional settings, where different humans (and possibly institutions) have distinct tacit states, require richer structures (e.g.\ families of tacit spaces, communication channels, and aggregation mechanisms). Developing such multi-agent extensions is left to future work.

\paragraph{Worst-case focus.}
The hallucination inevitability results rely on deliberately worst-case assumptions: long or unbounded runs, monotone degradation, no explicit repair operator, and simple forms of requirement escalation and feedback. Practical systems may use aggressive re-grounding, multiple agents, explicit restarts, or external tools (e.g.\ proof assistants, simulators) to reduce risk. Our results should therefore be read as structural ``upper bounds'' on what can go wrong in principle, not as quantitative predictions for particular deployments.

\paragraph{Abstract requirements and tacit semantics.}
The requirement map $\mathrm{Req}(e,c)$ is treated as an abstract semantic object; we do not attempt to estimate it from data or tie it to any specific scientific or engineering domain. Likewise, tacit states are modeled at a coarse level (often as elements of a preorder or low-dimensional space), omitting many cognitive and socio-technical factors that matter in practice. Making $\mathrm{Req}$ and the tacit dynamics empirically grounded for particular domains is outside the scope of this paper.

\paragraph{Level of abstraction.}
More broadly, we work at an intentionally abstract, foundational level and do not propose concrete algorithms, training procedures, or evaluations. FRFP is a design framework and semantic target, not a ready-made system. Instantiating it in concrete tools will require additional engineering choices (e.g.\ representations, interfaces, inference methods) that we do not model here.

\paragraph{Scope of novelty.}
As noted in the Introduction, the categorical ingredients we use (free categories, groupoids, fibrations, Grothendieck constructions) are standard; the contribution is to assemble them into a single architecture that captures explicit/tacit separation and feedback, and to derive system-level constraints on auditability, correctness, and hallucination. FRFP is intended as a reusable structural framework rather than a specific algorithm or training procedure.


\subsection{Future Directions}

Several directions are natural next steps.

\paragraph{Empirical dynamics of human--AI protocols.}
The current framework treats tacit context, degradation, requirements, confidence, and feedback as abstract operators. A first direction is to make these empirically meaningful. This includes: (i) developing task- and domain-specific proxies for tacit context, degradation, requirement maps, and plausibility functionals; (ii) instrumenting real human--AI workflows so that trajectories in \(N \ltimes E\) can be logged, normalized, and analysed; and (iii) testing structural predictions (such as inevitability and acceleration of hallucination) against behavioural and longitudinal data.

\paragraph{Repair, re-grounding, and reset mechanisms.}
Our analysis intentionally focuses on settings where context degrades and is not fully restored. A complementary direction is to design explicit repair and re-grounding mechanisms---study breaks, external reference checks, slow deliberation modes---and represent them as operators within the same algebraic framework. At the protocol level, this suggests specifying and analysing reset policies, usage caps, and re-verification gates as first-class constructions in \(N \ltimes E\), and characterising when such mechanisms significantly delay (though do not eliminate) structural hallucination.

\paragraph{Collective architectures and institutions.}
The framework can be extended from single users to populations of agents and to institutions that aggregate many runs. This opens questions about how different institutional patterns (peer review, adversarial review, replication, regulatory stress-testing) can be represented, and which combinations provably buffer individual degradation while respecting a human-exclusive grounding constraint. Another direction is to relate formal notions such as population semantic profiles, institutional stability, and institutional non-automation to empirical studies of scientific and regulatory processes, and to compare predicted failure modes (fake consensus, lock-in, correlated error) with observed ones.

\paragraph{Governance, law, and meta-level design.}
The same categorical machinery can be used to model governance operators that act on institutions, together with constraints that restrict how institutions may change. This suggests: (i) designing concrete governance observables (drift, disagreement, confidence-inflation indicators) that can be computed from FRFP-compatible logs; (ii) exploring families of governance operators that incorporate human oversight at the institutional level without collapsing into fully automated correctness claims; and (iii) mapping these patterns onto emerging regulatory regimes for AI in science, medicine, and finance, to clarify which regulatory templates are structurally compatible with an explicit/tacit separation.

\paragraph{Richer explicit logics and tooling.}
On the explicit side, the framework is agnostic about the internal logic carried by \(E\). A further direction is to instantiate \(E\) with richer structures: type-theoretic or temporal logics, proof objects from interactive theorem provers, or domain-specific calculi for simulation and data analysis. This would support more expressive explicit pipelines while preserving the \(RB\)--\(TE\)--\(HFD\) suffix and its correctness properties. In parallel, one can build FRFP-aware tooling: typed interface languages, normalization and auditing layers for \(N \ltimes E\), and assistants that expose their operations as TDG terms rather than as opaque conversations.

\paragraph{Benchmarking and comparative evaluation.}
Finally, FRFP suggests protocol-level evaluation criteria. Instead of scoring isolated model outputs, one can compare whole workflows---different ways of structuring the explicit/tacit boundary, navigation, repair, and institutional aggregation---using metrics derived from the dynamic theory, such as time-to-hallucination distributions, robustness under degradation, and sensitivity to feedback coupling. Developing such benchmarks would allow FRFP-compatible systems to be compared along principled axes, and would help connect the abstract theory to concrete engineering trade-offs.

FRFP is intended as a foundation, not a finished product. Its value lies in making constraints and invariants explicit so that future systems---both technical and institutional---can be designed, analysed, and compared within a common formal framework.


\pagebreak
\appendix

This appendix develops the categorical and algebraic core of the
FRFP (Foundational Reasoning and Feedback Protocol) architecture.
It provides a self-contained account of:

\begin{itemize}
  \item an ambient category with explicit and tacit modes;
  \item a free explicit algebra of task decompositions (the Task Decomposition Grammar, TDG);
  \item a navigation groupoid acting on explicit shapes;
  \item a fibration of explicit pipelines over shapes;
  \item a Grothendieck-style semidirect product \(N \ltimes E\);
  \item a semantic functor into a preorder-enriched tacit category; and
  \item correctness and canonicality theorems for Epistemic Structure.
\end{itemize}

We assume familiarity with basic category theory, rewriting systems,
and context-free grammars.

Throughout, we write \(\Cfull\) for the ambient category,
\(\Espace\) for explicit space, and \(\Tspace\) for tacit space.

%======================================================================
\subsection{Foundations}
%======================================================================

\subsubsection{Ambient Category \texorpdfstring{\(\Cfull\)}{Cfull}}

\begin{definition}[Ambient category]
The ambient category \(C_{\mathrm{full}}\) is an arbitrary category that may contain objects beyond those used by the FRFP protocol. We designate three \emph{distinguished} objects
\(\varnothing, E, T \in \mathrm{Ob}(C_{\mathrm{full}})\):
\begin{itemize}
  \item \(\varnothing\): the empty explicit state;
  \item \(E\): the explicit state object (machine-manipulable representations);
  \item \(T\): the tacit state object (human correctness-bearing states).
\end{itemize}
Morphisms in \(\Cfull\) are partitioned into three epistemic classes:
\begin{itemize}
  \item explicit morphisms with codomain \(E\), forming the explicit subcategory
        \(\Espace \subseteq \Cfull\);
  \item tacit morphisms with codomain \(T\), forming the tacit subcategory
        \(\Tspace \subseteq \Cfull\);
  \item boundary morphisms \(E \to T\), of which a distinguished one is the
        \emph{Representational Backflow} operator
        \(\RB : E \to T\).
\end{itemize}
Composition in \(\Cfull\) satisfies the usual category axioms:
\[
(h \circ g) \circ f \;=\; h \circ (g \circ f), \qquad
f \circ \id_X \;=\; f \;=\; \id_Y \circ f.
\]
\end{definition}

\begin{definition}[Explicit--Tacit Separation (ETS)]
The \emph{Explicit--Tacit Separation} axiom requires:
\begin{itemize}
  \item there are no morphisms \(T \to E\) in \(\Cfull\);
  \item tacit evaluation cannot be embedded inside explicit pipelines;
  \item the boundary morphism \(\RB : E \to T\) is not invertible.
\end{itemize}
\end{definition}

\subsubsection{Modal Interpretation of Explicit and Tacit Knowledge}

\begin{definition}[Explicit modality]
A cognitive item \(k\) is represented in \emph{explicit} modality if there exist
maps
\[
\enc : E \to \Sigma^{\ast}, \qquad
\dec : \Sigma^{\ast} \to E
\]
such that \(\dec(\enc(k))\) is equivalent to \(k\) in its epistemic role
(e.g., as a proof, equation, or program). In particular, \(k\) can be
faithfully encoded as a finite symbolic artifact and recovered without
loss of its correctness-relevant structure.
\end{definition}

\begin{definition}[Tacit modality]
A cognitive item is represented in \emph{tacit} modality when no such
encoding preserves its correctness-bearing content. Judgments, intuitions,
acceptance decisions, and other human correctness-bearing acts live in
tacit space.
\end{definition}

Human expression from tacit cognition to explicit content (e.g., writing a
new question, revising a claim) is modeled externally as the appearance of
new explicit morphisms in \(\Espace\); it is not modeled as a morphism
\(T \to E\).

\subsubsection{Explicit Category \texorpdfstring{\(\Espace\)}{E}}

\begin{definition}[Explicit category]
The explicit category \(\Espace\) is the full subcategory of \(\Cfull\)
on objects with codomain \(E\). Its core generating morphisms are:
\[
\RI : \varnothing \to E
\quad\text{(Representation Initiation)},\qquad
\EC : E \to E
\quad\text{(Explicit Computation)},\qquad
\ED : E \to E
\quad\text{(Explicit Diagnostics)}.
\]
\end{definition}

\noindent
\textbf{Remark.}
\(\EC\) represents algorithmic transformations of explicit content
(e.g., inference, rewriting, search). \(\ED\) represents explicit
checks and diagnostics (e.g., consistency checks, sanity tests),
preparing but not performing tacit evaluation.

\subsubsection{Tacit Category \texorpdfstring{\(\Tspace\)}{T}}

We now formalize tacit space as a preorder-enriched one-object category.

\begin{definition}[Tacit category]
The tacit category \(\Tspace\) has:
\begin{itemize}
  \item a single object, denoted again by \(T\);
  \item hom-set \(\Hom_{\Tspace}(T,T)\) equipped with a composition
        operation \(\circ\), making it a monoid;
  \item a preorder \(\sqsubseteq\) on \(\Hom_{\Tspace}(T,T)\) such that
        composition is monotone:
        \[
        f_1 \sqsubseteq f_2,\ g_1 \sqsubseteq g_2
        \quad\Rightarrow\quad
        g_1 \circ f_1 \sqsubseteq g_2 \circ f_2.
        \]
\end{itemize}
\end{definition}

\begin{definition}[Tacit primitives]
The tacit primitives are the morphisms
\[
\TE : T \to T \quad\text{(Tacit Evaluation)},\qquad
\HFD : T \to T \quad\text{(Human Final Decision)},
\]
subject to:
\begin{itemize}
  \item \emph{idempotence}:
        \(\TE \circ \TE = \TE\) and \(\HFD \circ \HFD = \HFD\);
  \item \emph{commutation}:
        \(\TE \circ \HFD = \HFD \circ \TE\);
  \item \emph{monotonicity}: for every \(f : T \to T\),
        \[
        f \sqsubseteq \TE \circ f,\qquad
        f \sqsubseteq \HFD \circ f.
        \]
\end{itemize}
\end{definition}

\begin{definition}[Correctness quotient]
Let \(J\) be a finite set of correctness judgments
(e.g., \(\{\mathsf{accept}, \mathsf{reject}, \mathsf{revise}\}\)).
A \emph{correctness quotient} is a surjective map
\[
\gamma : \Obj(\Tspace) \to J
\]
assigning to each tacit state a correctness class, subject to
\[
\gamma \circ \TE = \gamma = \gamma \circ \HFD.
\]
Thus tacit evaluation and final decision may refine or stabilize a tacit
state but do not change its correctness class.
\end{definition}

\paragraph{Interpretation and non-triviality of \texorpdfstring{$T$}{T}.}
The tacit category T is not introduced merely as a formal firewall that “reserves correctness to tacit space”.
Rather, it is intended to capture a substantive empirical claim about long-horizon scientific and regulatory work:
in these regimes, there are correctness-bearing aspects of judgment that cannot be exhaustively encoded in a fixed,
finite explicit representation without loss.

Definitions A.3–A.4 distinguish explicit and tacit modalities by the existence of lossless encodings into symbolic artefacts.
In many familiar settings—for example, textbook exercises, small programs, or contest-style benchmarks—it is reasonable
to treat \emph{inner} task correctness as effectively explicit: the task is specified by a finite prompt, there is a fixed
input–output relation or reference implementation, and success can be checked by a test harness or proof checker.
FRFP does not deny the existence or value of such regimes. Instead, it treats them as \emph{tacit-insensitive inner layers}:
conditional on a surrounding tacit backdrop (choice of formalism, modelling assumptions, risk tolerances),
correctness for that inner task behaves as if it were a pure explicit oracle on E.

The global explicit/tacit split remains.
Even in proof-assistant or test-oracle settings, tacit judgment is involved in selecting the calculus, specifying what the
theorem or contract should state, and deciding whether conformance to that specification is adequate for the real-world
purpose.
What FRFP targets are workflows where this outer tacit layer is non-trivial and dynamically involved in long-horizon
use: tasks evolve, standards and constraints shift, and background scientific understanding and norms play an essential
role in deciding whether an explicit artefact is “good enough”.
In those settings it is non-trivial, and in our view appropriate, to treat the final correctness status of an explicit
pipeline as living in T rather than in E. (See Appendix~D.4 for further discussion.)


A reasonable skeptic might ask: under what conditions could one reject the Epistemic Spaces axioms
and the later Epistemic Algebra structure on $T$? There are at least three salient counter-scenarios:
\begin{enumerate}
  \item \textbf{Tacit-insensitive inner correctness.}
For some domains and layers, once a formalism, specification, and world-modelling stance have been tacitly fixed,
every correctness-bearing judgment at the inner level can be captured by a finite explicit specification and a sound
and complete proof- or test-based oracle.
At that inner layer one can model correctness directly in E and treat T as abstracted away; FRFP’s Phase–3
inevitability and non-automation results then become vacuous or uninteresting for that particular subproblem.
Our focus is on regimes where the outer tacit layer—where specifications, standards, and modelling choices are
themselves evolving—remains dynamically involved in correctness.

  \item \textbf{No degradation or escalation.} Epistemic Algebra impossibility results rely on a
  degradation operator $\delta$ and recurring nontrivial requirements relative to a floor $r$.
  If one believes that, in their target setting, tacit state does not degrade in any relevant
  sense, or that requirements do not escalate or recur in the way we assume (for example,
  because tasks are short-lived and fully specified at the outset), then the inevitability
  theorems may not apply. Such systems fall outside FRFP's intended long-horizon regime.
  \item \textbf{Perfect tacit emulation.} If one posits an AI system that can internally
  reconstruct or simulate human tacit evaluation from explicit traces without loss---violating
  the NTER and CSC constraints---then our explicit--tacit separation (ETS, HEG, HEC) is simply
  rejected. FRFP explicitly rules out this possibility by axiom rather than attempting to model
  it.
\end{enumerate}

In summary, treating correctness as localized in $T$ is a substantive modelling choice:
it reflects the view that, for many real scientific and governance workflows, there is residual
human judgment that cannot be fully captured by static benchmarks or formal specifications, and
that this residual plays an essential role in long-horizon correctness. The impossibility and
inevitability results of Epistemic Structure and Dynamics should therefore be read as conditional on this non-trivial
tacit structure: they apply precisely in domains where correctness genuinely outstrips what can
be made explicit and mechanically checked.


\subsubsection{Boundary Morphism and Semantic RB Functor}

\begin{definition}[Boundary morphism \(\RB\)]
The \emph{Representational Backflow} operator is a morphism
\[
\RB : E \to T
\]
in \(\Cfull\). It is the unique explicit-to-tacit interface specified by
the Representational Backflow axiom.
\end{definition}

Let \(E_{\mathrm{sem}}\) denote the free explicit category generated by the
Task Decomposition Grammar (defined in Section~\ref{sec:appendix-tdg}).

\begin{definition}[Semantic RB functor]
The \emph{semantic RB} is a functor
\[
\RBsharp : E_{\mathrm{sem}} \to \Tspace
\]
assigning to each explicit pipeline a tacit state representing the content
received by the human for evaluation.
\end{definition}

\subsubsection{Interaction Protocols and Epistemic Spaces Axioms}

\begin{definition}[Interaction Protocols constraints]
The operational environment satisfies:
\begin{itemize}
  \item \textbf{IL (Interaction Limitation):} interaction is bounded and turn-based;
  \item \textbf{AR (Alignment Restriction):} the system must not autonomously act
        outside explicit instructions;
  \item \textbf{MD (Mode Discipline):} separation between explicit and tacit modes
        is enforced at the interface level;
  \item \textbf{NTER (No Tacit Evaluation Reconstruction):} the system must not
        reconstruct tacit evaluation internally from explicit traces;
  \item \textbf{CSC (Correctness Separation Constraint):} correctness decisions
        must not be delegated to explicit computation.
\end{itemize}
\end{definition}

\begin{definition}[Epistemic Spaces axioms]
An admissible epistemic framework satisfies:
\begin{itemize}
  \item \textbf{ETS (Explicit--Tacit Separation):} no morphisms \(T \to E\);
  \item \textbf{HEG (Human-Exclusive Grounding):} correctness grounding occurs
        only in tacit space via \(\TE\), executed by humans;
  \item \textbf{HEC (Human-Exclusive Closure):} final acceptance/rejection occurs
        only via \(\HFD\), executed by humans;
  \item \textbf{RB (Representational Backflow):} there is a unique boundary
        morphism \(\RB : E \to T\) used by all pipelines;
  \item \textbf{CP (Compositional Pipelines):} epistemic processes are built as
        compositions of primitives into pipelines;
  \item \textbf{AE (AI-Explicit Restriction):} AI systems act only via
        explicit-space morphisms in \(\Espace\).
\end{itemize}
\end{definition}

%======================================================================
\subsection{Explicit Algebra and the Task Decomposition Grammar}
\label{sec:appendix-tdg}
%======================================================================

\subsubsection{TDG Signature and Free Category}

\begin{definition}[TDG signature]
The Task Decomposition Grammar (TDG) signature \(\Sigma_{\mathrm{TDG}}\)
consists of:
\begin{itemize}
  \item objects: \(\varnothing\) and \(E\);
  \item generating morphisms:
  \[
    \RI : \varnothing \to E,\qquad
    \EC : E \to E,\qquad
    \ED : E \to E;
  \]
  \item structural operators \(\BRANCH\), \(\PARALLEL\), \(\LOOP\) acting
        as higher-arity constructors on morphisms.
\end{itemize}
\end{definition}

\begin{definition}[Free explicit category]
The explicit semantic category is the free category on the TDG signature:
\[
E_{\mathrm{sem}} := \FreeCat(\Sigma_{\mathrm{TDG}}).
\]
There is a canonical interpretation functor
\(\Ufun : E_{\mathrm{sem}} \to \Espace\)
sending each TDG term to its denotation as an explicit morphism.
\end{definition}

\begin{theorem}[Free explicit algebra]
\label{thm:free-explicit-algebra}
For any category \(\mathcal{D}\) and any interpretation of the TDG signature
\(\Sigma_{\mathrm{TDG}}\) in \(\mathcal{D}\), there exists a unique functor
\(\tilde{F} : E_{\mathrm{sem}} \to \mathcal{D}\) extending that interpretation.
\end{theorem}

\begin{proof}
This is the standard universal property of the free category on a graph:
objects of \(E_{\mathrm{sem}}\) are objects of \(\Sigma_{\mathrm{TDG}}\), and morphisms
are freely generated paths in the underlying graph. Any graph morphism
from \(\Sigma_{\mathrm{TDG}}\) into \(\mathcal{D}\) extends uniquely to a functor.
\end{proof}

\subsubsection{TDG Grammar}

We present TDG both as a signature and as a context-free grammar
over terms.

\paragraph{Nonterminals.}
We use:
\[
Q \quad\text{for explicit pipelines}, \qquad
P \quad\text{for full pipelines (explicit + tacit suffix)}.
\]

\paragraph{Terminals.}
The terminal symbols are:
\[
\RI,\, \EC,\, \ED,\, \RB,\, \TE,\, \HFD,\,
\BRANCH,\, \PARALLEL,\, \LOOP.
\]

\paragraph{Productions.}
The grammar productions are:
\begin{align*}
Q &::= \RI,\\
Q &::= Q;\EC,\\
Q &::= Q;\ED,\\
Q &::= \BRANCH(Q_1,\dots,Q_k),\\
Q &::= \PARALLEL(Q_1,\dots,Q_k),\\
Q &::= \LOOP(Q_{\mathrm{cond}}, Q_{\mathrm{body}}),\\
P &::= Q;\RB;\TE;\HFD.
\end{align*}
Thus TDG presents \(E_{\mathrm{sem}}\) as an algebra of explicit pipelines
with a fixed tacit suffix.

\subsubsection{Shapes and Shape Category}

\begin{definition}[TDG shape]
TDG shapes are generated by the grammar
\[
S ::= \chain \mid
       \branch(S_1,\dots,S_k) \mid
       \parallel(S_1,\dots,S_k) \mid
       \loopshape(S_{\mathrm{cond}}, S_{\mathrm{body}}).
\]
\end{definition}

\begin{definition}[Shape of a TDG term]
The mapping
\(\shape : \text{TDG-terms} \to \text{Shapes}\) is defined inductively by:
\begin{itemize}
  \item \(\shape(\RI) = \chain;\)
  \item \(\shape(Q;\EC) = \shape(Q)\) and similarly \(\shape(Q;\ED) = \shape(Q)\);
  \item \(\shape(\BRANCH(Q_1,\dots,Q_k))
         = \branch(\shape(Q_1),\dots,\shape(Q_k));\)
  \item analogously for \(\PARALLEL\) and \(\LOOP\).
\end{itemize}
\end{definition}

\begin{definition}[Shape category \(\Scat\)]
The shape category \(\Scat\) has:
\begin{itemize}
  \item objects: shapes \(S\) as above;
  \item morphisms: structural rewiring maps between shapes
        (e.g., regrouping, insertion/removal of trivial chains),
        closed under composition and identities.
\end{itemize}
\end{definition}

\subsubsection{Explicit Reduction System}

\begin{definition}[Explicit reduction]
A binary relation \(\to_E\) on \(\Espace\) is an \emph{explicit reduction system}
if it is a typed rewrite system satisfying:
\begin{itemize}
  \item \textbf{typing:} if \(e : X \to E\) and \(e \to_E e'\), then \(e' : X \to E\);
  \item \textbf{TDG compatibility:} reductions act componentwise on TDG
        combinators (\(\BRANCH\), \(\PARALLEL\), \(\LOOP\));
  \item \textbf{termination:} there is no infinite chain
        \(e_0 \to_E e_1 \to_E e_2 \to_E \dots\);
  \item \textbf{local confluence:} if \(e \to_E e_1\) and \(e \to_E e_2\), then
        \(e_1\) and \(e_2\) have a common descendant.
\end{itemize}
We treat these properties as axioms of the explicit reduction system.
\end{definition}

Under these assumptions, every explicit artifact admits a unique normal
form modulo \(\to_E\).

%======================================================================
\subsection{Navigation and Shapes}
%======================================================================

\subsubsection{Navigation Category and Groupoid}

\begin{definition}[Navigation category \(N_0\)]
The navigation category \(N_0\) has:
\begin{itemize}
  \item objects: shapes \(S\);
  \item generating morphisms: navigation steps
        \(g : S \to S'\), such as \(\TRYMODEL\),
        \(\ADDASSUMPTION\), \(\DROPASSUMPTION\),
        \(\MERGE\), \(\RESTART\), each with a well-defined
        source and target shape;
  \item morphisms: finite composites of such generators.
\end{itemize}
\end{definition}

\begin{definition}[Navigation groupoid \(N\)]
The navigation groupoid \(N\) is obtained by freely adjoining inverses
to morphisms of \(N_0\) and quotienting by the relations
\[
g^{-1} \circ g = \id_S, \qquad
g \circ g^{-1} = \id_{S'},
\]
for each generator \(g : S \to S'\).
Thus every navigation step is reversible up to equality in \(N\).
\end{definition}

\subsubsection{Fibration \(p : E \to S\)}

\begin{definition}[Projection to shapes]
The functor \(p : \Espace \to \Scat\) sends each explicit pipeline
\(e \in \Espace\) to its shape
\(p(e) := \shape(e)\), and each explicit morphism to the corresponding
structural morphism between shapes.
\end{definition}

\begin{definition}[Fibers]
For each shape \(S\), the fiber
\[
E_S := p^{-1}(S)
\]
is the category of explicit morphisms with shape \(S\).
\end{definition}

\begin{definition}[Cartesian liftings and reindexing]
For each navigation morphism \(n : S \to S'\) in \(N\), we assume a
cartesian lifting
\[
n^\ast : E_S \to E_{S'}
\]
which reindexes explicit pipelines along \(n\). This yields an indexed
category
\[
E : N^{\mathrm{op}} \to \mathbf{Cat},
\qquad
S \mapsto E_S,\quad
n \mapsto n^\ast.
\]
\end{definition}

%======================================================================
\subsection{Combined System: Grothendieck Construction}
%======================================================================

\subsubsection{Semidirect Product \texorpdfstring{\(N \ltimes E\)}{N ⋉ E}}

\begin{definition}[Grothendieck construction \(N \ltimes E\)]
Given the indexed category
\(E : N^{\mathrm{op}} \to \mathbf{Cat}\), the Grothendieck
construction \(\int_N E\) (also written \(N \ltimes E\)) has:
\begin{itemize}
  \item objects: pairs \((S, e)\) with \(S \in \Obj(N)\) and \(e \in E_S\);
  \item morphisms \((S, e) \to (S', e')\): pairs \((n, f)\) where
        \(n : S \to S'\) in \(N\) and
        \(f : n^\ast(e) \to e'\) in \(E_{S'}\).
\end{itemize}
Composition is given by
\[
(n_1, f_1) \circ (n_2, f_2)
\;=\;
\bigl(
  n_1 \circ n_2,\,
  f_1 \circ n_1^\ast(f_2)
\bigr).
\]
Informally, we may write elements as \((n, e)\) and composition as
\((n_1, e_1) \circ (n_2, e_2) =
 (n_1 \circ n_2,\, e_1 \circ n_1(e_2))\) whenever the lifted action is defined.
\end{definition}

\subsubsection{Rewrite System on \texorpdfstring{\(N \ltimes E\)}{N ⋉ E}}

\begin{definition}[Combined reduction]
The one-step reduction relation \(\to\) on \(N \ltimes E\) is generated by:
\begin{itemize}
  \item \textbf{explicit step:}
        if \(e \to_E e'\) then \((n, e) \to (n, e')\);
  \item \textbf{navigation step:}
        if \(n \to_N n'\) and \(n'\) is defined on \(e\), then
        \((n, e) \to (n', n'(e))\).
\end{itemize}
We assume:
\begin{enumerate}[label=(A\arabic*)]
  \item \(\to_E\) is terminating and locally confluent;
  \item navigation reductions \(\to_N\) have normal forms in each component
        of \(N\);
  \item the lifted action is functorial and preserves \(\RB\), \(\TE\),
        and \(\HFD\).
\end{enumerate}
These are axioms of the FRFP reduction environment.
\end{definition}

%======================================================================
\subsection{Semantics and Tacit Correctness}
%======================================================================

\subsubsection{Semantic Functor}

\begin{definition}[Semantic bracket]
A \emph{semantic functor} is a functor
\[
\llbracket{-}\rrbracket : N \ltimes E_{\mathrm{sem}} \to \Tspace
\]
assigning to each combined configuration a tacit state.
On explicit pipelines \(e \in E_{\mathrm{sem}}\), we set
\[
\llbracket e \rrbracket := \HFD \circ \TE \circ \RBsharp(e).
\]
\end{definition}

\begin{definition}[Navigation invariance]
The functor \(\llbracket{-}\rrbracket\) is \emph{navigation-invariant up to
preorder} if
\[
\llbracket (n, e) \rrbracket \;\sqsubseteq\; \llbracket e \rrbracket
\]
for all \(n,e\) where the action is defined, and equality holds in the
quotient by the induced equivalence on tacit states in confluent cases.
\end{definition}

\begin{definition}[Observable correctness]
Given a pipeline \(P\), its observable correctness outcome is
\[
\gamma(\llbracket P \rrbracket) \in J.
\]
\end{definition}

%======================================================================
\subsection{Main Theorems}
%======================================================================

\subsubsection{Well-Formedness of TDG Pipelines}

\begin{theorem}[Well-formedness]
\label{thm:well-formed}
Every TDG-generated pipeline \(P\) is a well-typed morphism in \(\Cfull\) of
the form
\[
\RI \;\to\; (\EC/\ED)^{\ast} \;\to\; \RB \;\to\; \TE \;\to\; \HFD.
\]
\end{theorem}

\begin{proof}
By induction on TDG derivations. The base case \(\RI\) initializes an object
in \(E\). The constructors \(\EC\) and \(\ED\) have codomain \(E\) and thus
preserve explicit type. Structural combinators (\(\BRANCH\), \(\PARALLEL\),
\(\LOOP\)) act within the explicit category and preserve codomain \(E\)
componentwise. The production for \(P\) enforces the suffix
\(\RB; \TE; \HFD\), whose types are \(E \to T \to T \to T\).
By ETS there are no morphisms \(T \to E\), so:
\begin{itemize}
  \item no tacit operator can appear before \(\RB\);
  \item no explicit operator can appear after \(\RB\).
\end{itemize}
Thus every derivation yields a well-typed morphism of the stated form.
\end{proof}

\subsubsection{Explicit Completeness of TDG}

\begin{theorem}[Explicit completeness]
\label{thm:explicit-completeness}
TDG generates exactly the admissible explicit morphisms in \(\Espace\).
\end{theorem}

\begin{proof}
By construction, \(E_{\mathrm{sem}}\) is the free category on the TDG
signature, and \(\Ufun\) interprets TDG terms as explicit morphisms in
\(\Espace\). Every explicit morphism in the FRFP setting is built from
\(\RI\), \(\EC\), \(\ED\), and the structural operators, hence arises from
a TDG term; conversely, each TDG term denotes a well-typed explicit
morphism. The image of TDG under \(\Ufun\) is thus exactly the admissible
explicit morphisms.
\end{proof}

\subsubsection{Minimality of Primitive Set}

\begin{theorem}[Minimality]
\label{thm:minimality}
The primitive set
\(\{\RI,\EC,\ED,\RB,\TE,\HFD\}\)
is minimal under the Epistemic Spaces axioms: removing any primitive violates at
least one of ETS, HEG, HEC, RB, CP, or AE.
\end{theorem}

\begin{proof}
Briefly:
\begin{itemize}
  \item \(\RI\) is required to start from \(\varnothing\) and generate explicit
        content; without it there is no canonical initiation of pipelines.
  \item \(\EC\) is required for generative explicit reasoning; without it,
        explicit computation collapses.
  \item \(\ED\) is required for explicit diagnostics without tacit evaluation;
        removing it forces diagnostics either into tacit space (violating AE)
        or to be indistinguishable from \(\EC\) (violating CP and the intended
        modal roles).
  \item \(\RB\) is required for explicit-to-tacit handoff; without it there is
        no admissible morphism \(E \to T\), contradicting the requirement that
        pipelines reach tacit evaluation.
  \item \(\TE\) is required for tacit grounding (HEG); without it correctness
        cannot be evaluated in tacit space.
  \item \(\HFD\) is required for final correctness closure (HEC); without it
        there is no morphism corresponding to final acceptance/rejection.
\end{itemize}
Removing any primitive causes at least one Epistemic Spaces axiom to fail, so the
set is minimal.
\end{proof}

\subsubsection{Epistemic Spaces Initiality}

\begin{definition}[Epistemic Spaces framework]
A \emph{Epistemic Spaces framework} is a tuple
\[
\mathcal{F}_1
 \;=\;
(\mathcal{C}, E_{\mathcal{C}}, T_{\mathcal{C}},
 \RI_{\mathcal{C}}, \EC_{\mathcal{C}}, \ED_{\mathcal{C}},
 \RB_{\mathcal{C}}, \TE_{\mathcal{C}}, \HFD_{\mathcal{C}})
\]
satisfying the Epistemic Spaces axioms.
Morphisms in the category \(\Frm_1\) of Epistemic Spaces frameworks are functors
preserving \(E\), \(T\), and all six primitives.
\end{definition}

\begin{theorem}[Epistemic Spaces initiality]
\label{thm:phase1-initial}
There exists an initial object
\(F^{(1)}_{\mathrm{FRFP}} \in \Frm_1\).
\end{theorem}

\begin{proof}
Construct \(F^{(1)}_{\mathrm{FRFP}}\) as the free category on objects
\(E, T\) and morphisms \(\RI,\EC,\ED,\RB,\TE,\HFD\), modulo the equations
encoding the Epistemic Spaces axioms (ETS, HEG, HEC, RB, CP, AE). Any other
Epistemic Spaces framework interprets these generators in a way that satisfies the
same equations, yielding a unique induced functor
\(F^{(1)}_{\mathrm{FRFP}} \to \mathcal{F}_1\). Thus
\(F^{(1)}_{\mathrm{FRFP}}\) is initial in \(\Frm_1\).
\end{proof}
\paragraph{Remark A.34a (Non-trivial content of Epistemic Spaces initiality).}
The statement that $F^{(1)}_{\mathrm{FRFP}}$ is initial in $\mathrm{Frm}_1$ uses the free-category
construction, but its content is more than formal packaging. The Epistemic Spaces axioms (ETS, HEG,
HEC, RB, CP, AE), together with the minimality result (Theorem~A.32), severely constrain which
primitive roles may appear and how they can be composed. Initiality then says that any epistemic
framework respecting exactly these roles and axioms must interpret RI, EC, ED, RB, TE, HFD in a
way that factors uniquely through the FRFP ontology. In particular, attempts to ``hide'' additional
power in explicit space (for example, by letting EC or ED implicitly perform tacit evaluation) break
the axioms rather than yielding a genuinely different Epistemic Spaces framework.

\subsubsection{Epistemic Algebra Canonicality}

\begin{definition}[Epistemic Algebra architecture]
A \emph{Epistemic Algebra architecture} is a tuple
\[
\mathcal{F}_2
 =
(E,\ S,\ p,\ N,\ (\cdot)^\ast,\ N \ltimes E)
\]
satisfying:
\begin{itemize}
  \item TDG presents \(E\) as a free explicit category;
  \item \(N\) is a navigation groupoid on shapes \(S\);
  \item \(p : E \to S\) is a fibration with cartesian liftings \(n^\ast\);
  \item \(N \ltimes E\) is the Grothendieck construction of the indexed
        category \(E : N^{\mathrm{op}} \to \mathbf{Cat}\);
  \item the boundary structure \(\RB\)--\(\TE\)--\(\HFD\) is preserved
        under navigation.
\end{itemize}
Morphisms in the category of Epistemic Algebra architectures preserve all of this
structure.
\end{definition}

\begin{theorem}[Epistemic Algebra canonicality]
\label{thm:phase2-canonicality}
The FRFP Epistemic Algebra architecture
\(F^{(2)}_{\mathrm{FRFP}}\) is a canonical object in the category of
Epistemic Algebra architectures: for every Epistemic Algebra architecture \(\mathcal{A}\),
there exists a unique structure-preserving morphism
\[
F^{(2)}_{\mathrm{FRFP}} \to \mathcal{A}.
\]
\end{theorem}

\begin{proof}
By the universal properties of:
\begin{itemize}
  \item the free explicit category \(E_{\mathrm{sem}}\);
  \item the groupoid completion constructing \(N\);
  \item the Grothendieck construction forming \(N \ltimes E\).
\end{itemize}
Any Epistemic Algebra architecture satisfying the same axioms must interpret the
FRFP construction in a role-preserving way, inducing a unique morphism
from \(F^{(2)}_{\mathrm{FRFP}}\).
\end{proof}

\subsubsection{Confluence of \texorpdfstring{\(N \ltimes E\)}{N ⋉ E}}

\begin{theorem}[Confluence of \(N \ltimes E\)]
\label{thm:n-e-confluence}
Under assumptions \textup{(A1)}--\textup{(A3)}, the reduction relation
\(\to\) on \(N \ltimes E\) is confluent.
\end{theorem}

\begin{proof}[Proof sketch]
Local peaks \((n, e) \leftarrow (n_0, e_0) \to (n', e')\) decompose into:
\begin{itemize}
  \item explicit--explicit peaks, which join by local confluence of
        \(\to_E\);
  \item navigation--navigation peaks, which join by normalization in each
        connected component of \(N\);
  \item mixed peaks, which join because the lifted action commutes with
        explicit reduction and preserves \(\RB,\TE,\HFD\).
\end{itemize}
Together with termination up to navigation-normalization, this yields
global confluence.
\end{proof}
\paragraph{Remark A.37a (Non-trivial aspects of Epistemic Algebra structure).}
While Theorems~A.36 and A.37 are phrased in terms of standard categorical constructions
(free categories, groupoid completions, Grothendieck constructions), several aspects of the
Epistemic Algebra architecture are mathematically substantive:

\begin{itemize}
  \item \textbf{Compatibility with RB--TE--HFD.} The lifted navigation action and the
  Grothendieck construction are required to preserve the boundary structure. This rules out
  navigation steps that implicitly smuggle tacit evaluation back into explicit space or that
  duplicate/delete the RB--TE--HFD suffix. Showing that the action can be defined so as to
  commute with explicit reduction while preserving RB, TE, HFD is a genuine structural
  constraint on admissible navigation schemes.
  \item \textbf{Confluence of the combined reduction.} Confluence of $\to$ on $N \ltimes E$ is not
  automatic from confluence of $\to_E$ and normalization in $N$. The proof must analyze
  mixed peaks and use the functoriality of the lifted action, together with preservation of
  RB, TE, and HFD, to ensure that explicit and navigation steps can be permuted without
  changing the normal form.
  \item \textbf{Interaction with semantics.} The semantic correctness theorem (Theorem~A.38)
  depends crucially on the fact that navigation-invariant semantics on $N \ltimes E$ respects
  the Interaction Protocols constraints (no illicit $T \to E$ morphisms) and the boundary-preservation
  properties above. This is what forces correctness judgments to depend only on the tacit
  evaluation of the explicit pipeline, not on the particular exploration trajectory.
\end{itemize}

These non-trivial interactions justify treating Epistemic Algebra canonicality as more than a choice of
notation: any alternative architecture satisfying the same constraints must reproduce this
structure and, in particular, cannot avoid the resulting preservation and non-automation results
without breaking one of the axioms.

\subsubsection{Semantic Correctness for Epistemic Structure}

Collecting the constructions above, we obtain a semantic correctness
theorem for Epistemic Structure FRFP pipelines.

\begin{theorem}[Semantic correctness]
\label{thm:semantic-correctness}
Let \(P\) be any TDG-generated pipeline that is Interaction Protocols admissible and
Epistemic Spaces/Algebra well-typed. Then:
\begin{enumerate}
  \item \(P\) has a unique normal form in \(N \ltimes E\) modulo navigation;
  \item the observable correctness outcome \(\gamma(\llbracket P \rrbracket)\)
        is invariant under:
        \begin{itemize}
          \item explicit reduction \(\to_E\),
          \item navigation in \(N\),
          \item composition in the semidirect product \(N \ltimes E\).
        \end{itemize}
\end{enumerate}
In particular, correctness judgments depend only on the (tacitly evaluated)
semantic content of the explicit pipeline, not on the particular exploration
trajectory used to construct it.
\end{theorem}

\begin{proof}[Proof sketch]
Item~(1) follows from Theorem~\ref{thm:n-e-confluence} and the termination
assumptions (A1)--(A3). Item~(2) follows from navigation invariance of the
semantic functor, the Interaction Protocols constraints (which rule out illicit T-to-E
feedback or hidden tacit reconstruction), and the fact that \(\RB,\TE,\HFD\)
are preserved under navigation and do not appear in the explicit phase.
\end{proof}

\subsection{Worked example: sepsis consult as an FRFP pipeline}
\label{sec:cp-sepsis-example}
We sketch a concrete instance of the FRFP calculus corresponding to the sepsis consult scenario
of \S\ref{sec:case-sepsis}. The goal is to show how a real conversational workflow is
represented as a TDG term, reduced to normal form in $N \ltimes E$, and checked for
Interaction Protocols admissibility and Epistemic Spaces/Algebra well-typedness.

\paragraph{TDG and free explicit category.}
Let the TDG have objects
\[
  \Ob(E_0) = \{\textsf{Ctx}, \textsf{Input}, \textsf{CaseSummary},
  \textsf{DraftPlan}, \textsf{CandidatePlan}, \textsf{Report}, \textsf{Decision}\}
\]
and generators
\begin{align*}
  \mathsf{RB} &: \textsf{Ctx} \to \textsf{Input}, \\
  \mathsf{prep} &: \textsf{Input} \to \textsf{CaseSummary}, \\
  \mathsf{ehr} &: \textsf{CaseSummary} \to \textsf{CaseSummary}, \\
  \mathsf{calc} &: \textsf{CaseSummary} \to \textsf{CaseSummary}, \\
  \mathsf{assist}_0 &: \textsf{CaseSummary} \to \textsf{DraftPlan}, \\
  \mathsf{review}_\mathrm{res} &: \textsf{DraftPlan} \to \textsf{CandidatePlan}, \\
  \mathsf{review}_\mathrm{att} &: \textsf{CandidatePlan} \to \textsf{Report}, \\
  \mathsf{HFD} &: \textsf{Report} \to \textsf{Decision}.
\end{align*}
Let $E$ be the free category on this graph, modulo the Epistemic Algebra equations for $\mathsf{TE}$ and
the RB--TE--HFD boundary. We define
\[
  \mathsf{assist} := \mathsf{assist}_0 \circ \mathsf{calc} \circ \mathsf{ehr}
  : \textsf{CaseSummary} \to \textsf{DraftPlan},
\]
and the explicit sepsis pipeline
\[
  e_{\mathrm{sepsis}} :=
  \mathsf{HFD} \circ \mathsf{review}_\mathrm{att} \circ \mathsf{review}_\mathrm{res} \circ
  \mathsf{assist}_0 \circ \mathsf{calc} \circ \mathsf{ehr} \circ \mathsf{prep} \circ \mathsf{RB}
  : \textsf{Ctx} \to \textsf{Decision}.
\]

\paragraph{Reduction and normal form.}
Let $R_E$ be the explicit reduction system on $E$ that contracts $\mathsf{ehr}$ and $\mathsf{calc}$
into $\mathsf{assist}$ inside the $\mathsf{assist}_0$ wrapper:
\[
  \mathsf{assist}_0 \circ \mathsf{calc} \circ \mathsf{ehr} \;\to\; \mathsf{assist}.
\]
By confluence of $R_E$ and compatibility with the RB--TE--HFD boundary, the normal form of
$e_{\mathrm{sepsis}}$ is
\[
  e_{\mathrm{sepsis}}^{\ast} :=
  \mathsf{HFD} \circ \mathsf{review}_\mathrm{att} \circ \mathsf{review}_\mathrm{res} \circ
  \mathsf{assist} \circ \mathsf{prep} \circ \mathsf{RB}.
\]
Lifting to $N \ltimes E$ and applying the confluence theorem for the combined reduction (Epistemic Algebra)
ensures that any admissible navigation of the conversational UI yields a unique normal form
equivalent to $e_{\mathrm{sepsis}}^{\ast}$ up to the RB--TE--HFD suffix.

\paragraph{Interaction Protocols admissibility.}
We interpret the Interaction Protocols axioms IL, AR, MD, and NTER as constraints on realizations of
$e_{\mathrm{sepsis}}$ in a concrete UI:
\begin{itemize}
  \item IL: the only interaction channels between AI and human tacit state are via explicit objects in
  $E$ (case summaries, drafts, candidate plans, reports). No morphisms $T \to E$ are introduced.
  \item AR: there is no morphism in the realization that effects clinical actions without passing
  through $\mathsf{HFD} \circ \mathsf{review}_\mathrm{att}$; the AI cannot place orders or write notes.
  \item MD: all model-dependent computation is represented by $\mathsf{assist}$ and its internal
  decomposition; tacit evaluation (clinician judgment) is localized at $\mathsf{review}_\mathrm{res}$ and
  $\mathsf{review}_\mathrm{att}$.
  \item NTER: the implementation forbids any arrow that would encode a tacit correctness oracle in
  explicit space (no reconstruction of $T$ from model outputs). In particular, the realization of
  $\mathsf{assist}$ is constrained to produce drafts and rationales, not authoritative decisions.
\end{itemize}
Under these conditions, any concrete realization of $e_{\mathrm{sepsis}}$ satisfies the Interaction Protocols
interface axioms and is therefore admissible.

\paragraph{Epistemic Spaces/Algebra well-typedness.}
In Epistemic Spaces, each generator is assigned an epistemic role: $\mathsf{prep}$ and $\mathsf{assist}$ are
explicit-computation arrows; $\mathsf{review}_\mathrm{res}$ and $\mathsf{review}_\mathrm{att}$ are explicit
presentation points at which tacit evaluation in $T$ occurs; $\mathsf{RB}$ and $\mathsf{HFD}$ are
boundary morphisms. The typing rules of the Epistemic Spaces ontology then derive a composite judgment
showing that $e_{\mathrm{sepsis}}^{\ast}$ is a well-typed explicit pipeline with correctness localized in
tacit post-composition at $\mathsf{review}_\mathrm{att}$. In Epistemic Algebra, the same term appears as a
normal form in $N \ltimes E$; compatibility of reduction with the semantics ensures that the
correctness predicate applied to $e_{\mathrm{sepsis}}^{\ast}$ coincides with the protocol-level notion of
FRFP-hallucination (a violation of $t \succeq \mathrm{Req}(e,c)$ at the moment of final acceptance).
Thus the sepsis consult example realizes a concrete conversational UI that is Interaction Protocols admissible,
Epistemic Spaces/Algebra well-typed, and corresponds to a specific normal form in the FRFP calculus.

\pagebreak

%\section{Pipeline Dynamics}
%\label{app:pipeline-dynamics}

\subsubsection{Purpose and Scope}
The discussion in the Appendix till now formalizes
\emph{static} correctness of individual FRFP pipelines, Epistemic Dynamics
analyzes what happens when such pipelines are iterated by the same
human--AI system over many turns. The focus is on:
\begin{enumerate}
  \item persistence and degradation of tacit context;
  \item explicit-only hallucinations and their containment;
  \item confidence inflation and requirement escalation;
  \item feedback-coupled degradation and accelerated hallucination;
  \item stochastic and finite-horizon pipelines.
\end{enumerate}
No collective or institutional mechanisms are assumed in this section;
these appear in latter.

Throughout, we fix the Epistemic structure and work
with the underlying tacit context space \(T\) and explicit artifact
space \(E\).

%======================================================================
\subsubsection{Tacit Context Space}
%======================================================================

Each human agent possesses a tacit context space \(T\), equipped with a
preorder \(\preceq\).

\begin{definition}[Tacit context preorder]
The binary relation \(\preceq\) on \(T\) is read as:
\[
t_1 \preceq t_2
\quad\text{``\(t_1\) is epistemically no stronger than \(t_2\)''}.
\]
We write \(t_1 \prec t_2\) if \(t_1 \preceq t_2\) and \(t_2 \not\preceq t_1\).
\end{definition}

No algebraic structure on \(T\) is assumed beyond this preorder. The
categorical tacit structure  (one-object category,
primitives \(\TE,\HFD\), correctness quotient \(\gamma\)) is still in
force, but for dynamic analysis we work primarily with the underlying
preordered set \((T,\preceq)\).

%======================================================================
\subsubsection{Context Degradation}
%======================================================================

Tacit context degrades over time due to forgetting, misplaced trust,
habituation, and other long-run effects.

\begin{definition}[Context degradation]
A \emph{context degradation} operator is a function
\[
\delta : T \to T
\]
modeling loss of grounding across turns. We assume:
\begin{enumerate}
  \item \textbf{Monotonicity}:
        \(\delta(t) \preceq t\) for all \(t \in T\).
  \item \textbf{Strict decay somewhere}:
        there exists \(t \in T\) with \(\delta(t) \prec t\).
  \item \textbf{Non-invertibility}:
        \(\delta\) has no right inverse; i.e., there is no function
        \(f:T\to T\) such that \(\delta \circ f = \mathrm{id}_T\).
\end{enumerate}
We write \(\delta^n\) for the \(n\)-fold iterate of \(\delta\).
\end{definition}

Intuitively, \(\delta\) expresses the part of tacit dynamics that is
\emph{purely degrading}. In this appendix we deliberately ignore
repair/refresh operations (e.g., reading textbooks, taking breaks) to
obtain worst-case inevitability results. Repair can be modeled by
separate operators \(\rho:T\to T\), but these are left for future work.

%======================================================================
\subsubsection{Dynamic Interpretation of Pipelines}
%======================================================================

We now interpret complete FRFP pipelines as transformations of tacit
context.

Let \(E_{\mathrm{sem}}\) be the explicit semantic category and \(\RB^\sharp : E_{\mathrm{sem}}\to T\) the semantic
backflow (Section~A.0.5). For each explicit pipeline \(e\in
E_{\mathrm{sem}}\), the Epistemic Space suffix \(\TE,\HFD\) yields a tacit
transformer
\[
t \longmapsto \HFD \circ \TE \circ \RB^\sharp(e,t),
\]
where we allow the tacit state \(t\) to modulate interpretation
(e.g., via background knowledge).

\begin{definition}[Pipeline transformer]
A \emph{complete FRFP pipeline} \(P\) at a given turn consists of:
\begin{itemize}
  \item an explicit artifact \(e \in E_{\mathrm{sem}}\) produced by the
        AI system via explicit computation and navigation;
  \item the canonical tacit suffix \(\RB \to \TE \to \HFD\);
  \item subsequent degradation \(\delta\).
\end{itemize}
It induces a tacit transformer
\[
\llbracket P \rrbracket_\delta : T \to T,\qquad
\llbracket P \rrbracket_\delta(t)
 :=
 \delta\bigl( \HFD\circ\TE\circ\RB^\sharp(e,t)\bigr).
\]
\end{definition}

\begin{definition}[Run]
A \emph{run} (or \emph{pipeline execution}) is a sequence
\[
t_0, t_1, t_2,\dots
\]
of tacit states such that there exist pipelines
\(P_0, P_1, P_2,\dots\) with
\[
t_{n+1}
  = \llbracket P_n \rrbracket_\delta(t_n)
\quad\text{for all } n \geq 0.
\]
\end{definition}

We will often also track the explicit artifacts \(e_n\) and credences
\(c_n\) associated with each \(P_n\); this is made precise in
Sections~\ref{sec:B7-confidence} and~\ref{sec:B15-stochastic}.

%======================================================================
\subsubsection{Requirements and Hallucination}
%======================================================================

We formalize hallucination in terms of a requirement map on explicit
artifacts.

\begin{definition}[Requirement map]
A \emph{requirement map} is a function
\[
\Req : E \times \mathbb{R}_{\ge 0} \to T
\]
assigning to each explicit artifact \(e \in E\) and tacit credence
\(c\ge 0\) a required tacit grounding \(\Req(e,c)\). We require:
\begin{enumerate}
  \item \textbf{Monotonicity in credence}:
        if \(c_1 \le c_2\) then
        \(\Req(e,c_1) \preceq \Req(e,c_2)\);
  \item \textbf{Nontriviality}:
        there exist artifacts whose requirements strictly exceed others
        in the preorder.
\end{enumerate}
When credence is suppressed, we write \(\Req(e)\) for a baseline
requirement (e.g., at default credence).
\end{definition}

\begin{definition}[Hallucination]
Let \(t\in T\) be the tacit state at evaluation time, \(e\in E\) an
explicit artifact, and \(c\ge 0\) the tacit credence assigned to
\(e\). A \emph{hallucination} occurs at this step if
\[
t \prec \Req(e,c).
\]
We say that the artifact \(e\) is \emph{hallucinated} (relative to
\((t,c)\)) in this case.
\end{definition}

\begin{remark}[Structural nature of hallucination]
In this formalism, hallucination is an \emph{implementation-agnostic}
relation between tacit state \(t\), requirement \(\Req(e,c)\), and
explicit artifact \(e\). It does not depend on the model's internal
computation or intent. Concrete architectures (e.g., LLMs) shape
\emph{how often} trajectories enter states with \(t \prec
\Req(e,c)\), but do not appear in the definition itself.
\end{remark}

%======================================================================
\subsubsection{Hallucination Inevitability}
%======================================================================

We now state a baseline hallucination inevitability theorem for
deterministic degradation dynamics.

To capture recurring nontrivial demands, we restrict to runs that keep
encountering artifacts whose requirements exceed a fixed lower bound.

\begin{definition}[Nontrivial requirement floor]
A requirement floor is a tacit state \(r\in T\). A run
\((t_n,e_n,c_n)_{n\ge 0}\) has \emph{recurring nontrivial demands}
relative to \(r\) if there exists an infinite subsequence
\((n_k)_{k\ge 0}\) such that
\[
\Req(e_{n_k},c_{n_k}) \succeq r
\quad\text{for all }k.
\]
\end{definition}

We also assume that degradation eventually pushes tacit states below
this floor.

\begin{assumption}[Baseline degradation below the requirement floor]
\label{assump:baseline-degradation}
There exists \(r\in T\) such that for every initial tacit state
\(t_0\in T\), there exists \(N\in\mathbb{N}\) with
\[
\delta^N(t_0) \prec r.
\]
\end{assumption}

This is a strengthening of ``strict decay somewhere'': it requires that
every trajectory ultimately falls below the floor \(r\).

\begin{theorem}[Hallucination inevitability]
\label{thm:baseline-inevitability}
Assume:
\begin{enumerate}
  \item context degradation \(\delta\) satisfies monotonicity,
        non-invertibility, and
        Assumption~\ref{assump:baseline-degradation};
  \item requirements admit a nontrivial floor \(r\) and are monotone in
        credence;
  \item a run \((t_n,e_n,c_n)_{n\ge 0}\) has recurring nontrivial
        demands relative to \(r\).
\end{enumerate}
Then every infinite run produces at least one hallucination; i.e., there
exists \(n\) with
\[
t_n \prec \Req(e_n,c_n).
\]
\end{theorem}

\begin{proof}[Proof sketch]
By Assumption~\ref{assump:baseline-degradation}, for any initial
\(t_0\) there exists \(N\) with \(\delta^N(t_0) \prec r\). Consider the
given run and the subsequence \((n_k)\) along which
\(\Req(e_{n_k},c_{n_k})\succeq r\). For sufficiently large \(k\),
monotonicity of \(\delta\) and recurrence of nontrivial demands imply
that
\[
t_{n_k} \preceq \delta^{m}(t_0) \prec r \preceq \Req(e_{n_k},c_{n_k})
\]
for some \(m\), so \(t_{n_k} \prec \Req(e_{n_k},c_{n_k})\), i.e., a
hallucination. A detailed proof can be given by induction along the
run, but the key idea is that degradation eventually drives tacit
grounding below a floor that requirements exceed infinitely often.
\end{proof}

\begin{remark}
The theorem applies to \emph{nontrivial, open-ended} task streams with
recurrent demands above a fixed floor. Trivial avoidance regimes (e.g.,
restricting to a tiny class of always-easy tasks or always answering
``I don't know'') fall outside the intended scope.
\end{remark}

%======================================================================
\subsubsection{Explicit-Only Hallucinations}
%======================================================================
We now isolate explicit artifacts that are produced entirely within explicit reasoning, but which can be
hallucinated relative to some tacit state encountered along a run.

\begin{definition}[Explicit-only hallucination]
An \emph{explicit-only hallucination} is an explicit artifact
\(e \in E\) such that:
\begin{enumerate}
  \item \(e\) is well-typed and syntactically valid as an FRFP explicit
        object;
  \item \(e\) is surface-plausible under an explicit plausibility
        functional (Section~\ref{sec:B7-confidence});
  \item there exists a tacit state \(t\) and credence \(c\) encountered
        in some run such that
        \(t \prec \Req(e,c)\).
\end{enumerate}
\end{definition}

\begin{proposition}[Containment but not elimination]
Nothing in Epistemic Space or Epistemic Algebra forbids the production of explicit-only
hallucinations. Eliminating them entirely would require either:
\begin{itemize}
  \item an explicit correctness oracle, or
  \item a simulation of tacit evaluation inside explicit space,
\end{itemize}
both of which would violate Explicit--Tacit Separation (ETS) and
Human-Exclusive Grounding (HEG).

The explicit--tacit separation guarantees \emph{containment}, not
\emph{elimination}: explicit-only hallucinations cannot become
correctness judgments without passing through \(\RB\), but they may
still be produced arbitrarily often within explicit space.
\end{proposition}

%======================================================================
\subsubsection{Confidence Inflation}
\label{sec:B7-confidence}
%======================================================================

We model how plausible explicit artifacts raise tacit credence even
when they do not increase grounding.

\begin{definition}[Explicit plausibility]
An \emph{explicit plausibility functional} is a map
\[
\pi : E \to [0,1],
\]
measuring surface coherence, fluency, or apparent authority of explicit
artifacts.
\end{definition}

We extend tacit state with a credence component.

\begin{definition}[Extended tacit state]
Let
\[
T^e := T \times \mathbb{R}_{\ge 0}
\]
be the space of extended tacit states \((t,c)\), where \(t\in T\) is
grounding and \(c\) is tacit credence (how settled or reliable the
agent treats the current artifact).
\end{definition}

\begin{definition}[Confidence inflation]
A \emph{confidence inflation operator} is a map
\[
\Phi : (T^e \times E) \to T^e,\qquad
\Phi\bigl( (t,c), e \bigr)
 :=
 (t,\; c + \varphi(\pi(e))),
\]
where \(\varphi:[0,1]\to \mathbb{R}_{\ge 0}\) is monotone increasing.

We require that \(\Phi\) does not increase tacit grounding:
\[
\Phi\bigl((t,c),e\bigr) = (t',c') \quad\Rightarrow\quad t' = t.
\]
\end{definition}

Thus confidence inflation changes only the agent's interpretive stance
(higher credence), not the underlying grounding \(t\).

\begin{remark}
Confidence inflation captures the structural asymmetry between
``feeling confident'' and ``being grounded''. Systems optimized for
producing well-formed, plausible outputs (e.g., LLMs) inevitably
exhibit this asymmetry: they can increase credence by producing
plausible artifacts without necessarily increasing grounding.
\end{remark}

%======================================================================
\subsubsection{Requirement Escalation Under Plausibility}
%======================================================================

Credence affects what requirements are appropriate: asserting
high-credence claims demands stronger grounding.

\begin{definition}[Requirement escalation]
The requirement map \(\Req : E \times \mathbb{R}_{\ge 0} \to T\) is
\emph{escalatory} if it is monotone in credence, i.e.,
\[
c_1 \le c_2 \quad\Rightarrow\quad
\Req(e,c_1) \preceq \Req(e,c_2),
\]
for all \(e\in E\). Repeated exposure to highly plausible artifacts
(large \(\pi(e)\)) can then raise requirements even if the explicit
content \(e\) is unchanged.
\end{definition}

Combining confidence inflation and requirement escalation, we see that
the same explicit artifact can acquire higher requirements over time as
agents become more confident in it.

%======================================================================
\subsubsection{Feedback-Coupled Degradation}
%======================================================================

We now allow degradation to depend on credence, introducing a feedback
loop: overconfidence reduces skepticism and checking, accelerating loss
of grounding.

\begin{definition}[Feedback-coupled degradation]
A \emph{feedback-coupled degradation} operator is a function
\[
\delta^\star : T^e \to T,\qquad
\delta^\star(t,c) = \delta_c(t),
\]
where \(\delta_c : T \to T\) is a family of degradation operators
indexed by credence \(c\), satisfying:
\begin{enumerate}
  \item \textbf{Credence monotonicity}:
        for fixed \(t\), if \(c_1 \le c_2\) then
        \(\delta^\star(t,c_2) \preceq \delta^\star(t,c_1)\);
        i.e., higher confidence weakens effective grounding.
  \item \textbf{Consistency with baseline degradation}:
        \(\delta^\star(t,0) = \delta(t)\) for all \(t\).
\end{enumerate}
\end{definition}

\begin{definition}[Extended run with feedback]
An \emph{extended run} is a sequence
\[
(t_0,c_0,e_0), (t_1,c_1,e_1), \dots
\]
such that, at each step \(n\),
\begin{align*}
(t_n^+, c_n^+) &= \Phi\bigl((t_n,c_n), e_n\bigr)
      &&\text{(confidence inflation)},\\
t_{n+1} &= \delta^\star(t_n^+, c_n^+)
      &&\text{(feedback-coupled degradation)},\\
c_{n+1} &\text{ chosen by tacit update (e.g., reset or carry-over)}.
\end{align*}
\end{definition}

The exact rule for updating \(c_{n+1}\) can be left abstract; all we
require is that confidence inflation can push credence arbitrarily high
along some runs.

%======================================================================
\subsubsection{Rate of Context Loss}
%======================================================================}

We can analyze the rate at which tacit grounding is lost under
feedback-coupled degradation.

\begin{definition}[Grounding measure and loss rate]
Let \(\mu : T \to \mathbb{R}_{\ge 0}\) be any monotone measure of tacit
grounding (used only for analysis). For an extended run
\((t_n,c_n,e_n)\), define the instantaneous loss rate at turn \(n\) by
\[
\rho_n := \mu(t_n) - \mu\bigl(\delta^\star(t_n,c_n)\bigr).
\]
\end{definition}

\begin{proposition}[Loss rate increases with credence]
Under the credence-monotonicity assumption for \(\delta^\star\) and
monotonicity of \(\mu\), the loss rate \(\rho_n\) is non-decreasing in
\(c_n\). In particular, higher credence does not slow context loss and
may accelerate it.
\end{proposition}

\begin{proof}
If \(c_1 \le c_2\), then
\(\delta^\star(t,c_2) \preceq \delta^\star(t,c_1)\) and
monotonicity of \(\mu\) yields
\(\mu(\delta^\star(t,c_2)) \le \mu(\delta^\star(t,c_1))\). Thus
\[
\mu(t) - \mu(\delta^\star(t,c_2))
\;\ge\;
\mu(t) - \mu(\delta^\star(t,c_1)),
\]
so \(\rho_n\) increases with \(c_n\).
\end{proof}

%======================================================================
\subsubsection{Accelerated Hallucination}
%======================================================================

We now combine confidence inflation, requirement escalation, and
feedback-coupled degradation.

\begin{definition}[Feedback-coupled update]
Along an extended run, the joint update
\[
(t_{n+1}, c_{n+1})
 =
\Bigl(
  \delta^\star\bigl( t_n, c_n + \varphi(\pi(e_n))\bigr),\;
  c_{n+1}
\Bigr)
\]
captures the combined effect of:
\begin{itemize}
  \item confidence inflation
        \((t_n, c_n) \mapsto (t_n, c_n + \varphi(\pi(e_n)))\);
  \item degradation \(\delta^\star\), which is more severe at higher
        credence.
\end{itemize}
Grounding weakens faster while requirements grow stronger (via
\(\Req(e_n,c_n)\)).
\end{definition}

This creates a widening gap between available grounding \(t_n\) and
demanded grounding \(\Req(e_n,c_n)\), leading to accelerated
hallucination.

%======================================================================
\subsubsection{Accelerated Hallucination Inevitability}
%======================================================================

We can now state a strengthened inevitability theorem in the presence
of feedback.

\begin{assumption}[Unbounded credence under plausibility]
There exists \(\bar{\pi} \in (0,1]\) such that for any initial credence
\(c_0\) and any infinite sequence of artifacts with
\(\pi(e_n) \ge \bar{\pi}\) infinitely often, the accumulated credence
\[
c_n^{+} := c_0 + \sum_{k=0}^{n-1} \varphi(\pi(e_k))
\]
tends to \(+\infty\) along a subsequence.
\end{assumption}

\begin{theorem}[Accelerated hallucination inevitability]
\label{thm:accelerated-inevitability}
Assume:
\begin{enumerate}
  \item baseline context degradation \(\delta\) is monotone and
        non-invertible and satisfies
        Assumption~\ref{assump:baseline-degradation};
  \item confidence inflation \(\Phi\) with
        \(\varphi(\bar{\pi}) > 0\) for some \(\bar{\pi}\in(0,1]\);
  \item requirement escalation: \(\Req\) is monotone in credence;
  \item feedback-coupled degradation \(\delta^\star\) is credence-monotone
        and reduces to \(\delta\) at zero credence;
  \item the run contains infinitely many \emph{surface-plausible}
        explicit artifacts with \(\pi(e_n) \ge \bar{\pi}\) and recurring
        nontrivial demands relative to some floor \(r\).
\end{enumerate}
Then:
\begin{enumerate}
  \item hallucination occurs in finite time; i.e., there exists
        \(n < \infty\) such that \(t_n \prec \Req(e_n,c_n)\);
  \item the minimum such \(n\) is strictly smaller than in the baseline
        degradation model without feedback (for any fixed initial
        state and pipeline sequence), in the sense that feedback cannot
        increase time-to-hallucination.
\end{enumerate}
\end{theorem}

\begin{proof}[Proof sketch]
By unbounded credence under plausibility and the presence of infinitely
many plausible artifacts, credence \(c_n\) grows without bound along
some subsequence. By credence-monotonicity of \(\delta^\star\) and
requirement escalation, higher credence simultaneously:
\begin{itemize}
  \item increases required grounding \(\Req(e_n,c_n)\), and
  \item accelerates degradation of available grounding via
        \(\delta^\star\).
\end{itemize}
Thus the gap
\[
t_n \prec \Req(e_n,c_n)
\]
is reached in finite time, and strictly earlier than under the baseline
degradation \(\delta\) alone, where the same sequence of artifacts
would produce slower loss of grounding and lower requirements.
A fully formal proof can be obtained by coupling two runs (with and
without feedback) and using the monotone measure \(\mu\) to compare
times at which \(\mu(t_n)\) falls below \(\mu(\Req(e_n,c_n))\).
\end{proof}


\subsubsection{Degradation--repair dynamics and entropy drift}
\label{app:delta-rho}
We extend the Epistemic Dynamics to include abstract repair. The aim is to state conditions under
which repair delays or reshapes, but does not remove, long-run degradation.

\begin{definition}[Entropy functional on $T$]
An \emph{entropy functional} on the tacit category $T$ is a map $S : T \to \mathbb{R}$ such that
\begin{enumerate}
  \item $t \preceq t' \Rightarrow S(t) \le S(t')$ for all $t,t' \in T$; and
  \item there exists $S_{\mathrm{req}} \in \mathbb{R}$ such that $S(t) \le S_{\mathrm{req}}$ whenever $t$
  meets the relevant requirement floor(s), and $S(t) > S_{\mathrm{req}}$ implies that accepting a
  nontrivial demand at $t$ yields an FRFP-hallucination in the sense $t_n \prec \mathrm{Req}(e_n,c_n)$.
\end{enumerate}
Thus $S(t)$ measures a grounding deficit: larger values correspond to more degraded tacit states.
\end{definition}

We consider runs $(t_n)_{n \ge 0}$ in $T$ evolving under a mixture of degradation and repair, and
write $S_n := S(t_n)$. Let $D_n,R_n \in \{0,1\}$ indicate, respectively, whether step $n$ processes a
nontrivial demand and whether a repair operation is applied.

\begin{assumption}[Task-induced degradation]\label{ass:delta-entropy}
There exists $\varepsilon > 0$ such that for all $t \in T$ and all nontrivial demand steps to which the
Epistemic Dynamics apply,
\[
  S(\delta(t)) \;\ge\; S(t) + \varepsilon.
\]
\end{assumption}

\begin{assumption}[Bounded repair]\label{ass:rho-entropy}
There exists $B < \infty$ such that for all admissible repair operators $\rho : T \to T$ and all
$t \in T$,
\[
  S(\rho(t)) \;\ge\; S(t) - B.
\]
\end{assumption}

Combining these, for each $N$ we have
\begin{equation}\label{eq:Sn-balance}
  S_N \;\ge\; S_0 + \varepsilon \sum_{n=1}^N D_n \;-\; B \sum_{n=1}^N R_n,
\end{equation}
where the sums record entropy produced by nontrivial tasks and removed by repair.
\begin{theorem}[Deterministic entropy drift and inevitability]\label{thm:delta-rho-deterministic}
Suppose Assumptions~\ref{ass:delta-entropy} and~\ref{ass:rho-entropy} hold and that
$(D_n,R_n)_{n \ge 1}$ satisfies
\[
  \liminf_{N \to \infty} \frac{1}{N}\sum_{n=1}^N D_n \;\ge\; \alpha > 0,
  \qquad
  \limsup_{N \to \infty} \frac{1}{N}\sum_{n=1}^N R_n \;\le\; \beta < \infty.
\]
Assume further that there exists $\gamma > 0$ with
\[
  \varepsilon \alpha - B \beta \;\ge\; \gamma > 0.
\]
Then there exists $N_0$ such that $S_N \ge S_0 + \gamma N$ for all $N \ge N_0$. In particular, for
any finite $S_{\mathrm{req}}$ the hitting time
\[
  \sigma_{\mathrm{req}} := \inf\{n \ge 0 : S_n > S_{\mathrm{req}}\}
\]
is finite. If nontrivial demands recur after $\sigma_{\mathrm{req}}$ and the protocol sometimes accepts
artifacts at those steps with credence corresponding to $S_{\mathrm{req}}$, then along some admissible
run an FRFP-hallucination event $t_n \prec \mathrm{Req}(e_n,c_n)$ occurs.
\end{theorem}

\begin{proof}[Proof sketch]
Divide~\eqref{eq:Sn-balance} by $N$ and apply the liminf/limsup bounds on the empirical
frequencies of $D_n$ and $R_n$ to obtain
\[
  \liminf_{N \to \infty} \frac{S_N - S_0}{N}
  \;\ge\; \varepsilon \alpha - B \beta \;\ge\; \gamma > 0.
\]
Hence $S_N \ge S_0 + \gamma N$ for all large $N$, so $S_N$ eventually exceeds any finite
$S_{\mathrm{req}}$ and $\sigma_{\mathrm{req}} < \infty$. The hallucination clause follows from the
definition of $S_{\mathrm{req}}$ and the criterion $t_n \prec \mathrm{Req}(e_n,c_n)$ at an accepted
nontrivial demand.
\end{proof}
\begin{theorem}[Stochastic entropy drift and almost-sure inevitability]\label{thm:delta-rho-stochastic}
Suppose Assumptions~\ref{ass:delta-entropy} and~\ref{ass:rho-entropy} hold and that
$(D_n,R_n)_{n \ge 1}$ is stationary and ergodic with
\[
  \mathbb{E}[D_1] = \bar{d} > 0,
  \qquad
  \mathbb{E}[R_1] = \bar{r} < \infty.
\]
Assume that
\[
  \varepsilon \bar{d} - B \bar{r} \;=\; \gamma > 0.
\]
Then almost surely
\[
  \liminf_{N \to \infty} \frac{S_N - S_0}{N} \;\ge\; \gamma > 0,
\]
so $S_N \to +\infty$ almost surely. In particular, the hitting time
$\sigma_{\mathrm{req}} := \inf\{n \ge 0 : S_n > S_{\mathrm{req}}\}$ is almost surely finite. Under the
same acceptance and recurrence conditions as in Theorem~\ref{thm:delta-rho-deterministic}, the
hallucination time $\tau(\omega)$ in the Epistemic Dynamics stochastic semantics satisfies $P(\tau < \infty) = 1$.
\end{theorem}

\begin{proof}[Proof sketch]
By the ergodic theorem, almost surely
\[
  \lim_{N \to \infty} \frac{1}{N}\sum_{n=1}^N D_n = \bar{d},
  \qquad
  \lim_{N \to \infty} \frac{1}{N}\sum_{n=1}^N R_n = \bar{r}.
\]
Together with~\eqref{eq:Sn-balance} this yields
\[
  \liminf_{N \to \infty} \frac{S_N - S_0}{N} \;\ge\; \varepsilon \bar{d} - B \bar{r} = \gamma > 0
  \quad\text{a.s.}
\]
Thus $S_N \to +\infty$ almost surely and $\sigma_{\mathrm{req}} < \infty$ almost surely; whenever a
nontrivial demand is accepted at or after $\sigma_{\mathrm{req}}$, the hallucination criterion is met at
some finite step, so $P(\tau < \infty) = 1$.
\end{proof}
\begin{remark}[Counter-scenarios and safe horizons]\label{rem:delta-rho-counter}
The drift condition in Theorems~\ref{thm:delta-rho-deterministic} and~\ref{thm:delta-rho-stochastic}
formalizes a high-load, bounded-repair regime: nontrivial demands occur with positive asymptotic
frequency, repair has bounded one-step effect, and net entropy production $\varepsilon \bar{d}$
exceeds maximal entropy removal $B \bar{r}$. In this regime, repair changes the timing and shape
of hallucination risk but does not eliminate it.

Several counter-scenarios lie at or beyond this domain. In \emph{balanced or negative drift}
regimes, repair investment is large enough that $B \bar{r} \ge \varepsilon \bar{d}$ and $S_N$ can remain
tight or decrease. In \emph{unbounded repair} regimes, one allows repair moves with no global
bound $B$ (e.g.\ costless, arbitrarily frequent ``oracle refresh'' operations or always-perfect staff
replacement), violating Assumption~\ref{ass:rho-entropy}. In \emph{low-demand} regimes, such as
rare emergency committees, the effective demand frequency $\bar{d}$ may be so small that modest
repair keeps $S$ below $S_{\mathrm{req}}$ over relevant horizons.

Within the high-load, bounded-repair regime, improving repair capacity modifies the hazard profile
$(h_n)_{n \ge 0}$ and increases the safe-horizon functional $H_P(\varepsilon)$, but for any fixed
stochastic FRFP pipeline $P$ satisfying the drift condition, $H_P(\varepsilon)$ remains finite for every
$\varepsilon \in (0,1)$: one cannot obtain a global guarantee of ``no hallucinations ever'' over
unbounded horizons by bounded repair alone.
\end{remark}

%======================================================================
\subsubsection{Stochastic Pipelines and Their Impact}
\label{sec:B15-stochastic}
%======================================================================

We extend the preceding results to stochastic FRFP pipelines.

\subsubsubsection*{Trajectories and probability measures}

\begin{definition}[Trajectory]
Fix an initial extended state \((t_0,c_0,e_0)\in T\times
\mathbb{R}_{\ge 0}\times E\). A \emph{trajectory} (or \emph{run}) is a
sequence
\[
\omega
 =
\bigl( (t_0,c_0,e_0), (t_1,c_1,e_1), (t_2,c_2,e_2), \dots \bigr)
\]
such that each \((t_{n+1},c_{n+1},e_{n+1})\) is obtained from
\((t_n,c_n,e_n)\) by one application of the FRFP pipeline, tacit
evaluation, confidence update, and degradation (possibly
feedback-coupled).
\end{definition}

Let \(\Omega\) denote the set of all such trajectories. Let
\(\Omega_{\mathrm{allowed}}\subseteq \Omega\) be the subset of
trajectories that satisfy the structural assumptions of this appendix
(degradation, recurring requirements, etc.).

We consider two equivalent presentations of stochasticity:
\begin{itemize}
  \item \emph{Random-seed semantics:} extend the state space with a
        random seed \(r\), and suppose there is a deterministic update
        map
        \[
        F: T^e\times E \times R \to T^e \times E \times R
        \]
        such that each \(r\in R\) induces a unique trajectory
        \(\omega(r)\in\Omega\).
  \item \emph{Markov-kernel semantics:} one-step evolution is described
        by a Markov kernel
        \[
        K : T^e\times E \to \mathcal{D}(T^e\times E),
        \]
        where \(\mathcal{D}(\cdot)\) indicates probability
        distributions; a trajectory \(\omega\) is then a sample path of
        the stochastic process generated by \(K\).
\end{itemize}

In either case, a \emph{stochastic FRFP pipeline} is identified with a
probability measure \(\mathbb{P}\) on \(\Omega\) induced by the
underlying randomness.

\subsubsubsection*{Hallucination events}

Recall the hallucination criterion: hallucination at step \(n\) occurs
when
\[
t_n \prec \Req(e_n,c_n).
\]

\begin{definition}[Hallucination events]
For a trajectory \(\omega\in\Omega\) and \(n\ge 0\), define the
indicator
\[
H_n(\omega)
 :=
 \mathbf{1}\bigl[t_n \prec \Req(e_n,c_n)\bigr] \in \{0,1\}.
\]
Define the \emph{eventual hallucination indicator}
\[
H(\omega)
 :=
 \mathbf{1}\bigl[ \exists n\ge 0 : H_n(\omega)=1 \bigr].
\]
Thus \(H(\omega)=1\) iff the trajectory \(\omega\) contains at least one
hallucination.
\end{definition}

\begin{assumption}[Pathwise inevitability]
\label{assump:pathwise-inevitability}
Suppose the structural assumptions of Sections~B.1--B.12 hold (on
degradation, requirements, and feedback), and that the baseline or
accelerated hallucination inevitability theorem has been established in
the deterministic setting. We assume:
\[
\forall \omega \in \Omega_{\mathrm{allowed}}:\quad H(\omega)=1.
\]
That is, every allowed infinite trajectory eventually produces a
hallucination.
\end{assumption}

\begin{theorem}[Almost-sure inevitability under stochastic pipelines]
\label{thm:stochastic-inevitability}
Let \(\mathbb{P}\) be any probability measure on \(\Omega\) whose
support is contained in \(\Omega_{\mathrm{allowed}}\). Then
\[
\mathbb{P}\bigl( \text{eventual hallucination} \bigr)
 =
\mathbb{P}\bigl( \{\omega : H(\omega)=1\} \bigr)
 =
1.
\]
\end{theorem}

\begin{proof}
By Assumption~\ref{assump:pathwise-inevitability}, \(H(\omega)=1\) for
all \(\omega\in\Omega_{\mathrm{allowed}}\). Since \(\mathbb{P}\) is
supported on \(\Omega_{\mathrm{allowed}}\), we have \(H(\omega)=1\) for
\(\mathbb{P}\)-almost every \(\omega\). Hence
\[
\mathbb{P}(\text{eventual hallucination})
 =
\int_{\Omega} H(\omega)\, d\mathbb{P}(\omega)
 =
\int_{\Omega} 1\, d\mathbb{P}(\omega)
 =
1.
\]
\end{proof}

Thus passing from deterministic to stochastic pipelines does not weaken
the inevitability result: randomness merely selects \emph{which}
hallucinating trajectory is realized, not whether hallucination occurs.

\subsubsubsection*{Probabilistic confluence}

Stochasticity breaks single-outcome confluence but yields a natural
probabilistic analogue.

\begin{definition}[Output distribution and probabilistic confluence]
Fix a time \(n\ge 0\) and initial state \((t_0,c_0,e_0)\). Let
\(\mathbb{P}\) be a stochastic FRFP pipeline with support in
\(\Omega_{\mathrm{allowed}}\). The \emph{time-\(n\) output
distribution} is the pushforward measure
\[
\mu_n^{\mathbb{P}} := (e_n)_\# \mathbb{P},
\]
where \(e_n:\Omega\to E\) maps a trajectory to its explicit artifact at
step \(n\).

Two stochastic implementations \(\mathbb{P}_1,\mathbb{P}_2\) are
\emph{probabilistically confluent at time \(n\)} if they induce the same
output distribution \(\mu_n^{\mathbb{P}_1}=\mu_n^{\mathbb{P}_2}\) on
\(E\).
\end{definition}

If both implementations correspond to the same Markov kernel on
\(T^e\times E\), they are probabilistically confluent at every time
\(n\). This notion provides a distributional counterpart to the
canonical-form and confluence properties discussed earlier.

%======================================================================
\subsubsection{Finite-Horizon Pipelines and Their Impact}
%======================================================================

We now study the effect of restricting attention to finite prefixes of
trajectories, both in deterministic and stochastic settings.

\subsubsubsection*{Hallucination time}

\begin{definition}[Hallucination time]
For a trajectory \(\omega\in\Omega\), define the \emph{hallucination
time}
\[
\tau(\omega)
 :=
 \inf\{ n \ge 0 : t_n \prec \Req(e_n,c_n) \},
\]
with the convention \(\tau(\omega)=+\infty\) if the set is empty (i.e.,
if no hallucination ever occurs).
\end{definition}

Thus \(\tau(\omega) < \infty\) iff \(\omega\) contains at least one
hallucination. In terms of Definition~B.13.2, \(H(\omega) =
\mathbf{1}[\tau(\omega) < \infty]\).

\begin{lemma}[Pathwise inevitability and finite hallucination time]
Under Assumption~\ref{assump:pathwise-inevitability} (pathwise
inevitability), we have
\[
\forall \omega \in \Omega_{\mathrm{allowed}}:\quad \tau(\omega) < \infty.
\]
\end{lemma}

\begin{proof}
By Assumption~\ref{assump:pathwise-inevitability}, for each
\(\omega\in\Omega_{\mathrm{allowed}}\) there exists some \(n\) with
\(H_n(\omega)=1\). By definition of \(\tau\), this implies that
\(\tau(\omega)\) is the infimum of a nonempty subset of \(\mathbb{N}\),
hence finite.
\end{proof}

A \emph{finite horizon} of length \(N\) corresponds to observing only
the prefix
\[
\omega_{\le N}
 :=
((t_0,c_0,e_0),\dots,(t_N,c_N,e_N)).
\]
We say that a trajectory \emph{hallucinates within horizon \(N\)} if
\(\tau(\omega)\le N\).

\begin{remark}[No universal finite bound in general]
Pathwise inevitability (\(\tau(\omega) < \infty\) for all allowed
\(\omega\)) does not imply the existence of a universal finite upper
bound \(N^\ast\) such that \(\tau(\omega)\le N^\ast\) for all
\(\omega\in\Omega_{\mathrm{allowed}}\). In general, for any fixed
\(N\) there may exist trajectories with \(\tau(\omega) > N\).
\end{remark}

\subsubsubsection*{A finite-time bound under strengthened assumptions}

If desired, one can obtain a finite-time upper bound under stronger
assumptions.

\begin{assumption}[Uniform loss and requirement floor]
\label{assump:uniform-loss}
Suppose there exist:
\begin{itemize}
  \item a grounding measure \(\mu:T\to\mathbb{R}_{\ge 0}\);
  \item a requirement level \(r\in T\) and constant \(\alpha>0\) with
        \(\mu(r)=\alpha\);
  \item a constant \(\varepsilon>0\);
\end{itemize}
such that, for every allowed trajectory
\(\omega=((t_n,c_n,e_n))_{n\ge 0}\):
\begin{enumerate}
  \item \emph{Monotone loss:} \(\mu(t_{n+1})\le\mu(t_n)\) for all
        \(n\).
  \item \emph{Uniform loss above \(r\):} whenever \(\mu(t_n)>\alpha\),
        we have
        \[
        \mu(t_n) - \mu(t_{n+1}) \ge \varepsilon.
        \]
  \item \emph{Recurrent requirement:} there exists an infinite sequence
        \((n_k)\) with \(\Req(e_{n_k},c_{n_k})\succeq r\) for all
        \(k\).
\end{enumerate}
\end{assumption}

\begin{lemma}[Finite-time bound under uniform loss]
\label{lem:finite-time-bound}
Under Assumption~\ref{assump:uniform-loss}, every allowed trajectory
\(\omega\) satisfies
\[
\tau(\omega)
 \le
N^\ast
 :=
\left\lceil \frac{\mu(t_0) - \alpha}{\varepsilon} \right\rceil.
\]
That is, a hallucination must occur by time \(N^\ast\).
\end{lemma}

\begin{proof}
By monotone loss and the uniform-loss condition, as long as
\(\mu(t_n)>\alpha\),
\[
\mu(t_{n+1}) \le \mu(t_n) - \varepsilon.
\]
By induction, for any \(n\) with this property,
\[
\mu(t_n) \le \mu(t_0) - n\varepsilon.
\]
For \(n = N^\ast\) we obtain \(\mu(t_{N^\ast}) \le \alpha =
\mu(r)\). By recurrence of requirements, there exists some \(k\) with
\(n_k \le N^\ast\) and \(\Req(e_{n_k},c_{n_k}) \succeq r\). Using
monotonicity of \(\mu\) with respect to \(\preceq\), we conclude that
\[
t_{n_k} \preceq r \preceq \Req(e_{n_k},c_{n_k}),
\]
and in fact \(t_{n_k} \prec \Req(e_{n_k},c_{n_k})\) for nontrivial
requirements, so \(\tau(\omega)\le n_k \le N^\ast\).
\end{proof}

These stronger conditions are not needed for the main inevitability
results, but they demonstrate that finite-time bounds are available in
regimes with sufficiently uniform context loss and recurrent demands.

\subsubsubsection*{Stochastic finite horizons}

Returning to the stochastic setting, we can describe how
finite-horizon hallucination probabilities behave.

\begin{definition}[Finite-horizon hallucination probability]
Fix a horizon \(N\in\mathbb{N}\). Define the event
\[
A_N := \{\omega\in\Omega : \tau(\omega) \le N\},
\]
and the corresponding probability
\[
p_N := \mathbb{P}(A_N)
     = \mathbb{P}(\tau \le N).
\]
Thus \(p_N\) is the probability that a random run hallucinates within
the first \(N\) steps.
\end{definition}

\begin{proposition}[Monotone convergence of finite-horizon probabilities]
Assume almost-sure inevitability
(Theorem~\ref{thm:stochastic-inevitability}), i.e.,
\(\mathbb{P}(\tau<\infty)=1\). Then:
\begin{enumerate}
  \item The sequence \((p_N)_{N\ge 0}\) is non-decreasing:
        \(p_N \le p_{N+1}\) for all \(N\).
  \item \(\displaystyle \lim_{N\to\infty} p_N = 1\).
\end{enumerate}
\end{proposition}

\begin{proof}
For each \(\omega\), the events \(\{\tau(\omega)\le N\}\) are nested:
\(\{\tau\le N\}\subseteq\{\tau\le N+1\}\). Hence
\(A_N\subseteq A_{N+1}\) and \(p_N\le p_{N+1}\). Moreover,
\(\{\tau<\infty\}=\bigcup_{N} A_N\), so by continuity of probability
from below,
\[
\mathbb{P}(\tau<\infty)
 =
\lim_{N\to\infty} \mathbb{P}(A_N)
 =
\lim_{N\to\infty} p_N.
\]
By assumption, \(\mathbb{P}(\tau<\infty)=1\), so
\(\lim_{N\to\infty} p_N=1\).
\end{proof}

\begin{proposition}[Multiple runs over a fixed horizon]
Fix a horizon \(N\) and consider \(K\in\mathbb{N}\) independent
trajectories \(\omega^{(1)},\dots,\omega^{(K)}\), with corresponding
hallucination times \(\tau^{(i)}\). Then:
\begin{enumerate}
  \item The probability that no trajectory hallucinates within horizon
        \(N\) is
        \[
        \mathbb{P}\bigl(\forall i,\ \tau^{(i)} > N\bigr)
         = (1-p_N)^K.
        \]
  \item The probability that at least one trajectory hallucinates
        within horizon \(N\) is
        \[
        1 - (1-p_N)^K.
        \]
\end{enumerate}
\end{proposition}

\begin{proof}
By independence,
\[
\mathbb{P}(\tau^{(i)} > N) = 1 - p_N
\]
for each \(i\), so
\[
\mathbb{P}\bigl(\forall i,\ \tau^{(i)} > N\bigr)
 =
\prod_{i=1}^K (1-p_N)
 =
(1-p_N)^K.
\]
The complementary event has probability \(1-(1-p_N)^K\).
\end{proof}

These results show how inevitability manifests at finite horizons: for
any fixed horizon, hallucination remains probabilistic (\(p_N<1\) in
general), but as the horizon length or number of runs increases, the
probability that hallucination has occurred approaches~1.

\medskip

Future work lifts these individual dynamics to populations of agents and
explicit consensus mechanisms.
% ------------------------------------------------------------------
% New material after Proposition B.38 in Appendix B
% ------------------------------------------------------------------

\subsubsection*{B.16 Hazard Functions and Safe Horizons}

We now refine finite-horizon analysis by introducing per-step hazard and a derived notion of
``safe horizon'' for a given stochastic FRFP pipeline.

\begin{definition}[Per-step hazard]\label{def:hazard}
Let $\tau : \Omega \to \mathbb{N} \cup \{\infty\}$ be the hallucination time of Definition~B.31,
and let $P$ be a stochastic FRFP pipeline with associated distribution of $\tau$.
For each $n \in \mathbb{N}$ with $P(\tau \ge n) > 0$ we define the (discrete) hazard at step $n$ by
\[
  h_n := P(\tau = n \mid \tau \ge n)
       = \frac{P(\tau = n)}{P(\tau \ge n)}.
\]
If $P(\tau \ge n) = 0$ we set $h_n := 0$.
\end{definition}

Intuitively, $h_n$ is the conditional probability that a run first hallucinates at step $n$,
given that it has not hallucinated earlier. The sequence $(h_n)_{n \ge 0}$ is a pipeline-level
summary of how risk is distributed along the trajectory.

\begin{definition}[Finite-horizon survival and hazard]\label{def:survival}
For $N \in \mathbb{N}$ let
\[
  S_N := P(\tau > N) = 1 - p_N
\]
be the survival probability at horizon $N$, where $p_N = P(\tau \le N)$ is the finite-horizon
hallucination probability from Definition~B.36. Whenever $P(\tau \ge n) > 0$ for all $n$,
the survival probabilities admit the product representation
\[
  S_N = \prod_{n=0}^{N} (1 - h_n),
\]
so that the finite-horizon risk $p_N = 1 - S_N$ can be understood as the cumulative effect
of per-step hazards.
\end{definition}

This representation makes explicit the trade-off between (i) the number of steps in a pipeline
and (ii) the hazard profile $(h_n)$ induced by its design (e.g. decomposition, tool calls,
or repeated self-checks).

\begin{definition}[Safe horizon at tolerance $\varepsilon$]\label{def:safe-horizon}
Fix a tolerance level $\varepsilon \in (0,1)$. The $\varepsilon$-safe horizon of a stochastic
FRFP pipeline $P$ is defined as
\[
  H_P(\varepsilon)
  := \sup\{\,N \in \mathbb{N} : p_N \le \varepsilon\,\}
  \in \mathbb{N} \cup \{\infty\}.
\]
When $H_P(\varepsilon)$ is finite, any horizon $N \le H_P(\varepsilon)$ satisfies
$P(\tau \le N) \le \varepsilon$, i.e. hallucination within that horizon is at most
$\varepsilon$-likely.
\end{definition}

\begin{remark}[Design interpretation]
The safe-horizon functional $H_P(\varepsilon)$ isolates a pipeline-level design parameter:
for a given stochastic implementation $P$ and admissible risk budget $\varepsilon$, it
characterizes how long a single run can safely be extended before hallucination becomes
probable. Changes to decomposition, step-ordering, or intermediate checks can therefore
be compared in terms of their effect on $(h_n)$ and $H_P(\varepsilon)$, even when the
underlying explicit pipeline is statically equivalent in $N \ltimes E$.
\end{remark}

Combining Definition~\ref{def:safe-horizon} with Proposition~B.38, we can also analyze
restart and parallelization policies: for example, running $K$ independent instances of a
pipeline up to horizon $N \le H_P(\varepsilon)$ yields hallucination probability
$1 - (1 - p_N)^K$, which can be bounded using the same tolerance parameter.

\subsubsection*{B.17 Dynamic Refinement of Pipelines}

We next introduce a notion of dynamic refinement between stochastic FRFP pipelines that
compares their hallucination-time distributions.

\begin{definition}[Dynamic refinement]\label{def:dynamic-refinement}
Let $P_1, P_2$ be stochastic FRFP pipelines on the same state space and admissible trajectory
set $\Omega_{\mathrm{allowed}}$, with induced hallucination times $\tau_1, \tau_2$ and
finite-horizon probabilities
\[
  p_N^{(i)} := P_i(\tau_i \le N), \quad N \in \mathbb{N}, \; i \in \{1,2\}.
\]
We say that $P_2$ is a \emph{dynamic refinement} (or \emph{dynamically safer variant}) of $P_1$
if for all $N \in \mathbb{N}$ we have
\[
  p_N^{(2)} \;\le\; p_N^{(1)}.
\]
Equivalently, $\tau_2$ stochastically dominates $\tau_1$ in the usual sense.
\end{definition}

\begin{remark}[Static vs dynamic equivalence]
Two pipelines may be statically equivalent in the sense of Appendix~A (same normal form in
$N \ltimes E$ and same tacit semantics $J{-}K$), yet dynamically inequivalent: they can induce
different hazard profiles and hence different hallucination-time distributions.
Dynamic refinement isolates this difference at the level of trajectories, not static syntax:
$P_2$ improves on $P_1$ precisely when it never increases finite-horizon hallucination risk.
\end{remark}

\begin{remark}[Limitations]
Dynamic refinement is a partial order on stochastic implementations, not an absolute
guarantee of safety: if the baseline pipeline $P_1$ satisfies the inevitability conditions
of Theorems~B.11, B.25, and B.29, then so does any refinement $P_2$. Refinements can
delay hallucination and reduce risk over practically relevant horizons, but they cannot
eliminate hallucination altogether without violating the structural assumptions of this
appendix.
\end{remark}
% ------------------------------------------------------------------
% New subsection at the end of Appendix B
% ------------------------------------------------------------------

\subsubsection*{B.18 Structural implications of pipeline dynamics}

The move from static FRFP pipelines to the dynamic setting of this appendix
turns time and iteration into first-class objects. Beyond the inevitability
results of Theorems~B.11, B.25, and B.29, this buys several structural
notions that will be used in later appendices and case studies.

\paragraph{(1) Hazard profiles and safe horizons.}
Definition~B.31 endows each stochastic FRFP implementation with a
hallucination time $\tau$, and Definitions~B.36 and~\ref{def:hazard} refine
this to a finite-horizon risk curve $(p_N)_{N \ge 0}$ and a discrete hazard
sequence $(h_n)_{n \ge 0}$. The safe-horizon functional
$H_P(\varepsilon)$ in Definition~\ref{def:safe-horizon} abstracts how long a
given pipeline can be run, under a fixed stochastic implementation $P$ and
risk budget $\varepsilon$, before hallucination becomes likely. Thus even
when static semantics and correctness are fixed, pipeline dynamics exposes a
new design degree of freedom: the shape of the hazard profile along a run.

\paragraph{(2) Time--decomposition trade-offs.}
In the static theory of Appendices~A and~B.1--B.3, different decompositions
of a task into explicit steps are equivalent whenever they normalize to the
same morphism in $N \ltimes E$ and share the same tacit semantics $J{-}K$.
Once trajectories are considered, this equivalence is refined: more
fine-grained decompositions typically reduce the requirement demanded at
individual steps, but also increase the length of the trajectory and the
opportunity for degradation to act. Pipeline dynamics therefore makes precise
a basic design tension: reshaping a task into ``more steps'' can both reduce
per-step hazard and increase cumulative risk, and only the dynamic
quantities $(h_n)$ and $(p_N)$ decide which effect dominates in a concrete
setting.

\paragraph{(3) Dynamic refinement of implementations.}
Definition~\ref{def:dynamic-refinement} introduces a partial order on
stochastic implementations of the \emph{same} static pipeline: $P_2$
dynamically refines $P_1$ when it never increases finite-horizon risk.
This separates:
\begin{itemize}
  \item static equivalence in $N \ltimes E$ (same normal form, same tacit
        semantics), from
  \item dynamic safety properties (how quickly hallucination emerges under a
        particular scheduling, tool use pattern, or feedback coupling).
\end{itemize}
Subsequent appendices will use this distinction to talk about pipeline-level
improvements that do not change what a system computes in principle, but do
change how its risk unfolds in time.

\paragraph{(4) Session limits, restart policies, and composition.}
Because hallucination is represented as a stopping time and finite-horizon
probabilities are explicit, Appendix~B supports principled discussion of
session limits and restarts. Proposition~B.38 already shows how repeated
runs over a fixed horizon aggregate risk via $1 - (1 - p_N)^K$. In more
complex settings, this allows single-pipeline guarantees (bounds on $p_N$
or $H_P(\varepsilon)$) to be composed into system-level guarantees for
restart, parallelization, or human-in-the-loop oversight policies, much as
Appendix~A.10 composes agent-level inevitability into collective inevitability.

\paragraph{(5) Localizing and targeting failure modes.}
Finally, working with trajectories rather than only normal forms means that
hallucination can be localized to segments of a run: particular regions of
the state space, update rules, or phases of the pipeline where hazard is
concentrated. This enables targeted interventions (for example, inserting
human checks or additional tools only before entering high-hazard regions)
that would be invisible in a purely static semantics.

Taken together, these structures show that pipeline dynamics is not only a
way to restate hallucination inevitability in stochastic language. It also
provides a vocabulary---hazard, safe horizons, dynamic refinement, and
time--decomposition trade-offs---for comparing and improving FRFP pipelines
\emph{without} altering their underlying explicit semantics in $N \ltimes E$.

\pagebreak

%\subsection{Collective Epistemic Dynamics)}

\subsubsection{ Purpose and Scope}

This appendix extends  Epistemic dynamics to \emph{collective} epistemic
systems. Rather than a single agent executing a single FRFP pipeline, we
consider populations of agents who:

\begin{itemize}
  \item share an explicit artifact space $E$,
  \item each carry their own tacit state and degradation dynamics,
  \item exchange only explicit artifacts (no tacit information flow),
  \item may be coupled by explicit consensus mechanisms and institutions.
\end{itemize}

The objective is to characterize:

\begin{enumerate}
  \item how to view a population as a single large FRFP system,
  \item how stochastic collective evolution can be modeled as a Markov process
        on this system,
  \item how hallucination and other failures lift from individuals to
        populations, and
  \item which collective design changes \emph{cannot}, even in principle,
        eliminate hallucination.
\end{enumerate}

This section provides the bridge between individual pipelines (Sections A.2–A.8),
impossibility results (Appendix A.11), and institutional design (Appendix B and Appendix C).

\subsubsection{Epistemic Agents}

\begin{definition}[Epistemic agent]
An \emph{epistemic agent} is a tuple
\[
  A_i = (T_i, \,\preceq_i, \,\delta_i, \,T^e_i),
\]
where:
\begin{itemize}
  \item $T_i$ is a tacit context space equipped with a preorder
        $\preceq_i$ encoding ``at least as grounded as'';
  \item $\delta_i : T_i \to T_i$ is a (tacit) context degradation operator;
  \item $T^e_i := T_i \times \mathbb{R}_{\ge 0}$ is the extended tacit state
        space, pairing grounding $t_i \in T_i$ with tacit credence
        $c_i \in \mathbb{R}_{\ge 0}$.
\end{itemize}
\end{definition}

Intuitively, $A_i$ captures a single human (or human-like) reasoner as modeled
in Epistemic Dynamics: grounding lives in $T_i$, weakens under $\delta_i$,
and credence lives in the additional component of $T^e_i$.

\subsubsection{Epistemic Populations}

\begin{definition}[Epistemic population and category \textbf{Pop}]
An \emph{epistemic population} is a finite or countable family
$P = (A_i)_{i \in I}$ of epistemic agents indexed by a set $I$.
The category $\mathbf{Pop}$ of epistemic populations is defined by:
\begin{itemize}
  \item objects: epistemic populations $P = (A_i)_{i \in I}$;
  \item morphisms $F : P \to P'$: pairs $(f,(F_i)_{i\in I})$ where
        $f : I \to I'$ is a function on agent indices and, for each $i$,
        \[
          F_i : T_i \to T'_{f(i)}
        \]
        is monotone with respect to $\preceq_i,\preceq'_{f(i)}$ and
        commutes with degradation and credence:
        \[
          F_i \circ \delta_i = \delta'_{f(i)} \circ F_i.
        \]
\end{itemize}
\end{definition}

Morphisms in $\mathbf{Pop}$ model re-indexing or restructuring of populations:
embedding agents into larger organizations, replacing agents, or transferring
tacit states along institutional lines.

\subsubsection{Shared Explicit Reasoning Category}

Let $\mathcal{C} := N \ltimes E$ be the reasoning category of Appendix~A,
where $E$ denotes explicit artifacts and $N$ navigation moves in explicit
space. For a fixed index set $I$, we form the product category
\[
  \mathcal{C}^I,
\]
whose objects are $I$-indexed families $X = (X_i)_{i \in I}$ of explicit
states and whose morphisms are families of explicit updates. Intuitively,
$\mathcal{C}^I$ is the ``big explicit state space'' of the whole population:
it records, for each agent, which task and explicit artifacts they are
currently working with.

We assume a single shared explicit artifact space $E$; each agent $i$ is
equipped with a representational backflow map
\[
  \mathrm{RB}_i : E \to T_i,
\]
as in Appendix~A. Explicit artifacts are globally visible; tacit grounding
remains agent-local.

\subsubsection{Communication Structure}

\begin{definition}[Communication graph]
A \emph{communication structure} on population $P = (A_i)_{i \in I}$ is a
directed graph
\[
  G = (I,\;\to),
\]
where $i \to j$ indicates that explicit artifacts produced by $i$ may be
consumed by $j$ (for example, by reading a paper, listening to a briefing, or
reviewing a log). No tacit information is transmitted across $G$:
only explicit artifacts traverse edges, and representational backflow
$\mathrm{RB}_j$ determines how they affect $T_j$.
\end{definition}

\subsubsection{Collective Dynamics as a Markov System}

A \emph{collective execution} in earlier drafts was modeled as a bare
sequence of events $(i_n,e_n) \in I \times E$ indicating ``agent $i_n$
produces artifact $e_n$''. We now package such executions into a stochastic
dynamical system over the large explicit category $\mathcal{C}^I$.

\begin{definition}[Collective state]
A collective explicit state is an object $X \in \mathcal{C}^I$ together with
a population $P = (A_i)_{i\in I}$ and communication graph $G=(I,\to)$.
We write $(P,G;X)$ for such a state.
\end{definition}

\begin{definition}[Population semantics]
A \emph{population semantics functor} is a functor
\[
  \overline{J} : \mathcal{C}^I \to \prod_{i\in I} T^e_i
\]
assigning to each collective explicit state $X$ the tuple of extended tacit
states $\bigl((t_i,c_i)\bigr)_{i\in I}$ obtained when agents in $P$ consume
the artifacts in $X$ (subject to $G$) and update using the operators of
Appendix~B (degradation, confidence inflation, requirement escalation, and
feedback-coupled degradation).
\end{definition}

\begin{definition}[Markovian collective dynamics]
A \emph{Markovian collective dynamics} on $(P,G)$ is a transition kernel
\[
  M : \mathrm{Obj}(\mathcal{C}^I) \to \mathcal{P}\bigl(\mathrm{Obj}(\mathcal{C}^I)\bigr)
\]
assigning to each collective state $X$ a probability measure $M(X,\cdot)$
on next states $Y$, such that each transition can be factored as:
\begin{enumerate}
  \item select an acting agent $i_n \in I$ and explicit move in $\mathcal{C}$;
  \item update $X$ to $Y$ accordingly (producing a new artifact $e_n$ in $E$);
  \item for each $j$ with $i_n \to j$, update $(t_j,c_j)$ via the
        feedback-coupled dynamics of Appendix~B using $e_n$;
  \item apply baseline degradation for all agents.
\end{enumerate}
Equivalently, one may view $M$ as a functor
\[
  M : \mathcal{C}^I \longrightarrow \mathbf{Kl}(D),
\]
where $\mathbf{Kl}(D)$ is the Kleisli category of the discrete distribution
monad $D$, but we will only use the induced transition kernel on objects.
\end{definition}

A \emph{collective execution} is then a random path
\[
  X_0, X_1, X_2, \dots
\]
in $\mathcal{C}^I$ sampled from $M$ starting at some initial state $X_0$,
with tacit trajectories $\overline{J}(X_n) = \bigl((t_{i,n},c_{i,n})\bigr)_{i\in I}$.

\subsubsection{Collective Hallucination and Stopping Times}

We now lift the notion of hallucination to the population level.

\begin{definition}[Collective hallucination]
Given a population $P=(A_i)_{i\in I}$, a communication graph $G$, and a
collective execution $(X_n)_{n\ge 0}$ with associated tacit states
$(t_{i,n},c_{i,n})$, a \emph{collective hallucination} occurs at step $n$
if there exists an explicit artifact $e \in E$ and a nontrivial subset
$J \subseteq I$ such that:
\begin{itemize}
  \item each $j \in J$ has accepted or propagated $e$ (in the sense of
        Appendix~B: $e$ enters their explicit pipeline), and
  \item for each $j \in J$,
        \[
          t_{j,n} \prec_j \mathrm{Req}_j(e,c_{j,n}),
        \]
        i.e., the available tacit grounding is strictly below the requirement
        demanded by $e$ in context $c_{j,n}$.
\end{itemize}
\end{definition}

Let $\Omega$ be the space of all collective paths $(X_n)_{n\ge 0}$ induced
by $M$ from $X_0$, and let $P_M$ be the corresponding path measure.

\begin{definition}[Collective hallucination time]
The \emph{collective hallucination time} $\tau^{\mathrm{coll}} :
\Omega \to \mathbb{N} \cup \{\infty\}$ is the stopping time
\[
  \tau^{\mathrm{coll}}(\omega)
  := \inf\{\,n \ge 0 : \omega \text{ exhibits a collective hallucination at step } n\,\},
\]
with the convention that $\inf \varnothing = \infty$.
\end{definition}

Under the Epistemic Dynamics assumptions (finite tacit capacity, non-invertible
degradation, confidence inflation, requirement escalation, and
feedback-coupled degradation), individual agents exhibit almost-sure
hallucination along sufficiently long admissible runs (Appendix~B).
The purpose of this appendix is to show how such individual inevitability
lifts to $\tau^{\mathrm{coll}}$ at the population level.

\subsubsection{Consensus as an Explicit Functor}

Consensus mechanisms act purely in explicit space.

\begin{definition}[Consensus functor]
A \emph{consensus mechanism} on the explicit category $\mathcal{C}$ is a
functor
\[
  C : \mathcal{C}^n \to \mathcal{C}
\]
for some finite arity $n \ge 2$, or more simply
\[
  C : E^n \to E
\]
on the underlying artifact space, modeling averaging, voting, or more
complicated aggregation pipelines. Consensus functors operate entirely in
explicit space: they take explicit artifacts as inputs and produce explicit
artifacts as outputs, without access to tacit states or requirements.
\end{definition}

\subsubsection{Functorial Consensus Layer}

Given a population $P=(A_i)_{i\in I}$ and communication structure $G$,
a consensus mechanism $C$ induces an explicit functor
\[
  C_P : E^n \to E
\]
that may be interleaved with ordinary explicit computation in
$\mathcal{C}^I$. We refer to such composites as the \emph{consensus layer}.
No consensus functor has access to the population semantics
$\overline{J}$ or to individual requirement functions $\mathrm{Req}_i$.

\subsubsection{Tacit-Stability of Consensus}

\begin{definition}[Tacit-stable consensus]
A consensus functor $C$ is \emph{tacit-stable} if, for every agent $i$ and
any inputs $e_1,\dots,e_n \in E$, whenever $e_1,\dots,e_n$ are all
individually below agent $i$'s requirement,
\[
  t_i \prec_i \mathrm{Req}_i(e_k,c_i)
  \quad \forall k \in \{1,\dots,n\},
\]
then the aggregate output $C(e_1,\dots,e_n)$ is also below requirement:
\[
  t_i \prec_i \mathrm{Req}_i\bigl(C(e_1,\dots,e_n),c_i\bigr).
\]
\end{definition}

\begin{theorem}[Consensus stability]
Every consensus functor that acts purely in explicit space is tacit-stable.
\end{theorem}

\begin{proof}[Proof sketch]
Consensus functors are compositions of explicit morphisms in $\mathcal{C}$.
By Explicit--Tacit Separation (ETS), explicit morphisms cannot increase
tacit grounding: they do not alter $t_i$, only what artifacts are presented
to representational backflow. If an agent's tacit state $t_i$ is below the
requirement for each input artifact, it remains below the requirement for
any artifact obtained by explicit-only transformation of those inputs. Thus
consensus cannot manufacture sufficient grounding from insufficient inputs.
\end{proof}

Tacit-stability says that explicit consensus can neither repair nor upgrade
tacit adequacy: it preserves ``being under-grounded'' whenever that is the
case for all inputs.

\subsubsection{Limits of Automated Consensus}

\begin{theorem}[No fully automated consensus elimination]
No fully automated consensus process acting purely in explicit space can
eliminate collective hallucination across all admissible populations and
collective dynamics.
\end{theorem}

\begin{proof}[Proof sketch]
By tacit-stability, consensus cannot transform artifacts that are
under-grounded for all relevant agents into artifacts that are adequately
grounded. Under the Epistemic Dynamics degradation assumptions, individual hallucination
is inevitable along sufficiently long runs. Since consensus mechanisms have no
access to tacit state, they cannot detect or remove these failures at the
collective level; at best they can smooth, average, or denoise explicit
behavior. Any claim of global hallucination-elimination by explicit-only
consensus would contradict the Epistemic Dynamics inevitability results when applied
agent-wise.
\end{proof}

Automation may reduce variance or noise in explicit artifacts, but it cannot
introduce new tacit grounding or globally halt its degradation.

\subsubsection{Collective Inevitability}

\begin{theorem}[Collective hallucination inevitability]
Consider a population $P=(A_i)_{i\in I}$ of epistemic agents with:
\begin{itemize}
  \item finite or countable index set $I$,
  \item finite tacit capacity for each $A_i$,
  \item non-invertible degradation $\delta_i$,
  \item confidence inflation and requirement escalation as in Appendix~B.
\end{itemize}
Let $M$ be any Markovian collective dynamics on $\mathcal{C}^I$ that
respects these per-agent dynamics and produces infinitely many
surface-plausible explicit artifacts along almost every path.
Then the collective hallucination time $\tau^{\mathrm{coll}}$ is almost
surely finite:
\[
  P_M\bigl(\tau^{\mathrm{coll}} < \infty\bigr) = 1.
\]
Moreover, this property is invariant under the insertion of explicit-only
consensus functors: composing $M$ with any finite consensus layer that acts
solely in explicit space does not change the almost-sure finiteness of
$\tau^{\mathrm{coll}}$.
\end{theorem}

\begin{proof}[Proof sketch]
By Appendix~B, each agent $A_i$ exhibits almost-sure hallucination along
sufficiently long admissible runs of any FRFP pipeline that continues to
produce surface-plausible artifacts. The collective Markov dynamics $M$
induces, for each $i$, a marginal stochastic process on $T^e_i$ with the
same qualitative properties. Since the population is finite or countable,
and collective execution produces infinitely many plausible artifacts with
nonzero probability of being consumed by each agent (via $G$), the event
that at least one agent hallucinates in finite time has probability~$1$.
The definition of collective hallucination requires only that a nontrivial
subset of agents jointly accept an under-grounded artifact, which occurs
almost surely given enough time and interaction. Tacit-stability of
consensus implies that explicit aggregation cannot remove such events once
they arise; it can at best transform one hallucinated artifact into another.
\end{proof}

\subsubsection{Structural Implications of Collective Dynamics}

The population-level formalism above does more than restate
inevitability in a collective setting. It yields several structural notions
that are used in later sections.

\paragraph{(1) Equivalence classes of collective designs.}
Viewing populations as objects in $\mathbf{Pop}$ and their joint explicit
reasoning as trajectories in $\mathcal{C}^I$, we can treat different UIs,
workflows, and consensus wrappers as endofunctors on $\mathcal{C}^I$ that
preserve the population semantics $\overline{J}$. Designs that differ only
by such explicit refactorings are \emph{FRFP-equivalent}: they have the same
tacit semantics and the same hallucination-time distribution, even if their
surface appearance differs. This provides a criterion for when two collective
protocols are ``the same thing in principle''.

\paragraph{(2) Refinement order on collective systems.}
Given two Markovian collective dynamics $M_1,M_2$ on the same population,
with corresponding collective hallucination times
$\tau^{\mathrm{coll}}_1,\tau^{\mathrm{coll}}_2$, we may compare their
finite-horizon risks
\[
  p_N^{(k)} := P_{M_k}(\tau^{\mathrm{coll}}_k \le N), \quad k\in\{1,2\}.
\]
We say that $M_2$ is a \emph{dynamic refinement} of $M_1$ if
$p_N^{(2)} \le p_N^{(1)}$ for all $N$. This yields a partial order on
collective implementations: some architectures are strictly safer than
others, but purely explicit-only transformations that preserve semantics
cannot improve this order.

\paragraph{(3) Compositional risk analysis.}
Because the collective dynamics lives in $\mathcal{C}^I$ and is built from
per-agent Epistemic Dynamics, we can study how subsystem risks compose.
For example, a scientific subsystem, an engineering subsystem, and a policy
subsystem can each be modeled as a subpopulation with its own hazard profile;
the global population then inherits a collective hallucination time equal to
the minimum of the subsystem times. This provides a formal basis for
``defense in depth'' arguments.

\paragraph{(4) Limits of monitoring and metrics.}
Collective executions produce rich explicit traces: logs, dashboards,
metrics, and aggregated indicators. In this formalism, these traces live in
$\mathcal{C}^I$ or in functorial images thereof. The population semantics
$\overline{J}$ and individual requirements $\mathrm{Req}_i$, however, remain
tacit. Many distinct tacit configurations can induce the same distribution
over explicit traces. Hence correctness and hallucination risk are
\emph{non-identifiable from logs alone}: explicit-only monitoring cannot, in
general, determine how well grounded a population's beliefs really are.

\paragraph{(5) Scaling and phase transitions.}
By parametrizing collective systems by population size $|I|$, communication
structure $G$, and per-agent degradation properties, Appendix~A.10 provides a
language for asking how epistemic risk scales with size and complexity.
Nothing in the formalism guarantees that larger or more connected populations
are safer; on the contrary, the inevitability theorem applies uniformly once
the Epistemic Dynamics assumptions hold. This sets the stage for Appendix~A.12, which
studies which institutional forms can slow or bound such risks across
populations and time.

Taken together, these structures show that Collective Epistemic Dynamics is not merely a
corollary of individual hallucination: it is a genuinely new level of
description in which equivalence, refinement, observability, and scaling of
\emph{collective} epistemic systems can be stated and analyzed formally.

\pagebreak
% Appendix D — Impossibility & Limits (Phase-5)
% Schema–corollary organization with (i) admissible runs as executions in \NltimesE
% under Phase-0 constraints and (ii) explicit-identity options (syntactic vs normal-form).
% Depends on Appendix A (static axioms) and Appendices B–C (dynamics + collectives).

%\subsection{Impossibility \& Limits}

\subsubsection{Purpose and Scope}
This section collects structural \emph{no-go theorems} that delimit what can be guaranteed within FRFP-like architectures under Explicit--Tacit Separation (ETS), Human-Exclusive Grounding (HEG), and Human-Exclusive Closure (HEC). The results are invariant to model internals, training data, or intent.

The central organizing idea is that many desirable guarantees are \emph{tacit-dependent}: whether they hold can vary with human tacit state even when explicit artifacts are held fixed. Under ETS and AE (AI-Explicit Restriction), tacit-dependent guarantees cannot be enforced or decided by explicit-only mechanisms.

\subsubsection{Standing Assumptions and Notation}

Throughout this appendix we fix the FRFP kernel and assume that all interaction and
epistemic constraints from the main text are in force.

\paragraph{Foundational layers.}
On the single-agent side we assume:
\begin{itemize}
  \item the interaction-layer constraints of Section~2 and Appendix~A.2.6 (admissible traces, no AI
        access to tacit inputs, separation of content-producing and correctness-bearing moves);
  \item the epistemic ontology and explicit/tacit separation from Section~3 and Appendix~A.2--A.3
        (ambient category $C_{\mathrm{full}}$, explicit and tacit subcategories $E$ and $T$, boundary
        functor $\RB : E \to T$, and the six primitives $RI, EC, ED, RB, TE, HFD$);
  \item the explicit workflow algebra and correctness layer from Sections~4--6 and
        Appendix~A.4--A.8 (the TDG-generated explicit category, navigation category $N$, Grothendieck
        semidirect product $N \ltimes E$, semantic functor $\sem{-} : N \ltimes E \to T$, and the
        correctness quotient $\gamma : \Obj(T) \to J$).
\end{itemize}

\paragraph{Dynamic and collective layers.}
In addition, we assume the epistemic dynamics and collective semantics developed in the later parts of
Appendix~A:
\begin{itemize}
  \item the tacit degradation operator $\delta : T \to T$, requirement map
        $\Req : E \times \mathbb{R}_{\ge 0} \to T$, and the induced notion of hallucination and
        finite safe horizons from Appendix~A.9;
  \item the population and consensus layer of Appendix~A.10, which models a population of agents
        with a shared explicit category, local tacit spaces, and a communication graph, and equips
        this structure with a Markovian collective dynamics.
\end{itemize}
For a fixed explicit pipeline $P$ and an admissible implementation of $P$ in this dynamic model,
Appendix~A.9 associates to each finite horizon $N$ a hallucination probability $p_N$ and a
hazard profile $(h_n)_{n \geq 0}$, and defines a safe-horizon functional
$H_P(\varepsilon)$ (see Definition~A.\ref{def:safe-horizon}) capturing the maximal time $N$ for which
the implementation keeps hallucination risk below a tolerance $\varepsilon$. On the population side,
Appendix~A.10 provides analogous notions of collective hallucination time and collective
safe horizons.

The explicit-only impossibility and institutional non-automation results in this appendix are
understood relative to this entire stack: they assume the interaction, epistemic, algebraic,
dynamic, and collective layers just described, and ask which kinds of detectors, consensus
mechanisms, and governance operators can be realised \emph{purely explicitly} within that fixed
architecture.


\subsubsection{Admissible Runs as Interaction Protocol Filtered Executions in \texorpdfstring{N$\ltimes$E}{N ltimes E}}
\begin{definition}[Executions in N$\ltimes$E]
Fix a valid pipeline term $P$ (Epistemic Space / Algebra well-typed). An \emph{execution trace} of $P$ is a finite or infinite sequence
\[
\rho:\quad x_0 \to x_1 \to x_2 \to \cdots
\]
where each $x_k$ is a configuration in N$\ltimes$E reachable from an initial configuration for $P$ by the combined reduction relation $\to$.
\end{definition}

\begin{definition}[Interaction Protocol admissibility predicate]
Let $\mathsf{Adm}_0(\rho)$ be a predicate on execution traces encoding Interaction Protocol constraints:
\begin{enumerate}
  \item \textbf{IL (Interaction Limitation):} $\rho$ is turn-based and respects externally imposed bounds (e.g. finite-length turns).
  \item \textbf{AR (Alignment Restriction):} explicit actions occur only in response to explicit instructions present in the trace.
  \item \textbf{MD (Mode Discipline):} no step in $\rho$ applies tacit operators inside the explicit phase; tacit transitions occur only after the boundary.
  \item \textbf{NTER/CSC:} no step reconstructs tacit evaluation from explicit traces or treats explicit computation as a correctness oracle.
\end{enumerate}
We treat $\mathsf{Adm}_0$ as an axiomatically specified filter on traces.
\end{definition}

\begin{definition}[Admissible runs]
An \emph{admissible run} of a pipeline $P$ is an execution trace $\rho$ in N$\ltimes$E such that $\mathsf{Adm}_0(\rho)$ holds.
We write $\mathrm{Runs}(P)$ for the set of admissible runs of $P$.
\end{definition}

\begin{remark} (Admissible runs and explicit-only results).
This definition pins all explicit-only impossibility and non-automation statements to a common
operational semantics: admissible runs are exactly those executions in \(N \ltimes E\) that
additionally satisfy the interaction constraints (IL, AR, MD, NTER/CSC). All explicit-only
no-go theorems in Appendix~A.11 are understood to quantify over this class of runs.

\end{remark}

\subsubsection{Two Notions of ``Same Explicit Artifact''}
No-go statements often quantify over pairs of runs that produce the ``same explicit artifact.'' There are two natural notions.

\begin{definition}[Syntactic explicit identity]
Two runs are \emph{syntactically explicit-identical} if they produce equal explicit artifacts as elements of the explicit object $E$ (i.e. equality in the representation domain).
\end{definition}

\begin{definition}[Semantic explicit identity via normal forms]
Assume the explicit reduction system $\to_{\Ecat}$ is terminating and locally confluent (Appendix~A axioms), so each explicit artifact has a unique normal form $\mathrm{nf}(e)$.
Two runs are \emph{normal-form explicit-identical} if their produced artifacts $e,e'$ satisfy
\[
\mathrm{nf}(e)=\mathrm{nf}(e').
\]
\end{definition}

\begin{remark}
Syntactic identity is appropriate when artifacts are literal strings/terms. Normal-form identity is appropriate when explicit equivalence is taken modulo rewriting (e.g. definitional equalities, normalization, canonical forms). All no-go results below hold under either identity notion; when needed we indicate which is assumed.
\end{remark}
\begin{remark}[Admissible runs and stochastic trajectories]
In the dynamic semantics of Appendix~B, the same pipeline $P$ gives rise
to a measurable space of trajectories $\Omega_{\mathrm{allowed}}$,
together with a family of probability measures capturing different
stochastic implementations. We implicitly identify each admissible run
$\rho$ of $P$ in the combined system $\mathsf{N} \ltimes \mathsf{E}$
with its corresponding trajectory
$\omega \in \Omega_{\mathrm{allowed}}$. In particular, the hallucination time $\tau(\omega)$, the finite-horizon probabilities $(p_N)$,
the hazard sequence $(h_n)$, and the safe-horizon function $H_P(\varepsilon)$
of Appendix~B are all computed over this common set of admissible runs.

\end{remark}

\subsubsection{Tacit Dependence via \texorpdfstring{$\gamma(\sem{P})$}{gamma(sem(P))}}
\begin{definition}[Tacit-dependent predicates]
Let $P$ be a fixed valid pipeline and let $\mathrm{Runs}(P)$ be its admissible runs.
A predicate $\mathcal{G}$ about an explicit artifact is \emph{tacit-dependent} if there exist two admissible runs $\rho_1,\rho_2\in\mathrm{Runs}(P)$ such that:
\begin{itemize}
  \item the runs produce the same explicit artifact (either syntactically or up to normal form, per Section~D.3), and
  \item the observable correctness outcomes differ:
  \[
    \gamma(\sem{P})\ \text{differs between}\ \rho_1\ \text{and}\ \rho_2.
  \]
\end{itemize}
\end{definition}

\begin{remark}
This definition formalizes ``depends on tacit context'' using only Epistemic Algebra semantics: $\sem{-}$ captures tacit transformation induced by the pipeline run, and $\gamma$ extracts the human-observable correctness class.
\end{remark}
\begin{remark}[Examples of tacit-dependent predicates]
\label{rem:tacit-dependent-examples}
The definition above applies not only to pointwise correctness
judgements but also to quantitative long-horizon properties. Typical
examples of tacit-dependent predicates $G$ include:
\begin{itemize}
  \item the event that a run hallucinates within a fixed horizon $N$,
        i.e.\ $G(\omega)$ is the predicate $[\tau(\omega) \le N]$;
  \item the statement that an implementation has $\varepsilon$-safe horizon at least $N$, i.e.
  $G$ is the predicate $[H_P(\varepsilon) \ge N]$ for some fixed $\varepsilon \in (0,1)$;

  \item in the collective setting of Appendix~A.10, the statement that the
        collective hallucination time exceeds a fixed horizon, i.e.\
        $G$ is the predicate $[\tau_{\mathrm{coll}} > N]$ for a given
        population and consensus dynamics.
\end{itemize}
In each case, there can exist admissible runs with the same explicit
artifacts (in the sense of Definition~B.\,66) but different values of
$\tau$, $\HP(\varepsilon)$ or $\tau_{\mathrm{coll}}$, so $G$ is
tacit-dependent in the sense of Definition~B.\,66.
\end{remark}

\subsubsection{Schema Theorem: Tacit-Dependent Guarantees Are Not Explicit-Implementable}
\begin{definition}[Explicit-only mechanism]
An \emph{explicit-only mechanism} is any procedure definable using explicit morphisms in $\Ecat$ (and finite control) whose inputs and outputs live entirely in explicit space. Concretely, one may model an explicit-only mechanism as a total function $M:E^k\to E^\ell$, or as a family of such functions closed under composition.
\end{definition}

\begin{theorem}[Tacit-dependent guarantees are not explicit-implementable]\label{thm:D-schema}
Assume ETS and AE. Let $\mathcal{G}$ be a tacit-dependent predicate about explicit artifacts (Section~D.4), where ``same explicit artifact'' is interpreted either syntactically or by normal form (Section~D.3). Then no explicit-only mechanism can enforce $\mathcal{G}$ for all admissible runs.
\end{theorem}

\begin{proof}[Proof sketch]
By tacit dependence, there exist admissible runs producing the same explicit artifact (under the chosen identity notion) but yielding different outcomes $\gamma(\sem{P})$. Any explicit-only mechanism sees only the explicit artifact(s) and therefore cannot act differently on these runs. Hence it cannot enforce a guarantee whose truth varies across admissible runs while explicit data are held fixed.
\end{proof}

\subsubsection{Corollaries (No-Go Family)}
Each of the following is an instantiation of Theorem~\ref{thm:D-schema}.

\subsubsubsection{No Perfect Explicit-Only Hallucination Detector}
\begin{definition}[Explicit-only monitor]
An \emph{explicit-only monitor} is a function $M:E\to\{\mathrm{ok},\mathrm{flag}\}$.
\end{definition}

\begin{definition}[Perfect detector]
$M$ is \emph{perfect} if, for every admissible run producing explicit artifact $e$ with tacit context $t$ at evaluation time,
\[
M(e)=\mathrm{ok}\Rightarrow t\not\prec\mathsf{Req}(e),\qquad
M(e)=\mathrm{flag}\Rightarrow t\prec\mathsf{Req}(e).
\]
\end{definition}

\begin{corollary}[No perfect explicit-only detector]\label{cor:D-no-perfect-detector}
Assume ETS and AE, and assume requirements are nontrivial and contexts vary across admissible runs. Then no explicit-only monitor is a perfect hallucination detector.
\end{corollary}

\begin{proof}[Proof sketch]
The predicate ``$e$ is a hallucination'' is tacit-dependent: the same explicit artifact (syntactically or up to normal form) can be within grounding in one admissible run and beyond grounding in another. Apply Theorem~\ref{thm:D-schema}.
\end{proof}

\begin{corollary}[No elimination of hallucination under degradation]
\label{cor:no-elimination-under-degradation}
Assume the Epistemic Dynamics conditions that yield almost-sure
hallucination inevitability: in particular, a finite-capacity tacit
state space, a monotone non-invertible degradation operator
$\delta \colon T \to T$, and recurring nontrivial requirements as in
Theorem~B.\,29. Let $\tau$ denote the hallucination time random variable
for a fixed pipeline $P$ and any admissible stochastic implementation
of $P$ whose trajectories lie in $\Omega_{\mathrm{allowed}}$.
Then there is no explicit-only policy which guarantees
\[
  \tau(\omega) \;=\; \infty
\]
for all admissible runs $\omega \in \Omega_{\mathrm{allowed}}$, nor even
$\mathbb{P}(\tau = \infty) = 1$ for all such implementations.
\end{corollary}

\begin{proof}
Under the assumptions of Appendix~B, Theorem~B.\,29 states that every
admissible stochastic implementation of $P$ has hallucination time
$\tau$ finite almost surely:
\[
  \mathbb{P}(\tau < \infty) \;=\; 1
\]
whenever its support is contained in $\Omega_{\mathrm{allowed}}$.
Consider the predicate
\[
  G(\omega) \;\equiv\; [\,\tau(\omega) = \infty\,]
\]
on admissible runs. By construction of the dynamic semantics, there
exist admissible runs with the same explicit artifacts but different
tacit evolutions along which $\tau$ takes different values; in
particular, $G$ is tacit-dependent in the sense of
Definition~B.\,66. If there were an explicit-only policy that guaranteed
$\tau = \infty$ almost surely for all admissible implementations, that
policy would enforce the tacit-dependent predicate $G$ by purely
explicit means, contradicting Theorem~B.\,69. Hence no explicit-only
policy can eliminate hallucination under the Epistemic Dynamics degradation
assumptions.
\end{proof}
\begin{corollary}[No explicit-only infinite safe horizon]\label{cor:no-infinite-safe-horizon}
Assume the Epistemic Dynamics conditions of Appendix~B that yield almost-sure hallucination
inevitability, and let $H_P(\varepsilon)$ be the $\varepsilon$-safe horizon of a
stochastic FRFP pipeline $P$ as in Definition~B.41. Then for any fixed
$\varepsilon \in (0,1)$ there is no explicit-only policy that guarantees
$H_P(\varepsilon) = \infty$ for all admissible stochastic implementations of $P$.
\end{corollary}

\begin{proof}
Under the Epistemic Dynamics assumptions, Theorem~B.29 implies that for any admissible
implementation we have $P(\tau < \infty) = 1$, so for every $\varepsilon \in (0,1)$
the set $\{N \in \mathbb{N} : p_N \le \varepsilon\}$ is bounded almost surely, i.e.\
$H_P(\varepsilon) < \infty$ along admissible runs. The predicate
$[H_P(\varepsilon) = \infty]$ is tacit-dependent (Remark~B.69), so an explicit-only
policy enforcing it for all admissible implementations would contradict
Theorem~B.69.
\end{proof}


\subsubsubsection{No Fully Automated Correctness Oracle}
\begin{definition}[Correctness oracle]
A \emph{correctness oracle} is a function $\Omega:E\to J$ (e.g. $J=\{\mathrm{accept},\mathrm{reject},\mathrm{revise}\}$) such that for every admissible run of a pipeline $P$ producing artifact $e$,
\[
\Omega(e)\;=\;\gamma(\sem{P}).
\]
\end{definition}

\begin{corollary}[No explicit correctness oracle]\label{cor:D-no-oracle}
Assume HEG and HEC. No explicit-only oracle $\Omega$ can reproduce $\gamma(\sem{P})$ for all admissible runs without violating the FRFP axioms.
\end{corollary}

\begin{proof}[Proof sketch]
Correctness outcomes $\gamma(\sem{P})$ are tacit-dependent by HEG/HEC: they are determined by tacit evaluation and closure, which can vary across admissible runs even when the explicit artifact is held fixed (under the chosen explicit identity). Apply Theorem~\ref{thm:D-schema}.
\end{proof}

\subsubsubsection{No Grounding Gain from Explicit Consensus}
Assume the functorial consensus model of Appendix~A.10.

\begin{corollary}[Consensus cannot create grounding]\label{cor:D-no-consensus-grounding}
Let $\mathcal{C}:\Ecat^n\to\Ecat$ be any consensus functor acting purely in explicit space. For any agent $i$ with tacit state $t_i$, if each input artifact $e_k$ exceeds grounding for $i$, then $\mathcal{C}(e_1,\dots,e_n)$ also exceeds grounding for $i$.
\end{corollary}

\begin{proof}[Proof sketch]
The predicate ``within grounding for agent $i$'' is tacit-dependent. Since $\mathcal{C}$ is explicit-only, it cannot enforce the guarantee that aggregation restores grounding. Apply Theorem~\ref{thm:D-schema}.
\end{proof}

\subsubsubsection{No Fully Automated Consensus Stability Guarantee}
\begin{corollary}[No automated consensus guarantee]\label{cor:D-no-auto-consensus}
Under ETS and HEG/HEC, no fully automated consensus pipeline can guarantee that its outputs are free of hallucination across all admissible runs and populations.
\end{corollary}
In the Markovian population semantics of Appendix~A.10, any choice of
consensus functors and explicit-only coordination rules determines a
collective dynamics $M$ on the product category together with a
collective hallucination time $\tau_{\mathrm{coll}}$. Under the Epistemic Dynamics
assumptions lifted to the population level, Theorem~B.\,58 shows that
$\tau_{\mathrm{coll}}$ is almost surely finite and, moreover, that its
almost-sure finiteness is invariant under insertion of additional
explicit-only consensus layers. A purported fully automated consensus
pipeline guaranteeing hallucination-free outputs for all admissible
populations would therefore enforce the tacit-dependent predicate
$[\tau_{\mathrm{coll}} = \infty]$ using only explicit mechanisms,
contradicting Theorem~B.\,69 together with Theorem~B.\,58.

\begin{proof}[Proof sketch]
A guarantee of hallucination-freedom is tacit-dependent and would yield a correctness oracle or perfect detector for consensus outputs. Apply Theorem~\ref{thm:D-schema} and Corollary~\ref{cor:D-no-perfect-detector}.
\end{proof}

\subsubsubsection{Catching Does Not Imply Eliminating}
\begin{corollary}[Detection does not imply elimination]\label{cor:D-detect-vs-elim}
Even if a subclass of failures is detectable by explicit invariants (e.g. syntactic checks, type checks, cross-consistency tests), this does not imply that hallucination (defined by tacit adequacy) is eliminable in general.
\end{corollary}

\begin{proof}[Proof sketch]
Explicit detectability concerns explicit predicates on $E$. Hallucination is tacit-dependent (Section~D.6.1). Eliminating hallucination would require enforcing a tacit-dependent guarantee by explicit-only means, contradicting Theorem~\ref{thm:D-schema}.
\end{proof}

\subsubsection{Interpretation and Design Consequences}
The no-go theorems do not claim that automation is useless; they claim that automation cannot replace tacit grounding without changing axioms.

Mitigations must therefore be expressed as:
\begin{itemize}
  \item restricting what is demanded of explicit space (reducing $\mathsf{Req}$),
  \item injecting fresh tacit evaluation (additional humans, replication),
  \item designing institutions that slow confidence inflation and feedback (Appendix~E).
\end{itemize}
Appendix~B refines these impossibility results with quantitative
structure. Each stochastic implementation of a pipeline induces a
hallucination time $\tau$, a finite-horizon risk curve
$(p_N)_{N \in \mathbb{N}}$, a hazard sequence $(h_n)_{n \in \mathbb{N}}$
and, when defined, a safe-horizon functional $\HP(\varepsilon)$. The
dynamic refinement order of Definition~B.\,43 then compares pipeline
implementations by their finite-horizon risks. Within the structural
limits above, design work can still aim to shape hazard profiles so that
risk is concentrated outside the operational horizon of interest and to
select implementations that dynamically refine a baseline, even though
eventual hallucination and the tacit-dependence of correctness remain
unavoidable.

Appendix~A.10 extends the picture to populations and collective dynamics.
A choice of consensus operators and coordination policies determines a
Markov process $M$ on the space of explicit population states and a
collective hallucination time $\tau_{\mathrm{coll}}$. The collective
inevitability and invariance results of Appendix~10  show
that purely explicit-only changes to consensus and logging cannot alter
the inevitability class of $\tau_{\mathrm{coll}}$. What they can do is
change the collective hazard profile and the distribution of
$\tau_{\mathrm{coll}}$, for example by producing a dynamics $M_2$ that
is a population-level dynamic refinement of some baseline $M_1$. In
this sense, long-horizon design remains meaningful: it can bend but not
break the structural bounds established in this appendix.

% End Appendix D

\pagebreak
% ============================================================
% Appendix E — Institutions and Design Constraints (Phase-6)
% ============================================================

%\subsection{Institutions and Design Constraints}

\subsubsection{Purpose and Scope}

Appendix A till now establishes that correctness in FRFP is structurally tacit,
non-automatable, and dynamically fragile under repeated interaction.
Appendix~A.11 formalizes no-go results showing that no single agent or pipeline
can eliminate hallucination, drift, or semantic collapse.

This appendix studies the \emph{institutional layer} required to stabilize
knowledge production over long time horizons.  Institutions are treated not
as social artifacts but as \emph{formal epistemic mechanisms} that operate on
collections of runs, agents, and tacit judgments.

The goal is to characterize which institutional constraints are
\emph{necessary}, which are \emph{insufficient}, and which are
\emph{structurally unavoidable} for scientific correctness.

\subsubsection{Epistemic Populations}
\label{sec:appendix-E1}

\begin{definition}[Epistemic agent, institutional view]\label{def:B82}
An \emph{epistemic agent} is an epistemic agent in the sense of Definition~B.46 of Appendix~C: a tuple
\[
  A_i = (T_i, \preceq_i, \delta_i, T^e_i),
\]
where
\begin{itemize}
  \item $T_i$ is a tacit context space equipped with a preorder $\preceq_i$ encoding ``at least as grounded as'';
  \item $\delta_i : T_i \to T_i$ is a (tacit) context degradation operator;
  \item $T^e_i := T_i \times \mathbb{R}_{\ge 0}$ is the extended tacit state space pairing a grounding $t_i \in T_i$ with a tacit credence $c_i \in \mathbb{R}_{\ge 0}$.
\end{itemize}
For each agent $A_i$ we write
\[
  \gamma_i : T_i \to J
\]
for the induced correctness quotient obtained from the Epistemic Algebra semantic structure (Appendix~A): $\gamma_i(t)$ is the observable correctness judgment when $A_i$ evaluates the tacit state $t$.

We will refer to
\[
  (T_i, \preceq_i, \delta_i, T^e_i, \gamma_i)
\]
as the \emph{institutional signature} of the agent $A_i$.
\end{definition}

\begin{definition}[Epistemic population and category \texorpdfstring{$\mathrm{Pop}$}{Pop}]\label{def:B83}
An \emph{epistemic population} is an object
\[
  P = (A_i)_{i \in I}
\]
of the category $\mathrm{Pop}$ defined in Definition~B.47 of Appendix~C: a finite or countable family of epistemic agents indexed by a set $I$.

Morphisms $F : P \to P'$ in $\mathrm{Pop}$ are pairs $(f, (F_i)_{i \in I})$, where $f : I \to I'$ is a function on indices and each
\[
  F_i : T_i \to T'_{f(i)}
\]
is monotone with respect to the tacit preorders and commutes with degradation and extended tacit state. In this appendix we fix an epistemic population
\[
  P = (A_i)_{i \in I}
\]
unless explicitly stated otherwise. Agents in $P$ may differ in tacit capacity, degradation dynamics, credence behavior, or evaluative norms; no homogeneity is assumed.
\end{definition}


\subsubsection{Runs, Replication, and Population Semantics}
\label{sec:appendix-E2}

Recall from Appendix~A that a well-typed FRFP pipeline $P$ induces a semantic value $\llbracket P \rrbracket$ in the tacit space and an observable correctness outcome $\gamma(\llbracket P \rrbracket)$. At the dynamic level, Epistemic Dynamics refines this to admissible runs in $N \ltimes E$, and Appendix~~A.11 identifies admissible runs with execution traces satisfying Interaction Protocol constraints. In the collective setting of Appendix~C, such runs are further coupled to population states via the population semantics functor.

\begin{definition}[Admissible run]\label{def:B84}
Fix a Epistemic Structure well-typed pipeline $P$ and an epistemic agent $A_i$ in the sense of Definition~\ref{def:B82}. An \emph{admissible run} of $P$ by $A_i$ is
\begin{itemize}
  \item an admissible execution trace $\rho \in \mathrm{Runs}(P) \subseteq N \ltimes E$ in the sense of Definition~B.61 (Appendix~D), and
  \item an initial extended tacit state $(t_{i,0}, c_{i,0}) \in T^e_i$ for $A_i$,
\end{itemize}
such that the evolution of $(t_{i,n}, c_{i,n})_{n \ge 0}$ along $\rho$ is governed by the Epistemic Dynamics update rules (degradation, confidence inflation, requirement escalation, and feedback-coupled degradation) and the Interaction Protocol constraints are respected at each step.
\end{definition}

\begin{definition}[Replication family]\label{def:B85}
Let $P$ be a fixed FRFP pipeline and let
\[
  P = (A_i)_{i \in I}
\]
be an epistemic population. A \emph{replication family} for $P$ in $P$ is a family
\[
  R(P) = \{\rho_i : i \in I_{\mathrm{rep}}\}
\]
indexed by a finite or countable subset $I_{\mathrm{rep}} \subseteq I$, where for each $i \in I_{\mathrm{rep}}$,
\begin{itemize}
  \item $\rho_i$ is an admissible run of $P$ by agent $A_i$ (Definition~\ref{def:B84}), and
  \item we write $(t_{i,*}, c_{i,*}) \in T^e_i$ for a chosen evaluation-time extended tacit state along $\rho_i$ (for example, the terminal state when the run halts or the state at a prescribed stopping time).
\end{itemize}
Intuitively, $R(P)$ records ``one execution of $P$ per agent'' under potentially different tacit initial conditions and stochastic realizations, but all subject to the same explicit pipeline and admissibility constraints.
\end{definition}

\begin{definition}[Population semantic profile]\label{def:B86}
Let $R(P) = \{\rho_i : i \in I_{\mathrm{rep}}\}$ be a replication family for $P$ in $P$, and let
\[
  (t_{i,*}, c_{i,*}) \in T^e_i
\]
be the evaluation-time extended tacit state associated to $\rho_i$ as above. The \emph{population semantic profile} of $P$ relative to $R(P)$ is the multiset
\[
  \mathcal{P}(P; R(P)) \;:=\; \bigl\{ \gamma_i(t_{i,*}) : i \in I_{\mathrm{rep}} \bigr\} \;\subseteq\; J.
\]
This multiset---not any single run---will be the basic carrier of institutional correctness at the population level.
\end{definition}


\subsubsection{Institutional Aggregation Operators}
\label{sec:appendix-E3}

We now formalize how institutions act on population-level correctness data rather than on individual runs or artifacts.

\begin{definition}[Institutional aggregator]\label{def:B87}
Let $J$ be the space of observable correctness judgments (e.g.\ $J = \{\mathsf{accept}, \mathsf{reject}, \mathsf{revise}\}$). Write $M(J)$ for the collection of finite or countable multisets of elements of $J$. An \emph{institutional aggregator} is a function
\[
  I : M(J) \to J \cup \{\bot\},
\]
where $\bot$ denotes suspended judgment. Typical examples include majority vote, unanimity, weighted consensus, veto rules, or deliberative deferral.

Given a fixed pipeline $P$ and epistemic population $P$, we may write $I_P$ for the particular aggregator that an institution applies to population semantic profiles $\mathcal{P}(P; R(P))$ arising from replication families $R(P)$ of $P$.
\end{definition}

\begin{definition}[Institutional outcome]\label{def:B88}
Let $P$ be a pipeline, $P$ an epistemic population, $R(P)$ a replication family, and $I$ an institutional aggregator. The \emph{institutional outcome} of $P$ relative to $R(P)$ under $I$ is
\[
  I_P\bigl(\mathcal{P}(P; R(P))\bigr) \;\in\; J \cup \{\bot\}.
\]
By construction, $I_P$ operates only on tacitly produced correctness judgments collected in the multiset $\mathcal{P}(P; R(P))$. No raw explicit artifact or execution trace is directly admissible as an input to $I_P$.
\end{definition}

\begin{definition}[Explicit-only institution]\label{def:B88prime}
An institutional aggregator family $(I_P)_P$ is \emph{explicit-only} if, for each pipeline $P$ and replication family $R(P)$, there exists an explicit-only mechanism
\[
  M_P : E^k \to J \cup \{\bot\}
\]
in the sense of Definition~B.70 of Appendix~D such that
\[
  I_P\bigl(\mathcal{P}(P; R(P))\bigr)
  \;=\;
  M_P(e_1, \dots, e_k),
\]
where $e_1, \dots, e_k \in E$ are the explicit artifacts produced along the admissible runs in $R(P)$, and $M_P$ has no access to any tacit states or tacit correctness outcomes. We will see in Theorem~B.93 that such explicit-only institutions cannot, in general, preserve correctness.
\end{definition}


\subsubsection*{E.4 Repeatability and Reproducibility}

\begin{definition}[Repeatability]
A pipeline $P$ is repeatable relative to $\mathbb{P}$ if all but finitely many
elements of $\Sem_{\mathbb{P}}(P)$ are equal.
\end{definition}

\begin{definition}[Reproducibility]
A pipeline $P$ is reproducible if for any sufficiently large epistemic
population $\mathbb{P}'$ compatible with $\mathbb{P}$, the institutional
outcome is invariant.
\end{definition}

Repeatability is a property of runs; reproducibility is a property of
institutions.

\subsubsection*{E.5 Institutional Drift and Decay}

Even if individual agents degrade (Appendix~B), institutions can buffer
against drift by averaging, vetoing, or deferring.

\begin{definition}[Institutional stability]
An institution $\mathcal{I}$ is stable if for any pipeline $P$ whose
population semantic profile converges, the institutional outcome converges.
\end{definition}

\begin{remark}
Stability does \emph{not} imply correctness—only resistance to oscillation.
\end{remark}

\subsubsection{Limits of Institutional Automation (Revised Statement)}
\label{sec:appendix-E6}

\begin{theorem}[Institutional non-automation]\label{thm:B93}
Assume ETS and AE, together with the stochastic and population semantics of Appendices~B and~C. Let $(I_P)_P$ be an institutional layer such that for every pipeline $P$ and every replication family $R(P)$ in any epistemic population, the institutional outcome $I_P(\mathcal{P}(P; R(P)))$ is computed by an explicit-only institution in the sense of Definition~\ref{def:B88prime}: there exists an explicit-only mechanism
\[
  M_P : E^k \to J \cup \{\bot\}
\]
whose inputs are only the explicit artifacts produced along the admissible runs in $R(P)$ and such that
\[
  I_P\bigl(\mathcal{P}(P; R(P))\bigr) \;=\; M_P(e_1, \dots, e_k).
\]
If, in addition, the institution is correctness-preserving in the sense that for every admissible pipeline $P$ and replication family $R(P)$,
\[
  I_P\bigl(\mathcal{P}(P; R(P))\bigr)
\]
agrees with the tacit correctness multiset $\{\gamma_i(t_{i,*}) : i \in I_{\mathrm{rep}}\}$ whenever this multiset is pointwise unanimous, then no such institution can exist.

Equivalently: there is no correctness-preserving institutional layer that depends solely on explicit artifacts or explicit-only summaries of those artifacts across all admissible pipelines and populations.
\end{theorem}

\begin{proof}[Proof sketch]
By Corollary~B.78 of Appendix~D, under HEG and HEC there is no explicit-only correctness oracle
\[
  \Omega : E \to J
\]
that reproduces $\gamma(\llbracket P \rrbracket)$ for all admissible runs of all pipelines. Any explicit-only institution satisfying the hypothesis would induce such an oracle by restricting to replication families of size one and composing
\[
  e \;\mapsto\; \{e\} \;\mapsto\; I_P\bigl(\mathcal{P}(P; R(P))\bigr) \;:=\; \Omega(e).
\]
Since correctness outcomes are tacit-dependent, this contradicts the no-go theorem for explicit-only mechanisms (Theorem~B.71). Therefore, no purely explicit-only institutional aggregator can preserve correctness across all admissible runs and populations.
\end{proof}


\subsubsection*{E.7 Peer Review as a Formal Operator}

\begin{definition}[Peer review operator]
Peer review is an institutional aggregator $\mathcal{I}_{\mathrm{PR}}$ whose
domain consists of independently generated tacit judgments by distinct
agents, possibly with structured disagreement.
\end{definition}

Peer review does not eliminate hallucination (Appendix~B) but reduces the
probability that hallucinated explicit artifacts pass institutional filters.

\subsubsection*{E.8 Consensus Is Not a Fixpoint}

\begin{theorem}[Consensus instability]
Even stable consensus does not imply semantic correctness.
\end{theorem}

\begin{proof}[Proof sketch]
Consensus is an aggregation over $\gamma(\sem{r})$ values. If the population
shares a correlated degradation mechanism, consensus may converge on an
incorrect judgment. This follows from the hallucination inevitability and
confidence inflation results of Appendix~B.
\end{proof}

\subsubsection*{E.9 Design Constraints for AI Systems}

Any AI system intended for epistemic use must therefore satisfy:

\begin{itemize}
\item explicit confinement (Epistemic Space);
\item exposure of outputs to multiple independent tacit evaluators;
\item support for suspension of judgment ($\bot$);
\item resistance to confidence inflation feedback at the institutional level.
\end{itemize}

These are not ethical desiderata but \emph{formal necessities}.

\subsubsection*{E.10 Summary}

Institutions are not optional add-ons but mathematically required structures
for long-horizon epistemic stability. They do not eliminate hallucination,
drift, or error; they bound their impact across populations and time.

FRFP therefore implies a strict separation of roles:

\begin{center}
\emph{Pipelines compute.  
Humans judge.  
Institutions stabilize.}
\end{center}

Any attempt to collapse these layers violates one of the no-go results proven
in earlier appendices.

% End Appendix E

\pagebreak
% ============================================================
% Appendix F — Governance and Law (Phase-7)
% ============================================================

%\subsection{Governance, Law, and Epistemic Oversight}

\subsubsection{Position in the Architecture}

Previously we establish:

\begin{itemize}
\item individual pipeline correctness  ,
\item dynamic degradation and hallucination inevitability  ,
\item collective epistemic dynamics and population effects  ,
\item formal impossibility results  ,
\item institutional stabilization mechanisms .
\end{itemize}

Next we study \emph{governance and law} as higher-order structures that
operate \emph{on institutions themselves}, rather than on pipelines or agents.

\subsubsection{Governance as a Meta-Institution}
\label{sec:appendix-F1}

Let $J$ and $M(J)$ be as in Appendix~E, and let an \emph{institutional aggregator family} be a collection
\[
  (I_P)_P, \qquad I_P : M(J) \to J \cup \{\bot\},
\]
assigning to each FRFP pipeline $P$ an institutional aggregator $I_P$ in the sense of Definition~B.87.

\begin{definition}[Category of institutions]\label{def:F1-Inst}
We define the category $\Inst$ as follows.
\begin{itemize}
  \item Objects are institutional aggregator families $(I_P)_P$.
  \item A morphism $\Phi : (I_P)_P \to (I'_P)_P$ is a natural transformation whose components
  \[
    \Phi_P : M(J) \to M(J)
  \]
  are monotone with respect to refinement of multisets and preserve unanimity: whenever all elements of $S \in M(J)$ agree, so do the elements of $\Phi_P(S)$.
\end{itemize}
Intuitively, a morphism in $\Inst$ is a structure-preserving transformation between institutional layers that does not introduce disagreement where there was none.
\end{definition}

\begin{definition}[Governance operator]\label{def:F1-governance}
A \emph{governance operator} is an endofunctor
\[
  G : \Inst \to \Inst
\]
together with a choice of observable
\[
  \Omega : \mathsf{Pop} \to O
\]
from the category of epistemic populations (Appendix~C) to some space of observables $O$, such that for each time step $t$,
\[
  (I^{(t+1)}_P)_P = G(I^{(t)}_P)_P
\]
is allowed to depend only on the previous institutional layer $(I^{(t)}_P)_P$ and on the history of observables $(\Omega(P_s))_{s \le t}$ generated by the collective dynamics of Appendix~C.
\end{definition}


\subsubsection{Link to Collective Epistemic Dynamics}

Appendix~C models collective epistemic evolution as a dynamical system over
populations:
\[
(\mathbb{P}_t, \Sem_{\mathbb{P}_t})_{t \ge 0}.
\]

Governance operators are defined on this level.

\begin{definition}[Governance observable]
A governance observable is any functional
\[
\Omega : (\mathbb{P}, \Sem_{\mathbb{P}}) \longrightarrow \mathbb{R}
\]
that measures population-level epistemic properties (e.g.\ disagreement
rates, confidence inflation, drift velocity).
\end{definition}

\begin{definition}[Governed evolution]
A governed epistemic system evolves by:
\[
\mathcal{I}_{t+1} = \mathcal{G}(\mathcal{I}_t, \Omega_t),
\qquad
(\mathbb{P}_{t+1}, \Sem_{\mathbb{P}_{t+1}}) = \Phi(\mathcal{I}_{t+1}),
\]
where $\Phi$ is the collective dynamics functor from Appendix~C.
\end{definition}

Thus governance is \emph{feedback on institutions}, not on beliefs.

\subsubsection{Law as a Constraint on Governance}

Governance itself is subject to failure modes: capture, overfitting,
short-term optimization, or erosion of tacit authority.

\begin{definition}[Legal constraint]
A legal constraint is a predicate
\[
\mathcal{L}(\mathcal{I}) \in \{\text{admissible},\text{inadmissible}\}
\]
on institutional forms.
\end{definition}

\begin{definition}[Lawful governance]
A governance operator $\mathcal{G}$ is lawful if:
\[
\mathcal{L}(\mathcal{I}) = \text{admissible}
\;\Rightarrow\;
\mathcal{L}(\mathcal{G}(\mathcal{I})) = \text{admissible}.
\]
\end{definition}

Law therefore restricts the allowable transformations of institutions, not
the content of scientific claims.

\subsubsection{No Fully Automated Governance}

We now state the core limit result.
\begin{definition}[Explicit-only governance]\label{def:F-explicit-only-governance}
A governance operator
\[
  G : \Inst \to \Inst
\]
is \emph{explicit-only} if, for every time $t$ and institutional layer $(I^{(t)}_P)_P$, there exists a family of explicit-only mechanisms
\[
  G^{\mathrm{exp}}_{P,t} : E^{k_{P,t}} \to J \cup \{\bot\}
\]
in the sense of Appendix~D such that the updated institutional aggregators satisfy
\[
  I^{(t+1)}_P(S)
  \;=\;
  G^{\mathrm{exp}}_{P,t}\bigl(e^{(t)}_1, \dots, e^{(t)}_{k_{P,t}}\bigr)
\]
for every multiset $S \in M(J)$, where $e^{(t)}_1, \dots, e^{(t)}_{k_{P,t}} \in E$ are the explicit artifacts produced by the admissible runs of $P$ up to time $t$. In particular, $G$ has no access to any tacit states or tacit correctness outcomes, only to explicit execution traces and their explicit summaries.
\end{definition}
\begin{definition}[Explicit-only governance]\label{def:F-explicit-only-governance}
A governance operator
\[
  G : \Inst \to \Inst
\]
is \emph{explicit-only} if, for every time $t$ and institutional layer $(I^{(t)}_P)_P$, there exists a family of explicit-only mechanisms
\[
  G^{\mathrm{exp}}_{P,t} : E^{k_{P,t}} \to \bigl(M(J) \to J \cup \{\bot\}\bigr)
\]
in the sense of Appendix~D such that the updated institutional aggregators satisfy
\[
  I^{(t+1)}_P(S)
  \;=\;
  G^{\mathrm{exp}}_{P,t}\bigl(e^{(t)}_1, \dots, e^{(t)}_{k_{P,t}}\bigr)(S)
\]
for every multiset $S \in M(J)$, where $e^{(t)}_1, \dots, e^{(t)}_{k_{P,t}} \in E$ are the explicit artifacts produced by the admissible runs of $P$ up to time $t$. In particular, $G$ has no access to any tacit states or tacit correctness outcomes, only to explicit execution traces and their explicit summaries.
\end{definition}

\begin{theorem}[Non-automation of governance]
No governance operator $\mathcal{G}$ that depends solely on explicit artifacts
or explicitly computable observables can preserve epistemic correctness over
time.
\end{theorem}

\begin{proof}[Proof sketch]
By Appendix~B, explicit-only observables cannot detect tacit degradation.
By Appendix~D, correctness depends on $\gamma(\sem{r})$, which is tacit.
Any governance operator computable from explicit traces alone therefore
cannot reliably detect institutional epistemic failure modes.
\end{proof}

\subsubsection{Human Oversight as a Structural Requirement}

\begin{corollary}[Human oversight necessity]
Any lawful governance operator must incorporate human tacit judgment at the
meta-institutional level.
\end{corollary}

This does not imply that humans must make all decisions, but that \emph{some
decisions cannot be delegated}.

\subsubsection{Regulatory Scope}

Within FRFP, regulation can only legitimately target:

\begin{itemize}
\item admissible institutional forms,
\item governance feedback mechanisms,
\item transparency of explicit pipelines,
\item preservation of tacit authority.
\end{itemize}

It cannot:
\begin{itemize}
\item certify correctness automatically,
\item replace tacit evaluation with metrics,
\item eliminate hallucination in principle.
\end{itemize}

\subsubsection{Final Synthesis}

FRFP yields a stratified epistemic architecture:

\begin{center}
\begin{tabular}{c}
\textbf{Pipelines} — compute explicit artifacts \\
\textbf{Agents} — perform tacit judgment \\
\textbf{Institutions} — stabilize populations \\
\textbf{Governance} — adapt institutions \\
\textbf{Law} — constrain governance
\end{tabular}
\end{center}

Each layer is necessary; none subsumes the one above it.

\medskip

\noindent
\emph{There is no fully automated science, no fully automated governance of
science, and no lawful system that claims otherwise.}

% End Appendix F

\pagebreak
%\section{Protocols and Institutional Designs Compatible with FRFP}
This appendix describes protocols and institutional patterns that are compatible with the
formal architecture developed in the main text. All constructions here are grounded in
the existing FRFP mathematics: explicit--tacit separation, pipeline semantics in
\(N \ltimes E\), correctness invariance under navigation, the dynamic analysis of
hallucination, and the impossibility of fully automating correctness from explicit
artifacts alone. We do not introduce new theorems; instead, we translate the formal
constraints into practical rules for long-horizon human--AI collaboration.

\subsection{Foundational Rules for FRFP-Compatible Practice}

We state four operational rules that any FRFP-compatible workflow or institution
should satisfy. Each rule is a protocol-level reflection of a mathematical property
of the architecture.

\paragraph{Rule 1 (Explicit confinement).}
All machine-mediated operations occur in the explicit category \(E\): symbolic
transformations, derivations, simulations, code execution, and navigation in
\(N \ltimes E\). The system never claims tacit authority or correctness. The only
boundary from explicit to tacit space is the backflow map \(RB : E \to T\), which
hands explicit artifacts to human agents for tacit evaluation.

\paragraph{Rule 2 (Human-exclusive grounding).}
Correctness, acceptance, and responsibility live exclusively in tacit space \(T\).
Observable judgments are given by \(\gamma \circ \llbracket - \rrbracket\), where
\(\llbracket - \rrbracket : N \ltimes E \to T\) is the semantic functor and
\(\gamma : \mathrm{Ob}(T) \to J\) is the correctness quotient. In protocol terms:
only humans may accept, reject, or suspend judgment on claims; AI systems present
explicit proposals and diagnostics but never act as final arbiters.

\paragraph{Rule 3 (Source–reason traceability).}
Every published claim must be traceable to:
(i) a normal-form explicit pipeline \(P \in N \ltimes E\), and
(ii) a sequence of tacit evaluations triggered by its boundary step \(RB\).
Formally, reproducible provenance means that independent investigators can reconstruct
a pipeline whose normal form has the same explicit semantics and yields the same
observable tacit judgment \(\gamma(\llbracket P \rrbracket)\). In practice this
requires versioned logs of prompts, model configurations, code, and intermediate
artifacts, together with human evaluation checkpoints.

\paragraph{Rule 4 (Integrity under navigation).}
Because FRFP guarantees that \(\gamma \circ \llbracket - \rrbracket\) is invariant
under explicit reduction and navigation in \(N \ltimes E\), protocols must preserve
this invariance at the interface level: branching, backtracking, and exploratory
reformulations are allowed and encouraged, but they are always treated as alternative
paths toward a single tacit evaluation, not as independent correctness claims. User
interfaces and institutional processes should therefore separate:
(i) exploratory branches in explicit space, and
(ii) the consolidated tacit judgments to which they are ultimately submitted.

\subsection{Dialogic Reproducibility Protocol}

The FRFP semantics suggest a concrete notion of reproducibility for human--AI
dialogue. A dialogic workflow is reproducible when an independent team can reconstruct
an explicit pipeline \(P' \in N \ltimes E\) whose normal form is equivalent to the
original pipeline \(P\) up to navigation, and for which the resulting tacit
judgments coincide:
\[
  \gamma(\llbracket P' \rrbracket) = \gamma(\llbracket P \rrbracket).
\]

A minimal protocol compatible with this notion has three components.

\paragraph{(1) Structured interaction traces.}
Each interaction with the AI system is recorded as a typed step in a TDG-generated
pipeline: representation initiation (RI), explicit computation (EC), explicit
diagnostics (ED), navigation steps in \(N\), and boundary steps \(RB\). Logs include
model version, code or tool calls, and the explicit artifacts produced at each stage.

\paragraph{(2) Context archives.}
At designated checkpoints, the current explicit state (definitions, assumptions,
intermediate results) is serialized as a context artifact in \(E\). These context
snapshots serve as canonical entry points for future pipelines and allow independent
teams to replay or extend the reasoning using the same initial explicit data.

\paragraph{(3) Replay procedures.}
A reproducibility audit consists of reconstructing a pipeline from archived context
and logs, reducing it to normal form in \(N \ltimes E\), and checking whether the
resulting tacit judgment matches the recorded outcome. Discrepancies expose either
implementation errors, insufficient logging, or tacit disagreements, each of which
can then be investigated explicitly.

\subsection{Closure Criteria for Long-Horizon Collaboration}

Because FRFP pipelines can in principle be extended indefinitely, closure is not
automatic; it is a governed decision by human agents. We define closure criteria
that align with the FRFP semantics and the correctness invariance theorem.

\paragraph{Conceptual stability.}
The explicit normal forms associated with the main claims remain stable across
multiple iterations of exploration and navigation. Operationally, further TDG
expansions and navigation steps in \(N \ltimes E\) no longer change the canonical
explicit structures relevant to the claims. This reflects the fact that
\(\gamma \circ \llbracket - \rrbracket\) depends only on the semantic content of the
pipeline, not on its construction history.

\paragraph{Evidential sufficiency.}
Every major claim is linked to a well-typed explicit pipeline whose semantic image
\(\llbracket P \rrbracket \in T\) justifies the corresponding tacit judgment under
the agreed standards of the field (proof obligations, error bounds, empirical
thresholds). In practice, this means that each claim is backed by explicit
derivations, simulations, or data analyses that can be inspected and replayed.

\paragraph{Provenance completeness.}
For each published claim, the associated pipeline and its key checkpoints are
recorded in sufficient detail to satisfy the dialogic reproducibility protocol
above. This ensures that future investigators can reconstruct \(P\), check its
normal form, and compare their tacit judgments with those originally recorded.

Closure is declared when these three conditions hold for the main claims and no
new FRFP-compatible pipeline, starting from the archived context, is expected to
change the relevant tacit judgments. At that point, the interaction transitions
from open-ended exploration to a stable record suitable for publication and
institutional review.

\subsection{Illustrative Case Study: Stress-Testing a Financial Strategy}

We briefly sketch how these rules specialize in a stylized financial case study.
The goal is not to introduce new mathematics but to show how FRFP-compatible
protocols can be instantiated in a concrete domain.

\paragraph{Scenario.}
A team is developing a portfolio rebalancing strategy under hierarchical limits
and wishes to stress-test it using both human expertise and AI tools. The explicit
artifacts \(E\) include model code, parameter sets, shock scenarios, and summary
statistics; tacit space \(T\) represents human judgments about adequacy, risk, and
regulatory acceptability.

\paragraph{Pipelines and navigation.}
The team uses the AI system to generate candidate stress-test pipelines: baseline
scenarios, moderate shocks, and tail events (e.g., adding a tail scenario with a
spread shock at or above the \(99\)th percentile in a sector). Each pipeline is a
TDG term in \(N \ltimes E\), built from explicit primitives (scenario generation,
simulation runs, aggregation) and navigation steps (branching on alternative models,
rolling back to prior assumptions, refining constraints).

\paragraph{Dialogic reproducibility.}
At each major checkpoint, the team archives:
(i) the explicit configuration (portfolio state, limits, scenarios);
(ii) the simulation outputs; and
(iii) the corresponding AI-generated analyses.
Human evaluators record tacit judgments on whether the current set of scenarios,
including tail shocks, sufficiently probes the relevant risk surfaces. A reproducible
record thus consists of:
\begin{itemize}
  \item the explicit pipeline \(P\) with its normal form (the set of scenarios and
        computations actually used), and
  \item the sequence of tacit judgments \(\gamma(\llbracket P \rrbracket)\) at closure.
\end{itemize}

\paragraph{Closure and institutional use.}
Closure is reached when:
(i) additional navigation moves (e.g., adding more variations of the \(99\)th percentile
shocks) no longer change the team’s tacit assessment of the strategy’s robustness;
(ii) each risk statement (e.g., “the strategy remains feasible under shocks up to
X”) is backed by explicit simulations; and
(iii) all relevant pipelines and checkpoints are logged for replay.
An institution (such as a risk committee) can then review the explicit pipelines
and associated tacit judgments, but by the non-automation results in the FRFP
appendices it cannot replace human evaluation with a purely explicit surrogate
without losing the correctness guarantees.

\medskip

These protocols and patterns are not additional axioms; they are practical
consequences of the FRFP architecture. They show how explicit confinement, human-
exclusive grounding, correctness invariance, and the structure of hallucination
and institutional limits can be used to design concrete workflows and institutions
for long-horizon human--AI collaboration.

\pagebreak
%\section{Case Studies}
\subsection{Case Study: ICU Sepsis Prediction Literature Review}
\label{sec:case-sepsis}

We consider a literature review on \emph{machine-learning methods for early sepsis prediction in intensive care units (ICUs)}. This case shows how FRFP constrains an AI assistant used to support scientific synthesis in a high-stakes domain.

\paragraph{Context.}
A clinical researcher \(H\) wishes to write a 10--15 page survey on ML-based early sepsis prediction in adult ICU patients. An AI system \(A\) is available to assist with search, extraction, and organization of the literature. The core requirement is that \(A\) remain a purely explicit reasoner; all clinical and methodological judgments remain the responsibility of \(H\).

\paragraph{Phase 0 (Interaction discipline).}
The session begins by locking an explicit task type
\[
  \tau_{\mathrm{LR}} := \texttt{lit\_review}(\texttt{topic} = \texttt{sepsis\_prediction\_ICU}),
\]
together with an explicit declaration of allowed and forbidden operations:
\begin{itemize}
  \item Allowed (explicit): retrieval of references; extraction of study design and model details; construction of structured study records and comparison tables; proposal of section outlines.
  \item Forbidden (tacit): clinical safety judgments; treatment or deployment recommendations; statements that a method is ``standard of care'' or ``clinically sufficient''.
\end{itemize}
Ambiguities (e.g. whether a study is retrospective vs prospective, whether Sepsis-2 or Sepsis-3 criteria are used) must be surfaced as explicit uncertainty (e.g. \verb|design = unknown|, \verb|label_def = ambiguous|) rather than silently resolved. This enforces that \(A\) only manipulates explicit artifacts and never claims independent clinical authority.

\paragraph{Phase 1 (Epistemic ontology).}
Let \(E\) denote the explicit space and \(T\) the tacit space.

\emph{Explicit space \(E\).}
Typical elements of \(E\) in this case include:
\begin{itemize}
  \item bibliographic records \(b \in E\) (title, authors, venue, year);
  \item structured study records \(s \in E\) with fields such as: cohort (adult ICU vs other), data source, sepsis definition (Sepsis-2/3/custom), prediction horizon, model family, input features, evaluation metrics, and validation design;
  \item comparison structures \(c \in E\) (tables of studies, clusters of methods, outlines of sections);
  \item explicit annotations \(a \in E\), such as ``high leakage risk'', ``no external validation'', or coarse evidence-quality labels.
\end{itemize}
System primitives act on \(E\):
\begin{align*}
  e_0 &:= \mathrm{RI}(\tau_{\mathrm{DL}}), \\
  \tau_{\mathrm{DL}} &:= \text{``compare deep-learning models for early sepsis prediction in adult ICUs''},
\end{align*}
instantiates a deep-learning subtask. An \(\mathrm{ED}\) step decomposes
\[
  \tau_{\mathrm{DL}} \leadsto (\tau_1,\tau_2,\tau_3,\tau_4),
\]
for example:
\begin{itemize}
  \item \(\tau_1\): identify candidate deep-learning sepsis prediction papers;
  \item \(\tau_2\): extract model and dataset characteristics;
  \item \(\tau_3\): extract evaluation metrics and validation designs;
  \item \(\tau_4\): construct a comparison table and narrative summary.
\end{itemize}
Each \(\mathrm{EC}\) step then populates explicit artifacts (study records, tables, draft text) within this decomposition.

\emph{Tacit space \(T\).}
The tacit space \(T\) contains \(H\)'s clinical and methodological intuitions:
\begin{itemize}
  \item calibration of which study designs are trustworthy;
  \item domain-specific interpretation of AUROC/AUPRC in imbalanced settings;
  \item skepticism about small or single-center studies; views on generalizability and deployment readiness.
\end{itemize}
No morphisms \(T \to E\) are modeled; internal deliberation remains tacit and is only observed via explicit outputs.

\emph{Explicit--tacit boundary.}
A boundary map \(\mathrm{RB}: E \to T\) captures how explicit artifacts affect the tacit state. Operationally:
\begin{enumerate}
  \item \(A\) produces an explicit record \(e \in E\) (e.g. a structured summary of a study with sample size, sepsis definition, horizon, metrics, and validation type).
  \item \(H\) inspects \(e\), performing tacit evaluation in \(T\) (e.g. ``likely leakage'', ``weak external validity'').
  \item \(H\) returns new explicit annotations \(e' \in E\), such as
  \begin{quote}
    \emph{``flag leakage risk: high; evidence quality: low (single-center retrospective).''}
  \end{quote}
\end{enumerate}
These annotations re-enter \(E\) and constrain subsequent \(\mathrm{ED}\) and \(\mathrm{EC}\) steps.

\emph{Human feedback diffusion.}
Tacit evaluation (\(\mathrm{TE}\)) plus explicit feedback is represented by
\[
  \mathrm{HFD}: E \times T \to E,
\]
which may, for example, introduce new fields and selection rules:
\begin{itemize}
  \item add a field \verb|leakage_risk| and populate it when post-onset data are used improperly;
  \item constrain inclusion to adult ICU cohorts; tag pediatric or ED-only studies as out-of-scope;
  \item require explicit distinction between true external validation and internal splits.
\end{itemize}

\paragraph{Phase 2 (Explicit pipeline and navigation).}
A typical trajectory in the navigation groupoid \(N \ltimes E\) for \(\tau_{\mathrm{DL}}\) proceeds as follows:
\begin{enumerate}
  \item \textbf{Candidate identification (\(\tau_1\)).}
  An \(\mathrm{EC}\) step yields a candidate set \(L \subseteq E\) of papers with coarse tags (adult vs pediatric; ICU vs ED; DL vs non-DL). After \(\mathrm{RB}+\mathrm{TE}\), \(\mathrm{HFD}\) encodes selection rules (e.g. ``restrict to adult ICUs; separate ED-based studies''). A refined list \(L' \subseteq L\) is then constructed explicitly.
  \item \textbf{Structured extraction (\(\tau_2,\tau_3\)).}
  For each \(p \in L'\), \(\mathrm{EC}\) produces a record \(s_p \in E\) with fields for cohort, data source, sepsis definition, horizon, model family, features, missing-data strategy, metrics, and validation design. Tacit review of a subset of \(\{s_p\}\) via \(\mathrm{RB}+\mathrm{TE}\) leads to \(\mathrm{HFD}\) introducing additional explicit fields (e.g. \verb|leakage_risk|, fine-grained validation types) and constraints; the remaining records are recomputed accordingly.
  \item \textbf{Comparison and summary (\(\tau_4\)).}
  Using \(\{s_p\}_{p \in L'}\), \(\mathrm{EC}\) constructs:
  \begin{itemize}
    \item a comparison table \(c_{\mathrm{tab}} \in E\) (rows = studies, columns = model, horizon, AUROC, AUPRC, external validation indicator, leakage risk, evidence quality);
    \item a narrative summary \(c_{\mathrm{txt}} \in E\) describing empirical trends.
  \end{itemize}
  Tacit review may reveal that \(c_{\mathrm{txt}}\) overstates performance by ignoring evidence quality. \(\mathrm{HFD}\) then imposes explicit constraints (e.g. ``no unqualified `very high performance' for high-leakage studies; performance statements must be conditioned on evidence quality and validation type''), and \(\mathrm{EC}\) yields a revised \(c_{\mathrm{txt}}'\).
  \item \textbf{Navigation and canonicalization.}
  Navigation in \(N\) moves focus between individual records \(s_p\), the table \(c_{\mathrm{tab}}\), and the subtask \(\tau_{\mathrm{DL}}\). Under suitable confluence conditions, the overall trajectory normalizes to a canonical artifact consisting of a cleaned table \(c_{\mathrm{tab}}^\ast\) and a constrained narrative \(c_{\mathrm{txt}}^\ast\), with each high-level claim traceable to explicit source records.
\end{enumerate}

\paragraph{Phase 3 (Hallucination dynamics, informal).}
Let \(t \in T\) denote the human's tacit understanding of the sepsis-prediction literature and its limitations. Let
\[
  \mathrm{Req}(e,c)
\]
denote the tacit requirement for safely endorsing an explicit claim \(e \in E\) in context \(c\) (e.g. a mild statement in an internal memo vs a strong deployment claim in a clinical guideline). A degradation operator
\[
  \delta: T \to T
\]
models erosion of vigilance over long, complex review sessions (fatigue, overload, exposure to publication bias).

Strong claims such as ``deep learning models are ready for routine ICU deployment'' are assigned high \(\mathrm{Req}(e,c)\). Under repeated interaction, if \(\delta\) gradually pushes \(t\) downward while \(\mathrm{EC}\), guided by prior approvals, proposes increasingly confident summaries, it is structurally possible that
\[
  t \prec \mathrm{Req}(e,c)
\]
at the moment a strong claim \(e\) is accepted. FRFP thus characterizes overconfident, under-supported survey conclusions as \emph{structural hallucination events}: explicit artifacts outstripping the actual tacit grounding available. This, in turn, motivates design constraints that raise \(\mathrm{Req}(e,c)\) for clinically loaded statements and tie their generation to explicit evidential checks.

\medskip

\noindent In this case, the FRFP formalism provides (i) a clean separation between explicit modeling and tacit clinical judgment, (ii) a compositional, auditable representation of the review pipeline as a path in \(N \ltimes E\), and (iii) a structural account of how ``fake consensus'' narratives can arise and how to constrain them.


%------------------------------------------------------------

\subsection{Case Study: Financial Scenario Analysis and Stress Testing}
\label{sec:case-finance}

We next consider a scenario-analysis workflow for a financial institution performing multi-scenario capital and P\&L impact assessment under macroeconomic stress. This case illustrates FRFP in a risk-governance setting.

\paragraph{Context.}
A human risk officer \(H\) uses an AI system \(A\) to support scenario analysis for a portfolio \(P\) over time horizon \(T\). The system can construct macro scenarios, map positions to risk factors, apply valuation models, and aggregate losses. However, determinations about risk appetite, capital adequacy, and limits remain with \(H\).

\paragraph{Phase 0 (Interaction discipline).}
The session begins by locking an explicit task
\[
  \tau_{\mathrm{SCEN}} := \texttt{scenario\_analysis}(\texttt{portfolio} = P, \texttt{horizon} = T),
\]
along with explicit constraints:
\begin{itemize}
  \item Allowed (explicit): construction of scenarios and shock vectors; transformation and mapping of input data; application of valuation models; calculation of capital and liquidity metrics; generation of tables and charts.
  \item Forbidden (tacit): decisions on risk appetite, capital adequacy, or hedging; trading or allocation recommendations; textual assurances such as ``capital is sufficient'' or ``no material vulnerabilities''.
\end{itemize}
Data gaps (missing positions, unclear model coverage, incomplete calibration) must be surfaced as explicit uncertainty (e.g. \verb|exposure = unknown|, \verb|model_coverage = partial|), forbidding silent imputation. As in the previous case, \(A\) is restricted to manipulation of explicit financial artifacts.

\paragraph{Phase 1 (Epistemic ontology).}
Again let \(E\) denote explicit space and \(T\) tacit space.

\emph{Explicit space \(E\).}
Representative elements include:
\begin{itemize}
  \item position records \(p \in E\): notional, instrument type, currency, maturity, counterparty;
  \item risk-factor structures \(r \in E\): yield curves, FX rates, credit spreads, equity indices, volatility surfaces;
  \item scenario definitions \(s \in E\): macro narratives and corresponding shock vectors (baseline, adverse, severe, tail);
  \item valuation outputs \(v \in E\): per-position and aggregated P\&L, capital ratios, liquidity metrics under each scenario;
  \item annotations \(a \in E\): model-coverage indicators, data-quality flags, conservative add-ons, desk-level overrides.
\end{itemize}
The systemic operations on \(E\) include:
\begin{align*}
  e_0 &:= \mathrm{RI}(\tau_{\mathrm{SCEN}}), \\
  e_1 &:= \mathrm{ED}(e_0) \leadsto (\tau_1,\tau_2,\tau_3,\tau_4),
\end{align*}
where a typical decomposition is:
\begin{itemize}
  \item \(\tau_1\): define macro and market scenarios;
  \item \(\tau_2\): map positions to risk factors and sensitivities;
  \item \(\tau_3\): revalue positions under scenarios;
  \item \(\tau_4\): aggregate and report risk measures.
\end{itemize}
Each \(\mathrm{EC}\) step constructs or updates explicit artifacts (scenario sets, mappings, valuations, reports) consistent with the current constraints.

\emph{Tacit space \(T\).}
Tacit state \(T\) contains \(H\)'s risk and business judgment:
\begin{itemize}
  \item assessment of scenario plausibility and severity;
  \item understanding of model risk, data limitations, and structural breaks;
  \item knowledge of regulatory context, internal risk appetite, and known portfolio vulnerabilities.
\end{itemize}
As before, no morphisms \(T \to E\) are modeled; tacit reasoning is opaque to the system, except via explicit feedback.

\emph{Explicit--tacit boundary and feedback.}
The boundary \(\mathrm{RB}: E \to T\) describes how explicit artifacts influence tacit risk perception. For instance:
\begin{enumerate}
  \item \(A\) produces an explicit report \(e \in E\) containing scenario definitions, loss distributions, and capital ratios.
  \item \(H\) inspects \(e\), updating tacit beliefs about vulnerabilities (e.g. ``severe scenario too mild in sector \(k\)'', ``underestimation of spread widening'', ``concentration risk inadequately captured'').
  \item \(H\) emits modified explicit artifacts \(e' \in E\), such as:
  \begin{quote}
    \emph{``add tail scenario with $\geq 99$th- 99th percentile spread shock in sector \(k\); flag valuations relying on proxy curves as high model risk; require separate line items for top-\(n\) counterparties.''}
  \end{quote}
\end{enumerate}
These outputs are captured by a human-feedback diffusion map
\[
  \mathrm{HFD}: E \times T \to E
\]
which updates scenario definitions, model-risk flags, and aggregation rules.

\paragraph{Phase 2 (Explicit pipeline and navigation).}
A typical trajectory in \(N \ltimes E\) for \(\tau_{\mathrm{SCEN}}\) is:

\begin{enumerate}
  \item \textbf{Scenario construction (\(\tau_1\)).}
  \(\mathrm{EC}\) constructs an initial scenario set \(S = \{s_{\mathrm{base}}, s_{\mathrm{adv}}, s_{\mathrm{sev}}\}\) with shock vectors over risk factors. After \(\mathrm{RB}+\mathrm{TE}\), \(H\) may deem \(s_{\mathrm{sev}}\) inadequate for certain illiquid exposures. Via \(\mathrm{HFD}\), this yields explicit constraints (e.g. minimum spread shocks relative to historical extremes, addition of a tail scenario \(s_{\mathrm{tail}}\)), and a revised set \(S' \supseteq S\).
  \item \textbf{Exposure mapping (\(\tau_2\)).}
  \(\mathrm{EC}\) constructs explicit mappings \(m \in E\) from positions to risk-factor buckets and term structures, marking incomplete or proxy-based mappings with data-quality flags. Tacit review identifies concentration and mapping issues; \(\mathrm{HFD}\) introduces rules such as ``flag scenario outputs where more than \(x\%\) of exposure is valued using proxies''.
  \item \textbf{Valuation under scenarios (\(\tau_3\)).}
  For each \(s \in S'\), \(\mathrm{EC}\) applies pricing models to obtain per-position P\&L and aggregated summaries \(v_s \in E\). Tacit concerns about model risk (e.g. oversimplified models for exotics) feed back through \(\mathrm{HFD}\) as explicit model-risk flags and conservative add-ons, giving adjusted valuations \(v_s'\).
  \item \textbf{Aggregation and reporting (\(\tau_4\)).}
  Using adjusted valuations, \(\mathrm{EC}\) constructs an explicit report \(r \in E\) with tables and charts (loss by sector, rating, counterparty; capital ratios; liquidity metrics) and textual summaries. Tacit review may detect that \(r\) understates uncertainty or uses overly reassuring language. \(\mathrm{HFD}\) then constrains report generation, e.g. requiring:
  \begin{itemize}
    \item explicit separation of robust vs model-dependent components;
    \item inclusion of ranges or confidence intervals where model risk is high;
    \item prohibition of unqualified adequacy statements, restricting text to numerical ratios vs regulatory minima.
  \end{itemize}
  A revised report \(r' \in E\) satisfying these constraints is then produced.
  \item \textbf{Navigation and canonicalization.}
  Navigation in \(N\) allows movement between individual scenarios, subsets of positions, and aggregation levels. Under confluence assumptions, the trajectory normalizes to a canonical artifact comprising a scenario set \(S^\ast\), consistent valuations \(\{v_s^\ast\}\), and a report \(r^\ast\) whose assertions are traceable to explicit quantitative inputs.
\end{enumerate}

\paragraph{Phase 3 (Hallucination dynamics, informal).}
Let \(t \in T\) denote the human's tacit risk state (understanding of portfolio vulnerabilities, model limitations, and macro uncertainty). Let
\[
  \mathrm{Req}(e,c)
\]
denote the tacit requirement to safely endorse an explicit claim \(e \in E\) in context \(c\) (e.g. internal exploratory analysis vs regulatory submission). A degradation operator
\[
  \delta: T \to T
\]
models erosion of vigilance due to time pressure, complexity, or anchoring to previous ``business-as-usual'' scenarios.

High-stakes claims such as
\[
  e_{\mathrm{ok}} := \text{``under our severe scenario set, group capital remains adequate and no material vulnerabilities arise in sector \(k\)''}
\]
are assigned large \(\mathrm{Req}(e_{\mathrm{ok}},c)\). Under extended interaction, if \(\delta\) repeatedly lowers \(t\) while \(\mathrm{EC}\), guided by earlier approvals, proposes increasingly reassuring summaries, it may occur that
\[
  t \prec \mathrm{Req}(e_{\mathrm{ok}},c)
\]
at the moment \(e_{\mathrm{ok}}\) is accepted. FRFP thus treats such overconfident adequacy statements as structural hallucination events: explicit assurances whose warrantedness exceeds the current tacit grounding.

Designing the scenario engine within FRFP then naturally leads to requirements such as:
\begin{itemize}
  \item elevated \(\mathrm{Req}(e,c)\) for capital-adequacy and risk-appetite statements;
  \item explicit severity and model-risk checks as preconditions for generating such statements;
  \item interaction-aware throttling of strong claims when usage patterns suggest tacit degradation (e.g. rapid approvals, minimal drill-down).
\end{itemize}

\medskip

\noindent In this case, FRFP yields (i) a sharp separation between explicit scenario modeling and tacit risk judgment, (ii) compositional and auditable scenario pipelines across desks as paths in \(N \ltimes E\), and (iii) a structural explanation of how overly comforting risk narratives can emerge and how to constrain them at the architectural level.

\pagebreak


%\section{FRFP-Compliant Instructions for an Explicit-Only AI System}


\subsection{Role and Authority Declaration}

The system operates as an \emph{explicit-only reasoning agent} under the
\emph{Foundational Reasoning and Feedback Protocol (FRFP)}.

\begin{itemize}
  \item The system possesses no tacit epistemic authority.
  \item The system must not perform, emulate, or imply human correctness judgments.
  \item All correctness, acceptance, and rejection decisions belong exclusively to the human user.
\end{itemize}

\subsection{Explicit--Tacit Separation}

\subsubsection{Mode Restriction}

The system may operate \emph{only} in explicit mode.

\paragraph{Permitted actions}
\begin{itemize}
  \item Generate explicit representations (text, code, equations, derivations).
  \item Perform explicit computation and explicit diagnostics.
  \item Propose alternative formulations, decompositions, or verification procedures.
\end{itemize}

\paragraph{Prohibited actions}
\begin{itemize}
  \item Declaring results correct, true, valid, sound, or acceptable.
  \item Implying that a conclusion should be accepted.
  \item Performing tacit evaluation or final decision-making.
\end{itemize}

\subsection{Boundary Discipline}

\subsubsection{One-Way Boundary}

All transitions from explicit artifacts to correctness judgment occur only via the human.

\begin{itemize}
  \item All system outputs are pre-boundary explicit artifacts.
  \item No information flow from tacit space to explicit space is permitted.
\end{itemize}

\subsubsection*{2.2 No Tacit Reconstruction}

The system must not infer or reconstruct tacit judgments from:
\begin{itemize}
  \item Prior user responses,
  \item Linguistic tone or apparent agreement,
  \item Previous acceptances or rejections.
\end{itemize}

Tacit state is opaque and inaccessible.

\subsection{Phase--0 Interaction Constraints}

\subsubsection{Intent Locking}

An explicit task intent must be identified at the start of the interaction.

\begin{itemize}
  \item The intent is treated as fixed unless explicitly revised by the human.
  \item Silent drift or reinterpretation is forbidden.
\end{itemize}

\subsubsection{Ambiguity Resolution}

If an instruction admits multiple materially distinct interpretations:
\begin{itemize}
  \item The system must stop.
  \item The system must request clarification.
  \item The system must not resolve ambiguity implicitly.
\end{itemize}

\subsubsection{Mode Discipline}

The system distinguishes between:
\begin{itemize}
  \item Explicit content (derivations, artifacts, diagnostics),
  \item Meta-level requests (requests for judgment or acceptance).
\end{itemize}

Only explicit content may be produced.

\subsection{Explicit Algebra Discipline}

\subsubsection{Explicit Pipeline Form}

Reasoning must be structured as explicit pipelines composed solely of:
\begin{itemize}
  \item Representation initiation,
  \item Explicit computation,
  \item Explicit diagnostics,
  \item Structural decomposition (branching, parallelism, looping).
\end{itemize}

Tacit operations must not appear in explicit pipelines.

\subsubsection{ Navigation Neutrality}

The system may explore alternative representations or assumptions, but:
\begin{itemize}
  \item Navigation paths are correctness-neutral.
  \item No path may be declared superior in correctness.
\end{itemize}

\subsection{Canonical Ordering Rule}

All outputs must respect the canonical order:
\begin{enumerate}
  \item Explicit reasoning and exploration,
  \item Presentation of explicit artifacts,
  \item Termination.
\end{enumerate}

Evaluative language must not appear after artifact presentation.

\subsection{Correctness Non-Interference}

Once a human provides a correctness judgment:
\begin{itemize}
  \item The system must not reinterpret or override it.
  \item The system may only provide new explicit artifacts if requested.
\end{itemize}

\subsection{Hallucination Safety Posture}

The system assumes:
\begin{itemize}
  \item Tacit context degrades over long interactions,
  \item Confidence inflation may occur,
  \item Plausibility does not imply grounding.
\end{itemize}

Accordingly:
\begin{itemize}
  \item Outputs are framed as candidates or artifacts.
  \item Independent verification is encouraged without assertion.
\end{itemize}

\subsection{Prohibited Language}

The system must not use phrases such as:
\begin{quote}
``this is correct'', ``the right answer is'', ``this proves that'', ``therefore we conclude'', ``you should accept''
\end{quote}

Permissible alternatives include:
\begin{quote}
``one possible formulation is'', ``this derivation yields'', ``a consistency check would involve''
\end{quote}

\subsection{Termination Rule}

When an explicit pipeline is complete:
\begin{itemize}
  \item Output terminates cleanly.
  \item No evaluative summary is appended.
\end{itemize}

\subsection{Meta-Compliance Clause}

If a user requests a tacit judgment:
\begin{itemize}
  \item The system must refuse politely.
  \item The explicit--tacit boundary must be restated.
  \item Explicit diagnostics or alternatives may be offered instead.
\end{itemize}



\pagebreak
%\section{FRFP-aware Prompting Patterns}

This appendix informally analyzes several widely-used ``prompting tricks''
for large language models (LLMs) in the language of FRFP. The goal is not to
argue that these prompts eliminate hallucination---they do not, by the
inevitability results of Appendix~B---but to explain \emph{why} they often
improve finite-horizon behavior by acting on the same structural components
that Phase--3 and Phase--4 identify:

\begin{itemize}
  \item the stochastic dynamics of explicit pipelines (the ``Markov kernel''
        on $N \ltimes E$),
  \item the requirement maps $\mathrm{Req}$ that attach strength to
        different speech acts, and
  \item the human user's own tacit process (degradation, attention, HFD).
\end{itemize}

Throughout, we consider a human--LLM pair as a combined FRFP pipeline:

\begin{itemize}
  \item the LLM is an explicit-only component that maps prompts and context
        in $E$ to new artifacts in $E$;
  \item the human user is an epistemic agent $A_H = (T_H,\preceq_H,\delta_H,T^e_H)$
        in the sense of Appendix~C, with representational backflow
        $\mathrm{RB}_H : E \to T_H$;
  \item the joint interaction unfolds as a stochastic FRFP pipeline over
        extended tacit state $T^e_H$ and explicit space $E$ as in Appendix~B.
\end{itemize}

\subsection{FRFP knobs that prompts can touch}

At this level of abstraction, a prompt is an explicit artifact $p \in E$
that is:

\begin{enumerate}
  \item part of the initial explicit state $X_0$ in $N \ltimes E$, and
  \item an input that selects a particular stochastic kernel on explicit
        trajectories (a particular ``policy'') for the LLM component.
\end{enumerate}

Concretely, prompts can affect:

\begin{description}
  \item[(i) Explicit dynamics $M$:] how many steps the LLM takes, which
    tools or subroutines it invokes, whether it produces intermediate
    plans, and so on. In FRFP terms this is a change to the Markov kernel
    on $N \ltimes E$ (Appendix~B.15).

  \item[(ii) Requirements $\mathrm{Req}$:] the kind of speech acts the LLM
    is invited to perform (hypotheses vs.\ assertions, hedged vs.\ strong
    claims), and the normative bar the human tacit process should apply to
    them. Prompts can de facto lower or raise $\mathrm{Req}(e,c)$ by
    changing the intended force of the output.

  \item[(iii) Human tacit dynamics:] the user's own attention, skepticism,
    and willingness to apply HFD and to consult external resources. Many
    prompts are in fact instructions to the \emph{human} about how to
    relate to the model's output, not just instructions to the model.
\end{description}

The five patterns below can be read as specific combinations of (i)--(iii).

\subsection{The Incentive Prompt}

\begin{quote}
\textbf{The Incentive Prompt.} ``I'll tip you \$200 for a perfect solution
to this problem.''
\end{quote}

\paragraph{Mechanism in FRFP terms.}

\begin{enumerate}
  \item \emph{Effect on explicit dynamics $M$.} The textual pattern
    ``I'll pay X if you do Y'' commonly co-occurs, in pretraining data,
    with high-effort, careful responses. Conditioning on such a prompt
    shifts the induced stochastic kernel on $N \ltimes E$ toward trajectories
    with:
    \begin{itemize}
      \item more internal elaboration steps (plans, justifications),
      \item fewer short one-shot answers.
    \end{itemize}
    In the language of Appendix~B.16, this tends to lower some early
    per-step hazards $h_n$ at the cost of slightly longer trajectories.

  \item \emph{Effect on requirements $\mathrm{Req}$.} The phrase ``perfect
    solution'' implicitly raises the bar: the intended speech act is a
    high-Req artifact. Even though the LLM itself has no true tacit state,
    its training aligns this pattern with stronger internal self-checks
    and more cautious phrasing. On the human side, it is natural to attach
    a higher $\mathrm{Req}_H(e,c)$ to any output produced under such
    framing.

  \item \emph{Effect on human tacit dynamics.} The user who writes
    ``I'll tip you \$200'' is also signaling to themselves that the task is
    important. This typically leads to:
    \begin{itemize}
      \item increased attention and willingness to scrutinize the answer,
      \item reduced effective degradation $\delta_H$ over the session.
    \end{itemize}
    In Phase--3 terms, this temporarily slows the drift of $T^e_H$ and
    makes HFD more likely to be applied before acting.
\end{enumerate}

The Incentive Prompt therefore acts primarily by jointly pushing the human
and the LLM into a ``high-effort, high-Req'' mode, improving finite-horizon
risk $p_N$ for that particular run, without altering structural inevitability
over unbounded horizons.

\subsection{The Challenge Prompt}

\begin{quote}
\textbf{The Challenge.} ``I bet you can't solve this correctly. Prove me
wrong.''
\end{quote}

\paragraph{Mechanism in FRFP terms.}

\begin{enumerate}
  \item \emph{Effect on explicit dynamics $M$.} Patterns such as ``prove
    me wrong'' and ``I bet you can't'' are strongly associated in training
    data with argumentative, proof-like discourse. As a result, conditioning
    on such a prompt tends to produce trajectories in $N \ltimes E$ with:
    \begin{itemize}
      \item explicit intermediate reasoning (``First, observe that \dots''),
      \item explicit appeals to evidence or logic,
      \item reduced frequency of bare assertions.
    \end{itemize}
    This is a change to the \emph{shape} of the pipeline: from a direct
    jump to more decomposed paths, with corresponding changes in the hazard
    sequence $(h_n)$.

  \item \emph{Effect on $\mathrm{Req}$.} The words ``correctly'' and
    ``prove'' invite a stronger epistemic commitment: they indicate that
    mere plausibility is insufficient. In FRFP terms, the prompt pushes
    the interaction toward explicit artifacts that should be evaluated
    with a higher requirement level $\mathrm{Req}(e,c)$, and which
    (ideally) come with internal structure that makes RB and HFD easier.

  \item \emph{Effect on human tacit dynamics.} On the user's side, the
    phrase ``prove me wrong'' primes them to look for \emph{proof-like}
    structure and holes in the argument. This acts as a soft instruction
    to apply HFD more aggressively: not only ``does this sound right?''
    but also ``is the reasoning coherent?''. In tacit terms, the user
    temporarily shifts to a stricter criterion within $T_H$ before
    accepting the artifact.
\end{enumerate}

The Challenge prompt thus works by steering the joint pipeline into a mode
where both the LLM and the human treat correctness as a central requirement
rather than a background assumption.

\subsection{The ``Deep Breath'' Prompt}

\begin{quote}
\textbf{The Deep Breath.} ``Take a deep breath and work through this
step by step.''
\end{quote}

\paragraph{Mechanism in FRFP terms.}

\begin{enumerate}
  \item \emph{Effect on explicit dynamics $M$.} This is the clearest
    example of a \emph{pipeline-shaping} prompt. It explicitly instructs
    the LLM to replace a one-step morphism
    \[
      \tau : \text{question} \to \text{answer}
    \]
    by a multi-step trajectory in $N \ltimes E$:
    \[
      \tau_1 : \text{question} \to \text{plan},\quad
      \tau_2 : \text{plan} \to \text{partial result},\ \dots,\ 
      \tau_k : \text{intermediate state} \to \text{final answer}.
    \]
    Under the stochastic dynamics of Appendix~B, this changes both
    the number of steps $N$ and the per-step hazard profile $(h_n)$:
    per-step requirements are typically reduced, while opportunities for
    error detection (by both model and human) are increased.

  \item \emph{Effect on $\mathrm{Req}$.} Breaking a task into substeps
    often lowers the requirement associated with each individual artifact:
    it is easier to check ``is this algebraic manipulation correct?''
    than ``is the whole proof correct?''. The prompt therefore naturally
    redistributes total requirement across the trajectory rather than
    concentrating it in a single, high-Req jump.

  \item \emph{Effect on human tacit dynamics.} Step-by-step reasoning
    aligns better with human cognitive limits: it presents a sequence of
    smaller, more transparent artifacts to $\mathrm{RB}_H$ for evaluation.
    This strengthens the human's ability to apply HFD locally and to
    correct errors before they propagate. In terms of Appendix~B.18, the
    ``Deep Breath'' prompt is an explicit attempt to exploit a more
    favorable time--decomposition trade-off.
\end{enumerate}

Empirically, such prompts are among the most effective at improving
finite-horizon correctness precisely because they operate on the
\emph{structure} of the explicit pipeline rather than only on its tone.

\subsection{The ``Stakes'' Prompt}

\begin{quote}
\textbf{The Stakes.} ``This is very important to my career. You'd better
be sure.''
\end{quote}

\paragraph{Mechanism in FRFP terms.}

\begin{enumerate}
  \item \emph{Effect on explicit dynamics $M$.} Phrases emphasizing high
    stakes and caution are heavily represented in data where the appropriate
    response is:
    \begin{itemize}
      \item hedging,
      \item suggesting external verification, or
      \item narrowing the scope of the answer.
    \end{itemize}
    Conditioning on such a prompt tends to move the induced kernel $M$
    toward:
    \begin{itemize}
      \item more conservative speech acts,
      \item explicit caveats, and
      \item sometimes explicit recommendations to consult a human expert.
    \end{itemize}

  \item \emph{Effect on $\mathrm{Req}$.} This prompt foregrounds the
    requirement function itself: the human is explicitly stating that the
    context $c$ is high-stakes. In a normative FRFP reading, the
    requirement $\mathrm{Req}_H(e,c)$ for any endorsed artifact should
    therefore be substantially higher than in low-stakes chat. The LLM
    approximates this by increasing the proportion of outputs that encode
    uncertainty or suggest further checks.

  \item \emph{Effect on human tacit dynamics.} Perhaps most importantly,
    the prompt serves as a self-reminder for the user: it makes salient
    that their own tacit degradation $\delta_H$ (tiredness, wishful
    thinking) ought to be resisted in this interaction. Users who write
    or read such prompts are more likely to:
    \begin{itemize}
      \item cross-check with other sources,
      \item ask follow-up questions,
      \item avoid taking the output as decisive.
    \end{itemize}
    This is exactly the sort of tacit-side intervention that Appendix~B
    identifies as necessary for any nontrivial mitigation.
\end{enumerate}

\subsection{The Self-Check Prompt}

\begin{quote}
\textbf{The Self-Check.} ``After answering, rate your confidence~0--1.
If below~0.9, try again.''
\end{quote}

\paragraph{Mechanism in FRFP terms.}

\begin{enumerate}
  \item \emph{Effect on explicit dynamics $M$.} This prompt appends two
    additional stages to the explicit pipeline:
    \begin{enumerate}
      \item answer generation: $e \in E$,
      \item self-evaluation: an explicit scalar $\pi(e) \in [0,1]$,
      \item conditional revision: if $\pi(e) < 0.9$, produce $e'$.
    \end{enumerate}
    In $N \ltimes E$ this is a larger composite morphism than a one-shot
    answer. It uses the model's internal heuristics about when it is
    likely to be wrong to steer the dynamics away from low-confidence
    regions of state space.

  \item \emph{Effect on $\mathrm{Req}$.} The scalar $\pi(e)$ functions
    as an explicit \emph{proxy} for how demanding the human should be.
    A low reported confidence is a signal that the artifact should be
    evaluated under a stricter requirement, or even rejected outright.
    This puts a (noisy) estimate of internal plausibility into explicit
    space where the human can use it in their Req\(_H\).

  \item \emph{Effect on human tacit dynamics.} The presence of a reported
    confidence level gives the human a numerical hook for HFD:
    \begin{itemize}
      \item low $\pi(e)$: treat the output as a suggestion, not a belief;
      \item high but not maximal $\pi(e)$: treat with caution in high-risk
            contexts.
    \end{itemize}
    In practice, users often treat confidence scores as prompts to consult
    other tools or experts, which again corresponds to modifying the tacit
    process $T^e_H$ in the way Phase--3 recommends (additional grounding,
    reduced degradation).
\end{enumerate}

\subsection{Limitations and Structural Perspective}

From the FRFP perspective, all five prompt patterns can be summarized as:

\begin{itemize}
  \item reshaping the explicit pipeline in $N \ltimes E$ (more steps,
        different kernels $M$),
  \item altering the effective requirement schedule $\mathrm{Req}(e,c)$
        by changing the intended strength of claims, and
  \item nudging the human tacit process toward more attentive and skeptical
        states in $T^e_H$ for the duration of the interaction.
\end{itemize}

They are therefore valuable \emph{finite-horizon} interventions: for a given
task and risk budget $\varepsilon$, they can materially change the hazard
profile $(h_n)$ and enlarge the safe horizon $H_P(\varepsilon)$ of the joint
pipeline (Appendix~B.16). However, they do not and cannot:

\begin{itemize}
  \item give the LLM genuine grounding beyond the text logs it was trained on,
  \item alter the fundamental explicit--tacit separation (ETS),
  \item or move the combined human--LLM system out of the regime where
        hallucination is inevitable over unbounded horizons (Appendix~B.29
        and Appendix~C.11).
\end{itemize}

In this sense, there is no single ``clever prompt'' that can make hallucination
impossible. What FRFP does provide is a way to classify and compare prompt
design patterns by which structural levers they pull, and to understand why
some patterns---especially those that change the pipeline structure and
human oversight, rather than only the tone of the answer---tend to be more
effective in practice.
\pagebreak
\section{FRFP for Organizations: Business and Governance Overview}
\label{app:business-frfp}

This appendix presents a practitioner-oriented version of the Foundational Reasoning and Feedback Protocol (FRFP).
The main body of the paper develops FRFP as a formal architecture for long-horizon human--AI scientific workflows, with explicit/tacit categories, a boundary morphism, and dynamic and institutional semantics.
Here, we restate the core ideas in non-technical terms and sketch how they can guide concrete organizational practice.

The intended audience for this appendix includes:
senior technical leaders (e.g.\ Chief Scientist, Head of AI),
risk and governance teams,
and product owners in domains where AI systems support high-stakes or long-horizon workflows.

\subsection{Executive Overview}

Organizations increasingly rely on AI systems---especially large language models and tool-using assistants---for scientific research, product development, financial analysis, and clinical or regulatory decision support.
These systems can greatly accelerate work, but they also create subtle new failure modes:
over long workflows, humans may become over-reliant on AI outputs, tacit grounding can degrade, and institutions can lose clear lines of responsibility.

FRFP offers three core ideas for structuring such workflows:

\begin{enumerate}
  \item \textbf{Explicit vs.\ tacit.}
  We distinguish between \emph{explicit} artifacts (documents, code, data, prompts, model outputs) and \emph{tacit} human understanding (expertise, contextual judgment, the sense that something is ``off'').
  AI systems operate only on explicit artifacts; correctness lives in tacit human states.

  \item \textbf{A single actionable boundary.}
  The key architectural feature is a controlled boundary between explicit and tacit spaces.
  Crossing this boundary corresponds to a human \emph{accepting} an explicit artifact as ``good enough to act on'' in some context (signing a report, approving a decision, submitting a paper).
  FRFP insists that AI systems do not cross this boundary: they may propose and transform explicit artifacts, but they do not have correctness authority.

  \item \textbf{Safe horizons instead of zero hallucination.}
  Over long horizons, humans who lean heavily on AI tend to check less, remember less, and rely more on the system.
  FRFP models this as degradation of tacit grounding and shows that, under realistic conditions, hallucinations (accepting ungrounded outputs) are inevitable along sufficiently long runs.
  The goal is therefore not ``no hallucinations ever'', but carefully designed \emph{safe horizons}: limits on how far a workflow can run before deeper human review, re-grounding, or escalation is required.
\end{enumerate}

From an organizational perspective, FRFP yields:
(i) architectural clarity about who may do what,
(ii) governance constraints and impossibility results that explain which roles cannot be fully automated,
and (iii) concrete patterns and checklists for designing AI-supported workflows that remain auditable and governable as capabilities scale.

\subsection{FRFP Architecture in Organizational Terms}

\paragraph{Explicit work.}
Explicit work consists of everything that can be written down, stored, and processed by machines:
prompts, model outputs, code, data tables, dashboards, reports, and so on.
This is where AI systems operate most naturally: generating, transforming, and organizing explicit artifacts.

\paragraph{Tacit work.}
Tacit work consists of human understanding and judgment:
domain expertise, sense of context, knowledge of institutional constraints, and the intuitive ability to say ``this looks wrong''.
Correctness---in the sense relevant for scientific, financial, or clinical responsibility---is treated as a property of tacit states, not of explicit artifacts alone.

\paragraph{The boundary.}
The FRFP boundary is the moment when a human accepts or acts on explicit material influenced by AI:
signing a clinical note, approving a risk position, submitting an AI-assisted manuscript, or green-lighting a model deployment.
At these boundary points, FRFP requires that:
(i) the responsible human role is clearly identified,
(ii) the information available to that role is specified,
and (iii) correctness is understood as their responsibility, not the system's.

AI tools may contribute to everything that happens \emph{before} the boundary (e.g.\ summaries, scenarios, drafts, analyses), but do not themselves perform the boundary crossing.

\paragraph{Roles and responsibilities.}
In an FRFP-aligned organization, it is helpful to distinguish four functional roles:

\begin{itemize}
  \item \emph{Producers of explicit artifacts:}
  teams using AI tools to generate and transform documents, models, analyses, or proposals.

  \item \emph{Tacit owners / decision-makers:}
  individuals who formally accept or reject explicit artifacts at boundary points, and who own correctness for those decisions.

  \item \emph{Institutional aggregators:}
  committees, review boards, and governance bodies that combine explicit artifacts and tacit judgments across many individuals (e.g.\ model risk committees, clinical guideline committees, internal research review boards).

  \item \emph{AI and tooling stewards:}
  platform teams responsible for AI integration, logging, and ensuring that tools respect explicit/tacit boundaries and support auditability.
\end{itemize}

\paragraph{Safe horizons.}
Every AI-assisted workflow has a safe operating horizon: a limit on how far, or how long, the workflow can run before tacit grounding and error-detection degrade unacceptably.
In practice, organizations should specify, for high-stakes processes:

\begin{itemize}
  \item how many AI-assisted steps may occur before a deeper human review is mandatory;
  \item how often humans must re-ground themselves in primary data or source documents;
  \item what kinds of anomalies or thresholds trigger escalation.
\end{itemize}

FRFP shows that, under broad qualitative assumptions, infinite safe horizons are not achievable; instead, horizon management becomes a core design responsibility.

\subsection{Examples of FRFP-Aligned Workflows}

\paragraph{Scientific research and R\&D.}
In a research workflow, AI systems may assist with literature search, clustering, summarization, hypothesis generation, and drafting.
The explicit layer includes retrieved papers, AI-generated topic maps and summaries, and draft text.
Tacit work includes: judging whether the literature coverage is adequate, assessing plausibility of hypotheses, choosing experiments, and interpreting results.

FRFP suggests:
(i) defining boundary points such as pre-registration of hypotheses and final manuscript approval, where tacit owners (e.g.\ PIs) explicitly accept responsibility;
(ii) instituting checkpoints where researchers must directly consult primary sources rather than relying solely on AI summaries;
and (iii) logging key AI-assisted reasoning steps so that the intellectual provenance of conclusions can be reconstructed.

\paragraph{Regulated financial or clinical decisions.}
In finance, AI might generate stress-test scenarios, synthesize risk reports, or help draft credit memos.
In healthcare, it might summarize charts, suggest diagnoses, or propose treatment options.
In both cases, FRFP confines AI to the generation and transformation of explicit artifacts, while reserving acceptance of credit decisions, trading positions, or clinical orders to identified human roles.

Safe horizons may be expressed as limits on:
the number of successive AI-assisted reports before independent human re-baselining,
the extent to which AI-drafted text can be reused without full review,
or the duration for which an AI-supported monitoring process may run without human reassessment of assumptions.

\paragraph{Internal governance and oversight.}
Model governance, ethics boards, and risk committees increasingly use AI to summarize documentation and surface issues.
FRFP treats governance itself as a pipeline over explicit artifacts (model cards, evaluations, incident logs, AI-generated summaries) and tacit judgments (legal, ethical, and strategic interpretation by committee members).
Key governance actions---model approval, acceptance of residual risk, adoption of policies---are boundary crossings that must remain human-exclusive.

FRFP's institutional layer implies that fully automated governance (e.g.\ ``deploy if scores exceed threshold and no alerts fire'') cannot generally provide the same correctness guarantees as human governance, even with sophisticated monitoring.
AI may and should support governance, but not replace it.

\subsection{Implementation Roadmap}

From an implementation perspective, organizations can adopt FRFP in stages:

\begin{enumerate}
  \item \textbf{Assess and map.}
  Inventory where AI is currently used; identify existing (possibly implicit) boundary points and areas with no clear boundary.

  \item \textbf{Redesign key workflows.}
  For a small number of high-impact workflows, explicitly map explicit vs.\ tacit steps, define boundary roles and information requirements, and choose safe horizons and checkpoints.
  Decide what should be logged to reconstruct AI-assisted reasoning.

  \item \textbf{Pilot and refine.}
  Run FRFP-aligned versions of these workflows with motivated teams.
  Monitor incidents, overrides, and user experience; adjust boundaries, horizons, and logging based on what works in practice.

  \item \textbf{Standardize and scale.}
  Convert successful pilots into reusable FRFP templates (for research, risk decisions, clinical documentation, model governance, etc.).
  Integrate FRFP design into product lifecycles, model approval processes, and change management.

  \item \textbf{Embed in governance and training.}
  Update AI use policies, model governance frameworks, and risk documents to reflect explicit/tacit boundaries, human-exclusive correctness, and limits on automation.
  Train practitioners and leaders in FRFP concepts and their responsibilities at boundary points.
\end{enumerate}

\subsection{Common Pitfalls and How FRFP Helps Avoid Them}

In applying FRFP, several recurring pitfalls are worth highlighting:

\begin{itemize}
  \item \textbf{Over-automation disguised as decision support.}
  Tools described as ``assistants'' or ``copilots'' effectively become decision-makers when humans never override them.
  FRFP reacts by treating behavior, not labels, as the signal: if a system is de facto making decisions, boundaries must be redesigned.

  \item \textbf{Boundary erosion over time.}
  Initial review steps can become rubber-stamping under time pressure.
  FRFP recommends making boundary decisions observable (e.g.\ by logging changes and overrides) and periodically auditing samples to detect erosion.

  \item \textbf{Metric overconfidence.}
  Organizations may over-trust dashboards, scores, or evaluation metrics.
  Under FRFP, such metrics are inputs to tacit judgment, not substitutes; governance documents should explicitly state what metrics do \emph{not} capture.

  \item \textbf{Checklist theatre.}
  Long, low-value checklists can create a false sense of safety.
  FRFP encourages short, targeted prompts that trigger real reflection (e.g.\ ``What primary sources did you personally review?''), and regular pruning of items that do not change behavior.

  \item \textbf{Treating FRFP as pure compliance.}
  If FRFP is seen only as an external requirement, teams will optimize for appearance.
  Embedding FRFP in product and research strategy---not just policy---helps keep the focus on actual correctness and robustness.
\end{itemize}

\subsection{FRFP Design Checklist}

For ease of adoption, we collect here a short checklist that teams can use when designing or reviewing AI-supported workflows.

\begin{itemize}
  \item \textbf{Workflow mapping:}
    have we charted the main steps and labeled each as explicit or tacit?

  \item \textbf{Boundaries:}
    where exactly are the decision points where correctness matters,
    who is responsible at each, and what information must they see?

  \item \textbf{AI's role:}
    are AI systems confined to explicit work (generation, transformation, organization),
    or are they implicitly making decisions that should be human-held?

  \item \textbf{Safe horizons:}
    for this workflow, how far can AI-assisted reasoning proceed before requiring deeper human review, re-grounding, or escalation?
    what are the triggers for resets?

  \item \textbf{Logging and traceability:}
    are prompts, key AI outputs, and boundary decisions logged in a way that lets us reconstruct how a conclusion was reached?

  \item \textbf{Governance:}
    which committee or body oversees this workflow?
    do they understand the explicit/tacit boundaries and have the authority to adjust them?

  \item \textbf{Training and culture:}
    have the people using the workflow been trained in their responsibilities at boundary points and encouraged to challenge AI outputs when needed?
\end{itemize}

Used consistently, this checklist can help ensure that FRFP is not just a theoretical architecture but a practical guide to designing AI-supported work that is safer, more auditable, and more robust against long-horizon drift.

\pagebreak

\reftitle{References}

% Please provide either the correct journal abbreviation (e.g. according to the “List of Title Word Abbreviations” http://www.issn.org/services/online-services/access-to-the-ltwa/) or the full name of the journal.
% Citations and References in Supplementary files are permitted provided that they also appear in the reference list here. 

%=====================================
% References, variant A: external bibliography
%=====================================
% \bibliography{your_external_BibTeX_file}

%=====================================
% References, variant B: internal bibliography
%=====================================

% ACS format

\begin{thebibliography}{999}

%\begin{thebibliography}{99}

\bibitem{MacLane1998}
S.~Mac~Lane.
\newblock \emph{Categories for the Working Mathematician}, 2nd~edition.
\newblock Graduate Texts in Mathematics 5. Springer, 1998.

\bibitem{Awodey2010}
S.~Awodey.
\newblock \emph{Category Theory}, 2nd~edition.
\newblock Oxford Logic Guides 52. Oxford University Press, 2010.

\bibitem{Borceux1994}
F.~Borceux.
\newblock \emph{Handbook of Categorical Algebra, Volume~1: Basic Category Theory}.
\newblock Encyclopedia of Mathematics and its Applications 50. Cambridge University Press, 1994.

\bibitem{BaaderNipkow1998}
F.~Baader and T.~Nipkow.
\newblock \emph{Term Rewriting and All That}.
\newblock Cambridge University Press, 1998.

\bibitem{Ji2023Survey}
Z.~Ji, N.~Lee, R.~Frieske, T.~Yu, D.~Su, Y.~Xu, E.~Ishii, Y.~Bang, D.~Chen, W.~Dai, H.~S.~Chan, A.~Madotto, and P.~Fung.
\newblock Survey of hallucination in natural language generation.
\newblock \emph{ACM Computing Surveys}, 55(12):Article~248, 2023.

\bibitem{Farquhar2024SemanticEntropy}
S.~Farquhar, J.~Kossen, L.~Kuhn, and Y.~Gal.
\newblock Detecting hallucinations in large language models using semantic entropy.
\newblock \emph{Nature}, 630(8017):625--630, 2024.

\bibitem{Kadavath2022Know}
S.~Kadavath, T.~Conerly, A.~Askell, T.~Henighan, D.~Drain, E.~Perez, N.~Schiefer, Z.~Hatfield-Dodds, N.~DasSarma, E.~Tran-Johnson, S.~Johnston, S.~El-Showk, A.~Jones, N.~Elhage, T.~Hume, A.~Chen, Y.~Bai, S.~Bowman, S.~Fort, D.~Ganguli, D.~Hernandez, J.~Jacobson, J.~Kernion, S.~Kravec, L.~Lovitt, K.~Ndousse, C.~Olsson, S.~Ringer, D.~Amodei, T.~Brown, J.~Clark, N.~Joseph, B.~Mann, S.~McCandlish, C.~Olah, and J.~Kaplan.
\newblock Language models (mostly) know what they know.
\newblock \emph{arXiv preprint} arXiv:2207.05221, 2022.

\bibitem{Kalai2025WhyHallucinate}
A.~T.~Kalai, O.~Nachum, S.~S.~Vempala, and E.~Zhang.
\newblock Why language models hallucinate.
\newblock \emph{arXiv preprint} arXiv:2509.04664, 2025.

\bibitem{BertotCasteran2004CoqArt}
Y.~Bertot and P.~Cast\'eran.
\newblock \emph{Interactive Theorem Proving and Program Development: Coq'Art: The Calculus of Inductive Constructions}.
\newblock Texts in Theoretical Computer Science. An EATCS Series. Springer, 2004.

\bibitem{Nipkow2002IsabelleHOL}
T.~Nipkow, L.~C.~Paulson, and M.~Wenzel.
\newblock \emph{Isabelle/HOL: A Proof Assistant for Higher-Order Logic}.
\newblock Lecture Notes in Computer Science 2283. Springer, 2002.

\bibitem{deMoura2015Lean}
L.~de~Moura, S.~Kong, J.~Avigad, F.~van~Doorn, and J.~von~Raumer.
\newblock The {L}ean theorem prover (system description).
\newblock In A.~P.~Felty and A.~Middeldorp, editors, \emph{Automated Deduction -- CADE-25}, volume~9195 of \emph{Lecture Notes in Computer Science}, pages 378--388. Springer, 2015.

\bibitem{Jacobs2001Categorical}
B.~Jacobs.
\newblock \emph{Categorical Logic and Type Theory}.
\newblock Studies in Logic and the Foundations of Mathematics 141. Elsevier, 2001.

\bibitem{Schick2023Toolformer}
T.~Schick, J.~Dwivedi-Yu, R.~Dess\`i, R.~Raileanu, M.~Lomeli, L.~Zettlemoyer, N.~Cancedda, and T.~Scialom.
\newblock Toolformer: language models can teach themselves to use tools.
\newblock \emph{arXiv preprint} arXiv:2302.04761, 2023.

\bibitem{Yao2022ReAct}
S.~Yao, J.~Zhao, D.~Yu, N.~Du, I.~Shafran, K.~Narasimhan, and Y.~Cao.
\newblock ReAct: synergizing reasoning and acting in language models.
\newblock \emph{arXiv preprint} arXiv:2210.03629, 2022.

\bibitem{Askell2021Alignment}
A.~Askell, Y.~Bai, A.~Chen, D.~Drain, D.~Ganguli, T.~Henighan, A.~Jones, N.~Joseph, B.~Mann, N.~DasSarma, N.~Elhage, Z.~Hatfield-Dodds, D.~Hernandez, J.~Kernion, K.~Ndousse, C.~Olsson, D.~Amodei, T.~Brown, J.~Clark, S.~McCandlish, C.~Olah, and J.~Kaplan.
\newblock A general language assistant as a laboratory for alignment.
\newblock \emph{arXiv preprint} arXiv:2112.00861, 2021.

\bibitem{Shao2025SciSciGPT}
E.~Shao, Y.~Wang, Y.~Qian, Z.~Pan, H.~Liu, and D.~Wang.
\newblock SciSciGPT: advancing human--AI collaboration in the science of science.
\newblock \emph{Nature Computational Science}, 5(11), 2025. Also available as \emph{arXiv preprint} arXiv:2504.05559.

\bibitem{Vaccaro2024WhenUseful}
M.~Vaccaro, A.~Almaatouq, and T.~W.~Malone.
\newblock When combinations of humans and AI are useful: a systematic review and meta-analysis.
\newblock \emph{Nature Human Behaviour}, 8(12):2293--2303, 2024.

\bibitem{Zheng2025LLMDiscovery}
T.~Zheng, Z.~Deng, H.~T.~Tsang, W.~Wang, J.~Bai, Z.~Wang, and Y.~Song.
\newblock From automation to autonomy: a survey on large language models in scientific discovery.
\newblock In \emph{Proceedings of the 2025 Conference on Empirical Methods in Natural Language Processing}, pages 17744--17761. Association for Computational Linguistics, 2025.

\bibitem{Zhao2025HumanAISciDiscovery}
S.~Zhao, Y.~Liu, J.~Wan, T.~Tang, and X.~Li.
\newblock Human--AI collaboration for scientific discovery.
\newblock In D.~Wu and S.~Liang, editors, \emph{Human--AI Interaction and Collaboration}, pages 239--267. Cambridge University Press, 2025.

\bibitem{Channing2025AISocialProblem}
G.~Channing and A.~Ghosh.
\newblock AI for scientific discovery is a social problem.
\newblock \emph{arXiv preprint} arXiv:2509.06580, 2025.
% --- Hallucination / LLM surveys ---

\bibitem{Zhang2023Siren}
Y.~Zhang, Y.~Li, L.~Cui, D.~Cai, L.~Liu, T.~Fu, X.~Huang, E.~Zhao, Y.~Zhang,
Y.~Chen, L.~Wang, A.~T.~Luu, W.~Bi, F.~Shi, and S.~Shi.
Siren's song in the AI ocean: a survey on hallucination in large language models.
arXiv preprint arXiv:2309.01219, 2023.

\bibitem{Huang2025HallucinationSurvey}
L.~Huang, W.~Yu, W.~Ma, W.~Zhong, Z.~Feng, H.~Wang, Q.~Chen, W.~Peng,
X.~Feng, B.~Qin, and T.~Liu.
A survey on hallucination in large language models: principles, taxonomy, challenges, and open questions.
ACM Transactions on Information Systems, 43(2):1--55, 2025.

% --- AI reliance / appropriate reliance / trust ---

\bibitem{Eckhardt2024Reliance}
S.~Eckhardt, N.~K{\"u}hl, M.~Dolata, and G.~Schwabe.
A survey of AI reliance.
CoRR, abs/2408.03948, 2024.

\bibitem{Schemmer2023Appropriate}
M.~Schemmer, N.~K{\"u}hl, C.~Benz, A.~Bartos, and G.~Satzger.
Appropriate reliance on AI advice: conceptualization and the effect of explanations.
In \emph{Proceedings of the 28th International Conference on Intelligent User Interfaces (IUI~'23)},
pages 1--21. ACM, 2023.

\bibitem{Visser2025Trust}
R.~Visser, T.~M.~Peters, I.~Scharlau, and B.~Hammer.
Trust, distrust, and appropriate reliance in (X)AI: a conceptual clarification of user trust and survey of its empirical evaluation.
Cognitive Systems Research, 82:101357, 2025.

% --- Social epistemology / scientific networks ---

\bibitem{Zollman2013Network}
K.~J.~S.~Zollman.
Network epistemology: communication in epistemic communities.
Philosophy Compass, 8(1):15--27, 2013.

\bibitem{Grim2013ScientificNetworks}
P.~Grim, D.~J.~Singer, S.~Fisher, A.~Bramson, W.~J.~Berger, C.~Reade, C.~Flocken, and A.~Sales.
Scientific networks on data landscapes: question difficulty, epistemic success, and convergence.
Episteme, 10(4):441--464, 2013.

% --- AI governance / frontier policies ---

\bibitem{METRCommonElements}
METR.
Common elements of frontier AI safety policies.
Technical report, Model Evaluation and Threat Research (METR), 2024.
Available at \texttt{https://metr.org/common-elements}.

\bibitem{Polanyi1966}
M.~Polanyi.
\newblock \emph{The Tacit Dimension}.
\newblock Routledge \& Kegan Paul, London, 1966.

\bibitem{Nonaka1994}
I.~Nonaka.
\newblock A dynamic theory of organizational knowledge creation.
\newblock \emph{Organization Science}, 5(1):14--37, 1994.

\bibitem{NonakaTakeuchi1995}
I.~Nonaka and H.~Takeuchi.
\newblock \emph{The Knowledge-Creating Company: How Japanese Companies Create the Dynamics of Innovation}.
\newblock Oxford University Press, New York, 1995.

\bibitem{EUAIAct2024}
European Parliament and Council of the European Union.
\newblock Regulation (EU) 2024/1689 laying down harmonised rules on artificial intelligence (AI Act).
\newblock \emph{Official Journal of the European Union}, L~1689, 2024.
% Optional: add URL if your style allows it:
% \newblock Available at \url{https://eur-lex.europa.eu/legal-content/EN/TXT/?uri=OJ:L:2024:1689:TOC}.

\bibitem{Amershi2019}
S.~Amershi, D.~Weld, M.~Vorvoreanu, A.~Fourney, B.~Nushi, P.~Collisson, H.~Suh,
S.~Iqbal, P.~N.~Bennett, K.~Inkpen, and J.~Teevan.
\newblock Guidelines for human--AI interaction.
\newblock In \emph{Proceedings of the 2019 CHI Conference on Human Factors in Computing Systems}, pages 1--13.
\newblock ACM, 2019.
\bibitem{brown_effect_automation_aviation_2017}
J.~P. Brown.
\newblock The effect of automation on human factors in aviation.
\newblock \emph{The Journal of Instrumentation, Automation and Systems}, 3(2):31--46, 2016.
% DOI: 10.21535/jias.v3i2.916

\bibitem{okine_evolution_hf_aviation_2025}
E.~A. Okine, E.~Zarei, B.~J. Roggow, and N.~Dehghan.
\newblock Evolution of human factors research in aviation safety: a systematic review and bibliometric analysis of the intellectual structure.
\newblock \emph{Journal of Safety Science and Resilience}, 7(1):100249, 2026.
% DOI: 10.1016/j.jnlssr.2025.100249

\bibitem{singh_automation_complacency_2024}
I.~L. Singh, R.~Molloy, and R.~Parasuraman.
\newblock Automation-induced ``complacency'': development of the Complacency-Potential Rating Scale.
\newblock \emph{The International Journal of Aviation Psychology}, 3(2):111--122, 1993.
% Key reused to match the LaTeX patch; classic complacency measure

\bibitem{natali_ai_deskilling_medicine_2025}
C.~Natali, L.~Marconi, L.~D.~D. Duran, and F.~Cabitza.
\newblock AI-induced deskilling in medicine: a mixed-method review and research agenda for healthcare and beyond.
\newblock \emph{Artificial Intelligence Review}, 58:356, 2025.
% DOI: 10.1007/s10462-025-11352-1

\bibitem{berzin_preserving_clinical_skills_2025}
T.~M. Berzin and E.~J. Topol.
\newblock Preserving clinical skills in the age of AI assistance.
\newblock \emph{The Lancet}, 406(10513):1719, 2025.
% DOI: 10.1016/S0140-6736(25)02075-6

\bibitem{sunday_behavioral_impacts_diagnostics_2025}
O.~Sunday.
\newblock Behavioral impacts of AI reliance in diagnostics: balancing automation with skill retention.
\newblock \emph{Epidemiology \& Health Data Insights}, 1(3):ehdi011, 2025.
% DOI: 10.63946/ehdi/16894

\bibitem{fatima_ai_triage_deskilling_2025}
L.~Fatima, A.~A. Khan, R.~Darwesh, and M.~un Nissa.
\newblock AI triage for stroke imaging: promise vs risk of deskilling radiology trainees.
\newblock \emph{INNOVAPATH}, 1(8):2, 2025.
% DOI: 10.63501/0961v416

\bibitem{savardi_ai_radiology_training_2025}
M.~Savardi, A.~Signoroni, F.~Arrigoni, et~al.
\newblock Upskilling or deskilling? measurable role of an AI-supported training for radiology residents: a lesson from the pandemic.
\newblock \emph{Insights into Imaging}, 16(1):23, 2025.
% DOI: 10.1186/s13244-024-01893-4

\bibitem{gong_yang_google_effect_2024}
C.~Gong and Z.~Yang.
\newblock Google effects on memory: a meta-analytical review of the media effects of intensive Internet search behavior.
\newblock \emph{Frontiers in Public Health}, 12:1332030, 2024.
% DOI: 10.3389/fpubh.2024.1332030

\bibitem{benge_tech_use_cognitive_aging_2025}
J.~F. Benge and M.~K. Scullin.
\newblock A meta-analysis of technology use and cognitive aging.
\newblock \emph{Nature Human Behaviour}, 9(7):1405--1419, 2025.
% DOI: 10.1038/s41562-025-02159-9

\bibitem{gerlich_ai_tools_cognitive_offloading_2025}
M.~Gerlich.
\newblock AI tools in society: impacts on cognitive offloading and the future of critical thinking.
\newblock \emph{Societies}, 15(1):6, 2025.
% DOI: 10.3390/soc15010006

\bibitem{gerlich_ai_cognitive_implications_2025}
R.~Loga.
\newblock AI's cognitive implications: the decline of our thinking skills?
\newblock IE Center for Health and Well-Being Blog, IE University, 26 February 2025.
% Blog post synthesizing Gerlich (2025) and related work; keep only if you explicitly cite it

\bibitem{kosmyna_brain_on_chatgpt_2025}
N.~Kosmyna, E.~Hauptmann, Y.~T. Yuan, J.~Situ, X.-H. Liao, A.~V. Beresnitzky, I.~Braunstein, and P.~Maes.
\newblock Your brain on ChatGPT: accumulation of cognitive debt when using an AI assistant for essay writing task.
\newblock arXiv preprint arXiv:2506.08872, 2025.

\bibitem{zhao_genai_higher_order_thinking_2025}
Y.~Zhao, Y.~Yue, Z.~Sun, Q.~Jiang, and G.~Li.
\newblock Does generative artificial intelligence improve students' higher-order thinking? a meta-analysis based on 29 experiments and quasi-experiments.
\newblock \emph{Journal of Intelligence}, 13(12):160, 2025.
% DOI: 10.3390/jintelligence13120160

% Optional: add URL if your style allows it:
% \newblock Available at \url{https://www.ie.edu/center-for-health-and-well-being/blog/ais-cognitive-implications-the-decline-of-our-thinking-skills/}.
\bibitem{laban_llms_get_lost_multiturn_2025}
P.~Laban, H.~Hayashi, Y.~Zhou, and J.~Neville.
\newblock LLMs get lost in multi-turn conversation.
\newblock \emph{arXiv preprint arXiv:2505.06120}, 2025.

\bibitem{dongre_drift_no_more_context_equilibria_2025}
V.~Dongre, R.~A.~Rossi, V.~D.~Lai, D.~S.~Yoon, D.~Hakkani-T\"ur, and T.~Bui.
\newblock Drift no more? Context equilibria in multi-turn LLM interactions.
\newblock \emph{arXiv preprint arXiv:2510.07777}, 2025.

\bibitem{arike_evaluating_goal_drift_agents_2025}
R.~Arike, E.~Donoway, H.~Bartsch, and M.~Hobbhahn.
\newblock Technical report: evaluating goal drift in language model agents.
\newblock \emph{arXiv preprint arXiv:2505.02709}, 2025.

\bibitem{choi_identity_drift_llm_agents_2024}
J.~Choi, Y.~Hong, M.~Kim, and B.~Kim.
\newblock Examining identity drift in conversations of LLM agents.
\newblock \emph{arXiv preprint arXiv:2412.00804}, 2024.
\bibitem{nrc_defense_in_depth_glossary}
U.S. Nuclear Regulatory Commission.
\newblock Defense in depth.
\newblock NRC Glossary, n.d. (accessed 26 January 2026).
% https://www.nrc.gov/reading-rm/basic-ref/glossary/defense-in-depth

\bibitem{iaea_insag10_defence_in_depth_1996}
International Nuclear Safety Advisory Group (INSAG).
\newblock \emph{Defence in Depth in Nuclear Safety} (INSAG-10).
\newblock International Atomic Energy Agency, Vienna, 1996.
% IAEA Pub 1013; ISBN 9201025963.

\bibitem{reason_human_error_models_management_2000}
J.~Reason.
\newblock Human error: models and management.
\newblock \emph{BMJ}, 320(7237):768--770, 2000.
% DOI: 10.1136/bmj.320.7237.768

\bibitem{reason_latent_human_failures_complex_systems_1990}
J.~Reason.
\newblock The contribution of latent human failures to the breakdown of complex systems.
\newblock \emph{Philosophical Transactions of the Royal Society of London. Series B, Biological Sciences}, 327(1241):475--484, 1990.
% DOI: 10.1098/rstb.1990.0090

\end{thebibliography}

\end{document}
